\documentclass[11pt]{article}
\usepackage[T1]{fontenc}
\usepackage[utf8]{inputenc}
\usepackage{lmodern,amsmath,amssymb,amsthm,mathtools}
\usepackage[margin=1in]{geometry}
\usepackage{xurl}
\usepackage[hidelinks]{hyperref}
\usepackage{microtype}
\hypersetup{pdftitle={Universality of the transport law for supercritical quasi-periodic operators},pdfauthor={Seth Douglas and Nidhal Mghirbi}}
\newcommand{\C}{\mathbb C}
\newcommand{\N}{\mathbb N}
\newcommand{\Z}{\mathbb Z}
\newcommand{\Q}{\mathbb Q}
\newcommand{\R}{\mathbb R}
\newcommand{\T}{\mathbb T}
\newcommand{\un}{\mathbf 1}
\newcommand{\tr}{\operatorname{tr}}
\newcommand{\ad}{\operatorname{ad}}
\newcommand{\dist}{\operatorname{dist}}
\newcommand{\supp}{\operatorname{supp}}
\newcommand{\Lip}{\operatorname{Lip}}
\newcommand{\Ga}{G_\alpha}
\newcommand{\au}{\alpha_u^+}
\newcommand{\eps}{\varepsilon}
\newtheorem{theorem}{Theorem}[section]
\newtheorem{lemma}[theorem]{Lemma}
\newtheorem{proposition}[theorem]{Proposition}
\newtheorem{corollary}[theorem]{Corollary}
\newtheorem{maintheorem}{Theorem}
\renewcommand{\themaintheorem}{\Alph{maintheorem}}
\newtheorem{maincorollary}[maintheorem]{Corollary}
\theoremstyle{definition}
\newtheorem{definition}[theorem]{Definition}
\newtheorem{remark}[theorem]{Remark}
\setlength{\emergencystretch}{2em}
\raggedbottom
\title{Universality of the transport law for\\ supercritical quasi-periodic operators\thanks{The authors developed this work with AI models from more than one developer and take responsibility for every claim. How the work was done is described at the end of the paper.}}
\author{Seth Douglas \and Nidhal Mghirbi}
\date{Version 1.0.0, October 2026}
\begin{document}
\maketitle
\begin{abstract}
We determine how fast wave packets spread under quasi-periodic Schr\"odinger operators $(Hu)(n)=u(n+1)+u(n-1)+v(\theta+n\alpha)u(n)$ with real-analytic $v$ and positive Lyapunov exponent $L$ on the spectrum. Let $\beta(\alpha)=\limsup_n q_n^{-1}\log q_{n+1}$, let $\ell_*$ be the minimum of $L$ on the spectrum, and put $F(\ell)=\max\{0,1-\ell/\beta(\alpha)\}$, with $F\equiv0$ if $\beta(\alpha)=0$ and $F\equiv1$ if $\beta(\alpha)=\infty$. If $v$ is of type~I in the sense of Avila and of Ge and Jitomirskaya, then for almost every phase the fastest transport exponent $\alpha_u^+$ and the exponent of every moment, both upper and Abel-averaged, equal $F(\ell_*)$; the lower bound for $\alpha_u^+$ holds for every phase. The law also holds energy by energy: for packets filtered ever more narrowly near an energy $E$ of the spectrum, every moment exponent tends to $F(L(E))$, and the spectral measure has upper local dimension $F(L(E))$ at almost every $E$. These conclusions extend to $v(x)=w(sx)$, with $\beta(s\alpha)$ in place of $\beta(\alpha)$, when $w$ at frequency $s\alpha$ is supercritical and of type~I. Beyond type~I, we show, conditionally on one hypothesis, that every real-analytic potential with positive Lyapunov exponent on its spectrum has $\alpha_u^+=F(\ell_*)$, computed with $\beta(s\alpha)$, for almost every phase, where $s$ is the period reduction of the potential. The hypothesis: at fixed energy, the zeros of the rational discriminants of its primitive part, as functions of the phase, do not cluster at an exponential rate near the drifting phase. The fastest-transport law for analytic potentials is thereby reduced to this weak separation hypothesis for primitive potentials with acceleration at least two. The proofs combine a covering argument over all times, an exact $q$-step resolvent identity, elliptic runs of Weyl solutions, and an arcsine law for the rational discriminant.
\end{abstract}

{\small\tableofcontents}

\section{Introduction}\label{sec:intro}

Consider a one-frequency quasi-periodic Schr\"odinger operator on $\ell^2(\Z)$,
\begin{equation}\label{eq:qp}
(H_{v,\alpha,\theta}u)(n)=u(n+1)+u(n-1)+v(\theta+n\alpha)\,u(n),
\end{equation}
with a real-analytic sampling function $v$ on $\T=\R/\Z$, an irrational frequency $\alpha$ and a phase $\theta\in\T$. The best-known example is the almost Mathieu operator, $v(x)=2\lambda\cos2\pi x$, known to physicists as the Aubry--Andr\'e model~\cite{AubryAndre1980} and realized with cold atoms in bichromatic optical lattices~\cite{Roati2008} and with light in quasi-periodic waveguide arrays~\cite{Lahini2009}. When the Lyapunov exponent is positive on the spectrum, the folklore physical picture is that every state is localized.

That picture depends on the arithmetic of $\alpha$. Let $p_n/q_n$ be the continued-fraction convergents of $\alpha$ and
\begin{equation}\label{eq:beta}
\beta(\alpha)=\limsup_{n\to\infty}\frac{\log q_{n+1}}{q_n}\in[0,\infty].
\end{equation}
Avila, You and Zhou~\cite{AYZ2017} proved that the almost Mathieu operator with $\lambda>1$ has Anderson localization for almost every phase when $\lambda>e^{\beta(\alpha)}$, and purely singular continuous spectrum for every phase when $1<\lambda<e^{\beta(\alpha)}$. Ge and Jitomirskaya extended sharp arithmetic transitions of this kind to supercritical operators of type~I~\cite{GeJitomirskaya2024}, a class defined through Avila's global theory~\cite{Avila2015}. Strongly Liouville frequencies thus destroy localization even though the Lyapunov exponent stays positive.

What does a wave packet do in that regime? This paper answers the question for every supercritical type~I analytic potential, with one formula. Three points summarize the answer.
\begin{itemize}
\item Transport is a competition between two exponential rates, the localization strength $L$ and the Liouville strength $\beta(\alpha)$, and its exponent is $\max\{0,1-L/\beta(\alpha)\}$.
\item The law is universal and local in energy. The acceleration does not enter; the farthest tails are set by the smallest Lyapunov exponent, every moment follows the same law, and the part of the packet at energy $E$ spreads with exponent $\max\{0,1-L(E)/\beta(\alpha)\}$.
\item Beyond type~I, the fastest-tail law follows from a single geometric hypothesis on the zeros of rational discriminants, which we isolate exactly.
\end{itemize}

\subsection{The transport exponent}

Let $\psi_t=e^{-itH}\delta_0$ and average in time with the exponential weight $(2/T)e^{-2t/T}$:
\begin{equation}\label{eq:abel}
a(j,T)=\frac2T\int_0^\infty e^{-2t/T}\,|\langle\delta_j,\psi_t\rangle|^2\,dt,\qquad
P(N,T)=\sum_{|j|>N}a(j,T),\qquad M_p(T)=\sum_{j\in\Z}|j|^p a(j,T).
\end{equation}
$P(N,T)$ is the time-averaged probability of finding the particle beyond distance $N$, and $M_p(T)$ is the averaged $p$-th moment of the position. The upper transport exponents are
\begin{equation}\label{eq:exponents}
\beta^+(p)=\limsup_{T\to\infty}\frac{\log M_p(T)}{p\log T},\qquad
\au=\sup_{p>0}\beta^+(p).
\end{equation}
The quantity $\au$ is the fastest upper exponent of Damanik and Tcheremchantsev~\cite{DT2007}. It sees even polynomially small tails: $\au\le a$ holds when $P(T^a,T)$ decays faster than every power of $T$ (Lemma~\ref{lem:moments}), and a tail of size $T^{-C}$ at distance $T^{a}$ already forces $\au\ge a$. Always $0\le\beta^+(p)\le\au\le1$, since hopping is bounded. For $0<\beta<\infty$ we write
\[
F_\beta(\ell)=\max\Bigl\{0,\ 1-\frac{\ell}{\beta}\Bigr\},
\]
with $F_0\equiv0$ and $F_\infty\equiv1$.

\subsection{Main results}

Write $L(E)$ for the Lyapunov exponent of $H_{v,\alpha,\theta}$ at energy $E$ and $\Sigma$ for its spectrum, which does not depend on $\theta$. Following Avila~\cite{Avila2015}, an energy is \emph{supercritical} if $L(E)>0$. The acceleration of $E$ is the right derivative at $y=0$ of $y\mapsto L(E,y)/(2\pi)$, where $L(E,y)$ is the Lyapunov exponent of the cocycle with phase complexified by $iy$; it is a positive integer at supercritical spectral energies. Following Ge and Jitomirskaya~\cite{GeJitomirskaya2024}, a supercritical energy is \emph{of type~I} if its acceleration is one (Section~\ref{sec:typeI} recalls the precise definition). The almost Mathieu operator with $\lambda>1$ is supercritical and of type~I at every spectral energy~\cite{Avila2015}, and so is every sufficiently small analytic perturbation of it.

\begin{maintheorem}[The transport law]\label{thm:A}
Let $v$ be real-analytic and $\alpha$ irrational, and assume that every energy in $\Sigma$ is supercritical and of type~I. Let $\ell_*=\min_{E\in\Sigma}L(E)>0$.
\begin{enumerate}
\item If $\beta(\alpha)<\infty$, then for every $\theta$ in a set $G^v_\alpha$ of full measure,
\[
\au(H_{v,\alpha,\theta})=\max\Bigl\{0,\ 1-\frac{\ell_*}{\beta(\alpha)}\Bigr\},
\]
read as $0$ when $\beta(\alpha)=0$, and $P(T^a,T)=O(T^{-m})$ for every $a>F_{\beta(\alpha)}(\ell_*)$ and $m>0$. When $v$ is even, $G^v_\alpha$ contains the explicit set $\Ga$ of~\eqref{eq:Ga}. When $\beta(\alpha)=0$ the equality holds for every phase.
\item If $\beta(\alpha)=\infty$, then $\au=1$ for every phase.
\end{enumerate}
The lower bound $\au\ge1-\ell_*/\beta(\alpha)$ holds for every phase.
\end{maintheorem}

The set $G^v_\alpha$ is explicit (Definition~\ref{def:tI-good}): it consists of the phases $\theta$ for which $q_n\theta$ stays at distance at least $2q_n^{-3}$ from the turning points of the band centres of the $p_n/q_n$-periodic approximants, for all large $n$. The exponent vanishes when $\beta(\alpha)\le\ell_*$, including the threshold itself. When $\beta(\alpha)>0$ it grows to $1$ as $\ell_*\downarrow0$; when $\beta(\alpha)=0$ it vanishes. For the almost Mathieu operator the threshold $\beta(\alpha)=\log\lambda$ is the arithmetic spectral transition of Avila, You and Zhou.

The formula has a simple heuristic origin. At a resonant scale $q=q_n$ the potential is almost $q$-periodic, with a slow drift of size $|q\alpha-p|^{-1}\approx e^{\beta q}$ in space. On that scale the packet moves like a particle in a nearly periodic band of width about $e^{-L(E)q}$, so its speed is about $e^{-L(E)q}$. The slow drift turns this into a Bloch-type oscillation that lasts for a time about $T=e^{\beta q}$. In that time the packet covers a distance about $Te^{-L(E)q}=T^{1-L(E)/\beta}$. The fastest spreading comes from the energies with the smallest Lyapunov exponent. The theorems show that no other mechanism moves any polynomial fraction of the packet faster, at any time.

\begin{remark}[The almost Mathieu operator]\label{rem:amo}
For $v(x)=2\lambda\cos2\pi x$ with $\lambda>1$, Theorem~\ref{thm:A} applies with $\ell_*=\log\lambda$ and gives $\au=\max\{0,1-\log\lambda/\beta(\alpha)\}$ for every $\theta\in\Ga$. The companion paper~\cite{DMamo} proves this special case independently, by an elementary route built on an exactly conserved quantity: at each denominator $q$, a degree-$q$ polynomial of $H$ is a slowly varying diagonal plus an operator of range $q$. It also determines every moment and the packing dimension of every spectral measure of the almost Mathieu operator, and it shows that the transport exponent can be noncomputable.
\end{remark}

The next theorem extends the law from the fastest tails to every moment, and resolves it in energy. Let $\mu_\theta$ be the spectral measure of $\delta_0$, and let $\ell_\theta\ge\ell_*$ and $\ell^\theta$ be the minimum and the maximum of $L$ on its support. For a finite measure $\mu$ the upper local dimension is $d^+_\mu(E)=\limsup_{r\to0}\log\mu([E-r,E+r])/\log r$, and the upper and lower packing dimensions $\dim_P^\pm(\mu)$ are its $\mu$-essential supremum and infimum.

\begin{maintheorem}[Every moment, energy by energy]\label{thm:B}
Assume the hypotheses of Theorem~\ref{thm:A} with $0<\beta(\alpha)<\infty$, write $\beta=\beta(\alpha)$, and let $\theta\in G^v_\alpha$.
\begin{enumerate}
\item For every $p>0$, $\beta^+(p)=F_\beta(\ell_*)$. If $\ell_*<\beta$, then $L$ attains $\ell_*$ on $\supp\mu_\theta$, and for every $\gamma>0$ a fixed positive fraction of the packet lies beyond distance $e^{(b_n-\ell_*-\gamma)q_n}$ at the times $T_n=|q_n\alpha-p_n|^{-1}$, for all large $n$ with $b_n=q_n^{-1}\log T_n\ge\ell_*+15\gamma$.
\item Let $E_0\in\Sigma$, $0<\gamma<\ell_*/2$, and let $g\in C_c^\infty(\R)$ be real, supported in $\{E:|L(E)-L(E_0)|<\gamma\}$, with $g(H)\delta_0\ne0$. Then the packet started at $g(H)\delta_0$ satisfies, for every $p>0$,
\[
F_\beta\bigl(L(E_0)+\gamma\bigr)\le\beta^+_{g(H)\delta_0}(p)\le F_\beta\bigl(L(E_0)-\gamma\bigr).
\]
\item $d^+_{\mu_\theta}(E)=F_\beta(L(E))$ for $\mu_\theta$-almost every $E$. Consequently $\dim_P^+(\mu_\theta)=F_\beta(\ell_\theta)$ and $\dim_P^-(\mu_\theta)=F_\beta(\ell^\theta)$.
\end{enumerate}
\end{maintheorem}

Part~1 says that, at every phase of $G^v_\alpha$, the fastest-tail exponent and every upper moment exponent coincide, and that the critical distance is reached by a fixed fraction of the packet, not only by a polynomially small tail. Part~2 says that a packet filtered near an energy $E$ spreads, in every moment, with exponent $F_\beta(L(E))$, in the limit of narrow filters. Part~3 is its spectral counterpart: the upper local dimension of the spectral measure is an explicit function of the energy. It is nonconstant, and the two packing dimensions differ, as soon as $\ell_*<\beta$ and $L$ is not constant on $\supp\mu_\theta$. For $\theta\in G^v_\alpha\cap(G^v_\alpha-\alpha)$ every nonzero spectral measure $\mu_\varphi$ satisfies $d^+_{\mu_\varphi}=F_\beta\circ L$ almost everywhere, so $\dim_P^+(\mu_\varphi)\le F_\beta(\ell_*)$ (Corollary~\ref{cor:allmeasures}). For frequencies with $\beta(\alpha)>\max_\Sigma L$, lower bounds of the form in parts~1 and~3 hold at every phase for every Lipschitz potential whose Lyapunov exponent is continuous on the spectrum (Theorem~\ref{thm:U1}).

The type~I hypothesis singles out the correct frequency. If $v(x)=w(sx)$ for an integer $s\ge2$, then $v(\theta+n\alpha)=w(s\theta+ns\alpha)$, so $H_{v,\alpha,\theta}=H_{w,s\alpha,s\theta}$: the operator is really driven by the frequency $s\alpha$, and $\beta(s\alpha)$ can differ from $\beta(\alpha)$. Such a $v$ has acceleration at least $s$ and is never of type~I.

\begin{maintheorem}[Covariance]\label{thm:Cov}
Let $v$ be real-analytic and nonconstant, let $s$ be the greatest common divisor of the frequencies $k\ne0$ with $\hat v(k)\ne0$, and write $v(x)=w(sx)$. If $H_{w,s\alpha}$ is supercritical and of type~I, then for almost every $\theta$,
\[
\au(H_{v,\alpha,\theta})=\max\Bigl\{0,\ 1-\frac{\ell_*}{\beta(s\alpha)}\Bigr\},
\]
read as $0$ when $\beta(s\alpha)=0$, and for every $\theta$ if $\beta(\alpha)\in\{0,\infty\}$. Moreover $\beta(\alpha)\le\beta(s\alpha)\le s\beta(\alpha)$.
\end{maintheorem}

If the primitive part $w$ has acceleration at least two somewhere on the spectrum of $H_{w,s\alpha}$, then $H_{w,s\alpha}$ is not of type~I and Theorem~\ref{thm:Cov} does not apply. For them the rational discriminant, as a function of the phase at a fixed energy, can have several zeros near the real axis, and the zeros can cluster. Our last result shows that excluding such clusters is the only input our argument still needs. For a convergent $p/q$ let $t_q(x,E)$ be the trace of the transfer matrix over one period of the $p/q$-periodic approximant at the phase $x/q$. We say that WXS$(v,\alpha)$ holds (weak separation) if for every $\eps>0$ and $c>0$, for almost every $\theta$ and all large $n$ with $q_n^{-1}\log|q_n\alpha-p_n|^{-1}\ge c$, no energy within $1/q_n$ of the spectrum has two zeros of $t_{q_n}(\cdot,E)$, real or complex and counted with multiplicity, within $e^{-\eps q_n}$ of each other with one of them within $e^{-\eps q_n}$ of $q_n\theta$ (Definition~\ref{def:WXS}).

\begin{maintheorem}[A conditional complete law]\label{thm:D}
Let $v$ be real-analytic with positive Lyapunov exponent on its spectrum, and write $v(x)=w(sx)$ as in Theorem~\ref{thm:Cov}. If WXS$(w,s\alpha)$ holds, then for almost every $\theta$,
\[
\au(H_{v,\alpha,\theta})=\max\Bigl\{0,\ 1-\frac{\ell_*}{\beta(s\alpha)}\Bigr\},
\]
read as $0$ when $\beta(s\alpha)=0$, where $\ell_*$ is the minimum of the Lyapunov exponent on the spectrum.
\end{maintheorem}

By Theorems~\ref{thm:A} and~\ref{thm:Cov}, the hypothesis is needed only when the primitive part $w$ has acceleration at least two at some point of the spectrum. For real-analytic potentials with positive Lyapunov exponent on the spectrum, the transport law is therefore reduced to one precise open hypothesis: \emph{WXS for primitive potentials with acceleration at least two.} We do not know whether it holds. For the two-harmonic family $2\lambda\cos2\pi x+2\mu\cos4\pi x$ we reduce it further to a large-deviation estimate, (E-$\Delta$), for the measure of the energies at which one explicit polynomial of degree at most $2q$ is exponentially small (Proposition~\ref{prop:EDelta}). The hypothesis must be imposed on the primitive part: for non-primitive potentials it can fail by symmetry alone (Remark~\ref{rem:primitive}).

\subsection{Relation to earlier work}

For the almost Mathieu operator two endpoints were known. Damanik and Tcheremchantsev proved $\au=0$ for $\lambda>1$ when $\beta(\alpha)=0$~\cite{DT2007}, and Jitomirskaya and Zhang proved $\au=1$ for every phase when $\beta(\alpha)=\infty$~\cite{JZ2022}. Jitomirskaya and Kr\"uger proved exponential dynamical localization in phase expectation when $\beta(\alpha)$ is sufficiently small, including positive values~\cite{JK2013}, and Bradshaw, de Jong, Liu, Wang, Wang and Yang refined the growth rates for weakly Liouville frequencies with $\beta(\alpha)=0$~\cite{BdJLWWY2025}. In between, lower bounds of the form $1-C\log\lambda/\beta$ were obtained by Jitomirskaya and Zhang and by Haeming~\cite{JZ2022,Haeming2025}, and upper bounds on fractal dimensions of spectral measures near the transition by Jitomirskaya, Liu and Tcheremchantsev~\cite{JLT2025}; Levi extended their packing and generalized-dimension bounds, with the same factor two, to potentials that are monotone on the period~\cite[Theorems~1.1 and~1.3]{Levi2025}. Jitomirskaya, Liu and Tcheremchantsev also wrote that they believe the same formulation, with $\log\lambda$ replaced by $\min L(E)$, should work in the other cases where a sharp arithmetic transition has been established~\cite[Section~1]{JLT2025}. For supercritical type~I analytic potentials, Theorem~\ref{thm:B}(3) and Corollary~\ref{cor:allmeasures} confirm this, with factor one in place of two, at phases in $G^v_\alpha\cap(G^v_\alpha-\alpha)$; potentials that are only monotone on the period, as in Levi's work, are not covered. Haeming also proved lower bounds for every moment, at every phase and along sequences of times, with exponent $1-d$ when $\beta(\alpha)>3\ell_*/d$ and $0<d<1/2$~\cite[Theorem~1.1]{Haeming2025}. The exact value of $\au$ for finite positive $\beta(\alpha)>\log\lambda$ was not known even for the almost Mathieu operator; the companion paper~\cite{DMamo} determines it. Here we treat general supercritical analytic potentials. For them the endpoint $\beta(\alpha)=0$ was observed by Jitomirskaya and Mavi~\cite[Remark~1.12]{JitomirskayaMavi2017}. The structural inputs are Avila's global theory~\cite{Avila2015}, the type~I class of Ge and Jitomirskaya~\cite{GeJitomirskaya2024}, the finite harmonic structure of rational discriminants of Jitomirskaya and Marx~\cite{JitomirskayaMarx2012} and Shamis~\cite{Shamis2011}, and, for the spectral consequences, the inequality of Barbaroux, Germinet and Tcheremchantsev between transport and generalized dimensions~\cite{BGT2000}.

\subsection{How the proofs work}

\emph{The lower bound} (Section~\ref{sec:lower}) constructs, at a resonant scale $q=q_n$ and for any $h$ slightly above $\ell_*$, a band of energies on which complex-energy transfer matrices stay bounded by $e^{2hq}$ over a distance $e^{(b_n-h)q+O(\log q)}$. It applies to any Lipschitz sampling function whose Lyapunov exponent is continuous near a spectral minimizer.

\emph{The upper bound} has two layers. The outer layer (Section~\ref{sec:covering}) is a covering argument that works for any potential. It reduces the upper bound to two estimates about the operator. The first, (M1), confines the packet to distance $O(q_n)$ up to times $e^{\kappa q_n}$ at every denominator; it comes from the complex-energy method of Damanik and Tcheremchantsev~\cite{DT2007}. The second, (M2), describes the motion at a scale where $|q_n\alpha-p_n|$ is exponentially small: speed $e^{-\ell_*q_n}$, up to a stopping radius. The inner layer (Section~\ref{sec:typeI}) proves (M2) without any conserved operator. Its single object is the resolvent vector at a fixed complex energy, which satisfies an exact $q$-step recursion whose coefficient is the rational discriminant at the local phase, with an error controlled by the Lyapunov exponent at that energy. Where the discriminant is large the recursion forces decay; where it is small its imaginary part has a fixed sign, and a weighted estimate bounds the speed. The type~I hypothesis enters only through the geometry of the zeros of the discriminant near the phase of the initial site.

\emph{Every moment} (Section~\ref{sec:moments}). Two new ingredients give a fixed fraction of the packet, not a polynomially small one, at the critical distance. First, an arcsine law: at resonant scales, the rational discriminant evaluated along the spectral measure is distributed like the discriminant of a periodic operator, whose law is exactly the arcsine law by the periodic magic formula of Damanik, Killip and Simon~\cite{DKS2010}. Second, an elliptic run: at an energy where the transfer matrix over one period is elliptic with a margin, and at the damping $\operatorname{Im}z=|q_n\alpha-p_n|$, the Weyl solution keeps its size, in a frame that follows the slowly drifting periods, over at least $e^{(b_n-h)q_n}$ sites for every $h>L(E)$, and a flux identity controls everything beyond the run. A Taylor expansion of energy filters localizes both statements in energy, and the window estimate of Section~\ref{sec:typeI} makes the arcsine law local.

\emph{Beyond type~I} (Section~\ref{sec:beyond}). The rational discriminant of an analytic potential is, up to an exponentially small error, a trigonometric polynomial of fixed degree in the phase. Its exact factorization over its zeros shows that the window estimate needs only weak separation of the zeros near the drifting phase.

\subsection{Organization}

Section~\ref{sec:prelim} fixes notation and collects elementary facts. Section~\ref{sec:covering} states (M1) and (M2) and proves the covering theorem. Section~\ref{sec:lower} proves the lower bound. Section~\ref{sec:typeI} treats supercritical type~I potentials and proves Theorems~\ref{thm:A} and~\ref{thm:Cov}. Section~\ref{sec:moments} proves Theorem~\ref{thm:B}. Section~\ref{sec:beyond} proves Theorem~\ref{thm:D} and treats the two-harmonic family. Section~\ref{sec:discussion} discusses open problems.

\section{Preliminaries}\label{sec:prelim}

\subsection{Notation}

Throughout, $v:\T\to\R$ is continuous, and $H=H_{v,\alpha,\theta}$ is the operator~\eqref{eq:qp}. When $v$ is Lipschitz we write $\Lip(v)$ for its Lipschitz constant. We write $S$ for the shift $(Su)(n)=u(n+1)$, $X$ for the position operator $X\delta_j=j\delta_j$, and $\un_A$ for multiplication by the indicator of $A\subset\Z$. For a rational frequency $p/q$ we use $H_{v,p/q,\varphi}$ in the same way. The spectrum $\Sigma$ of $H_{v,\alpha,\theta}$ does not depend on $\theta$ when $\alpha$ is irrational, since lattice translations conjugate $H_{v,\alpha,\theta}$ to $H_{v,\alpha,\theta+k\alpha}$, the orbit $\theta+k\alpha$ is dense, and the operator depends norm-continuously on the phase.

\emph{Transfer matrices.} For $E\in\C$ and $x\in\T$ put $T(E,x)=\begin{pmatrix}E-v(x)&-1\\1&0\end{pmatrix}$ and, for $r\ge0$,
\[
A_r(E,x)=T(E,x+(r-1)\alpha)\cdots T(E,x),\qquad A_0=I.
\]
A solution of $Hu=Eu$ satisfies $(u(r),u(r-1))^{\mathsf T}=A_r(E,\theta)(u(0),u(-1))^{\mathsf T}$. For a rational frequency $p/q$ we write $A^{(p/q)}_r$ for the same products with $\alpha$ replaced by $p/q$. The Lyapunov exponent is
\[
L(E)=\lim_{r\to\infty}\frac1r\int_\T\log\|A_r(E,x)\|\,dx=\inf_{r\ge1}\frac1r\int_\T\log\|A_r(E,x)\|\,dx .
\]
As an infimum of continuous functions of $E$, it is upper semicontinuous. For real-analytic $v$ it is continuous in $E$~\cite[Theorem~1]{BourgainJitomirskaya2002}.

\emph{Continued fractions.} Let $p_n/q_n$ be the convergents of $\alpha$ and put
\[
\delta_n=q_n\alpha-p_n,\qquad D_n=|\delta_n|^{-1},\qquad \beta_n=\frac{\log q_{n+1}}{q_n},\qquad b_n=\frac{\log D_n}{q_n}.
\]
The classical inequalities $(q_{n+1}+q_n)^{-1}<|\delta_n|<q_{n+1}^{-1}$ give
\begin{equation}\label{eq:bn}
q_{n+1}<D_n<2q_{n+1},\qquad 0<b_n-\beta_n<\frac{\log2}{q_n}.
\end{equation}
In particular $\limsup b_n=\limsup\beta_n=\beta(\alpha)$. Denominators grow at least like Fibonacci numbers, and $|\alpha-p_n/q_n|=|\delta_n|/q_n<q_n^{-2}$. When the scale $n$ is fixed we drop the index and write $q,p,\delta,D,b$.

\emph{Transport for a general initial state.} For a vector $\psi\in\ell^2(\Z)$ put $\psi_t=e^{-itH}\psi$ and define $a_\psi(j,T)$, $P_\psi(N,T)$ and $M_{p,\psi}(T)$ as in~\eqref{eq:abel} with $\psi_t$ in place of $e^{-itH}\delta_0$. For $\psi\ne0$ the exponents $\beta^+_\psi(p)$ and $\alpha^+_{u,\psi}$ are defined as in~\eqref{eq:exponents}; they do not change when $\psi$ is multiplied by a nonzero constant. We say that $\psi$ has \emph{finite moments} if $\sum_j|j|^p|\psi(j)|^2<\infty$ for every $p>0$. When $\psi=\delta_0$ we drop the index. We write $\mu_\psi$ for the spectral measure of $\psi$ and $\mu_\theta=\mu_{\delta_0}$ for that of $\delta_0$ and $H_{v,\alpha,\theta}$.

\begin{lemma}[Parseval]\label{lem:parseval}
Let $\psi\in\ell^2(\Z)$, $T>0$, and write $G_\psi(n;z)=\langle\delta_n,(H-z)^{-1}\psi\rangle$. Then
\[
a_\psi(n,T)=\frac1{\pi T}\int_\R\bigl|G_\psi(n;E+i/T)\bigr|^2\,dE .
\]
Moreover $\sum_n|G_\psi(n;z)|^2=\operatorname{Im}\langle\psi,(H-z)^{-1}\psi\rangle/\operatorname{Im}z$ for $\operatorname{Im}z>0$, and for every bounded Borel function $f$,
\[
\frac1{\pi T}\int_\R f(E)\,\bigl\|(H-E-i/T)^{-1}\psi\bigr\|^2\,dE=\int_\R(\mathcal P_{1/T}*f)\,d\mu_\psi,\qquad \mathcal P_\epsilon(x)=\frac{\epsilon}{\pi(x^2+\epsilon^2)} .
\]
\end{lemma}

\begin{proof}
Let $\phi(t)=e^{-t/T}\langle\delta_n,\psi_t\rangle$ for $t\ge0$ and $\phi(t)=0$ for $t<0$. Its Fourier transform is $\int_0^\infty e^{iEt}\phi(t)\,dt=\langle\delta_n,\int_0^\infty e^{-it(H-E-i/T)}dt\,\psi\rangle=-iG_\psi(n;E+i/T)$. Plancherel's theorem gives $\int_\R|\phi|^2dt=(2\pi)^{-1}\int_\R|G_\psi(n;E+i/T)|^2dE$, and $a_\psi(n,T)=(2/T)\int|\phi|^2dt$. The second statement is the spectral theorem: both sides equal $\int|x-z|^{-2}d\mu_\psi(x)$. For the third, $\pi^{-1}T^{-1}|x-E-i/T|^{-2}=\mathcal P_{1/T}(x-E)$; integrate against $f(E)\,dE$ and $d\mu_\psi(x)$ and use Fubini's theorem.
\end{proof}

This is the Abel form of the Parseval formula used by Damanik and Tcheremchantsev~\cite{DT2007}.

\subsection{The phase set}

For an irrational $\alpha$ put
\begin{equation}\label{eq:Ga}
\Ga=\bigl\{\theta\in\T:\ |\sin(2\pi q_n\theta)|\ge q_n^{-2}\ \text{for all sufficiently large }n\bigr\},
\end{equation}
and for $\theta\in\Ga$ let $n_G(\theta)$ be an index from which on this holds.

\begin{lemma}\label{lem:Ga}
For every irrational $\alpha$ the set $\Ga$ has full Lebesgue measure.
\end{lemma}

\begin{proof}
The map $\theta\mapsto q_n\theta$ preserves Lebesgue measure on $\T$, and $\{x\in\T:|\sin 2\pi x|<u\}$ has measure at most $u$. Hence $\{\theta:|\sin 2\pi q_n\theta|<q_n^{-2}\}$ has measure at most $q_n^{-2}$. Since $\sum_nq_n^{-2}<\infty$, the Borel--Cantelli lemma shows that almost every $\theta$ lies in only finitely many of these sets.
\end{proof}

\subsection{Moments}

\begin{lemma}\label{lem:moments}
Let $\psi$ be a unit vector with finite moments, and let $H$ have a bounded real potential.
\begin{enumerate}
\item For every $p>0$ there is $C_{p,\psi}$, depending only on $p$ and on finitely many moments of $\psi$, such that $M_{p,\psi}(T)\le C_{p,\psi}(1+T)^p$ for all $T>0$.
\item For $0<a<1$ and $T\ge1$, $M_{p,\psi}(T)\le T^{ap}+\bigl(M_{2p,\psi}(T)\,P_\psi(T^a,T)\bigr)^{1/2}$.
\item If $0<a<1$ and $P_\psi(T^a,T)=O(T^{-m})$ for every $m>0$, then $\alpha^+_{u,\psi}\le a$.
\end{enumerate}
\end{lemma}

\begin{proof}
1. Let $k\ge1$ be an integer. We have $[X,H]=S^{-1}-S$ and $[X^k,H]=\sum_{i=0}^{k-1}X^{i}[X,H]X^{k-1-i}$. Moving the shifts to the left with $c(X)S^{\pm1}=S^{\pm1}c(X\mp1)$ writes $[X^k,H]$ as $\sum_{\pm}S^{\pm1}c_\pm(X)$, where $c_\pm$ are polynomials of degree at most $k-1$ with coefficients bounded by a constant depending only on $k$. Hence $\|[X^k,H]\phi\|\le C_k\sum_{j<k}\|X^j\phi\|$. Duhamel's formula $X^k\psi_t=e^{-itH}X^k\psi-i\int_0^te^{-i(t-s)H}[X^k,H]\psi_s\,ds$, justified by truncating $X$ to bounded sets, and induction on $k$ give $\|X^k\psi_t\|\le C_{k,\psi}(1+|t|)^k$. For real $p>0$ choose an integer $k$ with $2k\ge p$. Jensen's inequality for the probability distribution $|\psi_t(j)|^2$ gives
\[
\sum_j|j|^p|\psi_t(j)|^2\le\Bigl(\sum_j|j|^{2k}|\psi_t(j)|^2\Bigr)^{p/(2k)}=\|X^k\psi_t\|^{p/k}\le C(1+|t|)^p.
\]
Averaging against $(2/T)e^{-2t/T}$ proves the claim.

2. Split $M_{p,\psi}(T)$ into $|j|\le T^a$ and $|j|>T^a$. The first part is at most $T^{ap}$, because $\sum_ja_\psi(j,T)=1$. Apply the Cauchy--Schwarz inequality to the second part, for the probability measure $a_\psi(\cdot,T)$ on $\Z$.

3. By parts 1 and 2, $M_{p,\psi}(T)\le T^{ap}+C(1+T)^{p}T^{-m/2}$. Taking $m\ge2p$ gives $M_{p,\psi}(T)\le CT^{ap}$, so $\beta^+_\psi(p)\le a$ for every $p$.
\end{proof}

\subsection{Uniform growth and rational products}

\begin{lemma}\label{lem:growth}
Let $K\subset\R$ be compact and $h>\max_KL$. There are a bounded open set $U\supset K$ and a constant $C$ such that
\[
\|A_r(E,x)\|\le Ce^{hr}\qquad\text{for all real }E\in U,\ x\in\T,\ r\ge0.
\]
\end{lemma}

The maximum exists because $L$ is upper semicontinuous.

\begin{proof}
Fix $E_0\in K$ and $h_1=(\max_KL+h)/2$. Since $L(E_0)<h_1$, there is $r_0$ with
\[
\frac1{r_0}\int_\T\log\|A_{r_0}(E_0,x)\|\,dx<h_1.
\]
The function $F_E(x)=\log\|A_{r_0}(E,x)\|$ is jointly continuous, so the same strict inequality holds for $E$ in a compact neighborhood $N_0$ of $E_0$. Rotation by $r_0\alpha$ is uniquely ergodic, so Birkhoff averages of the continuous functions $F_E$ converge to $\int F_E$ uniformly in $x$; by equicontinuity of $\{F_E\}_{E\in N_0}$ they converge uniformly in $E\in N_0$ as well. Hence there is $m_0$ with $\sum_{k<m}F_E(x+kr_0\alpha)\le mr_0h$ for all $m\ge m_0$, $x\in\T$ and $E\in N_0$. Writing a product of length $r=mr_0+s$ with $0\le s<r_0$ as $m$ blocks of length $r_0$ and a remainder, submultiplicativity gives $\|A_r(E,x)\|\le C e^{hr}$ on $N_0$. A finite cover of the compact set $K$ by the interiors of such neighborhoods gives $U$ and $C$.
\end{proof}

For a convergent $p/q$ and a phase $\varphi$, the \emph{rational discriminant} is $\Delta_\varphi(E)=\tr A^{(p/q)}_q(E,\varphi)$, the discriminant of the $q$-periodic operator $H_{v,p/q,\varphi}$. Its spectrum is $\{E:|\Delta_\varphi(E)|\le2\}$. The shift $\varphi\mapsto\varphi+p/q$ permutes the $q$ factors of $A^{(p/q)}_q$ cyclically and leaves the trace unchanged. Since $\gcd(p,q)=1$, $\Delta_\varphi$ is invariant under $\varphi\mapsto\varphi+1/q$.

\begin{lemma}[Rational products]\label{lem:ratprod}
Let $v$ be Lipschitz, and let $U$, $h$ and $C$ be as in Lemma~\ref{lem:growth}. There is $C'$ such that for every convergent $p/q$ of $\alpha$, every real $E\in U$ and all phases $x,\varphi,\varphi'$:
\begin{enumerate}
\item $\|A^{(p/q)}_r(E,x)\|\le C'e^{hr}$ for $0\le r\le q$;
\item $|\partial^r_E\Delta_\varphi(E)|\le 2C'^{r+1}q^re^{hq}$ for every integer $r\ge1$, and
\[
|\Delta_\varphi(E)-\Delta_{\varphi'}(E)|\le 2C'^2\Lip(v)\,q\,e^{hq}|\varphi-\varphi'|;
\]
\item $\|A_q(E,x)-A^{(p/q)}_q(E,x)\|\le C'^2\Lip(v)\,q\,e^{hq}|\delta|$.
\end{enumerate}
\end{lemma}

\begin{proof}
1. For $0\le k<q$ the potentials $v(x+kp/q)$ and $v(x+k\alpha)$ differ by at most $d:=\Lip(v)\,k|\alpha-p/q|\le\Lip(v)|\delta|$. Write each rational factor as the irrational one plus an error of norm at most $d$, and expand the product of $r\le q$ factors. A term with $j$ errors contains at most $j+1$ maximal irrational segments, each bounded by Lemma~\ref{lem:growth}. Hence
\[
\|A^{(p/q)}_r(E,x)\|\le\sum_{j=0}^r\binom rj C^{j+1}e^{h(r-j)}d^j=Ce^{hr}\bigl(1+Cde^{-h}\bigr)^r\le Ce^{hr}e^{C\Lip(v)e^{-h}},
\]
because $rd\le q\Lip(v)|\delta|\le\Lip(v)$.

2. Each factor is affine in $E$ with derivative $\bigl(\begin{smallmatrix}1&0\\0&0\end{smallmatrix}\bigr)$. Differentiating the product $r$ times inserts this matrix at $r$ distinct positions; there are at most $q^r$ ordered choices, and each resulting product consists of at most $r+1$ rational segments of total length at most $q$, so it is bounded by $C'^{r+1}e^{hq}$ by part~1. The trace is at most twice the norm. For the phase, telescope the difference of the two products: the term with index $k$ is a product of two rational segments of total length $q-1$ and a one-step difference of norm at most $\Lip(v)|\varphi-\varphi'|$.

3. Telescope as in part~2. The one-step differences are at most $\Lip(v)|\delta|$ by part~1, and the segments are bounded by Lemma~\ref{lem:growth} and part~1.
\end{proof}

\subsection{Rational spectra}

For a rational $p/q$ in lowest terms let $\Sigma_{p/q}=\bigcup_{\varphi\in\T}\sigma(H_{v,p/q,\varphi})$.

\begin{lemma}\label{lem:window}
Let $H_0$ be a bounded Schr\"odinger operator on $\ell^2(\Z)$ with hopping one and real potential, and let $E\in\sigma(H_0)$. For every integer $M\ge1$ there is a unit vector $u$ supported in an interval of length $2M$ with $\|(H_0-E)u\|\le 8/M$.
\end{lemma}

\begin{proof}
By Weyl's criterion there is a finitely supported vector $w$ with $\|(H_0-E)w\|\le M^{-1}\|w\|$. For $k\in\Z$ let $\chi_k(n)=\sin\bigl(\pi(n-kM)/(2M)\bigr)$ for $kM\le n\le (k+2)M$ and $\chi_k(n)=0$ otherwise. Then $\sum_k\chi_k(n)^2=1$, because on $[kM,(k+1)M]$ the two windows $\chi_k$ and $\chi_{k-1}$ are a sine and a cosine of the same angle. Also $|\chi_k(n+1)-\chi_k(n)|\le\pi/(2M)$, and for each $n$ at most three indices $k$ have $\chi_k(n)$ or $\chi_k(n+1)$ nonzero. Since $([S,\chi_k]w)(n)=(\chi_k(n+1)-\chi_k(n))w(n+1)$, this gives $\sum_k\|[S^{\pm1},\chi_k]w\|^2\le 3\pi^2(4M^2)^{-1}\|w\|^2$. Hence $\sum_k\|\chi_kw\|^2=\|w\|^2$ and
\[
\sum_k\|(H_0-E)\chi_kw\|^2\le2\sum_k\|\chi_k(H_0-E)w\|^2+2\sum_k\|[H_0,\chi_k]w\|^2\le\frac{2+6\pi^2}{M^2}\,\|w\|^2.
\]
Some $k$ therefore has $\|(H_0-E)\chi_kw\|\le (8/M)\|\chi_kw\|$.
\end{proof}

\begin{lemma}\label{lem:ratspec}
Let $v$ be Lipschitz. There is $C$ such that for every rational $p/q$:
\begin{enumerate}
\item every point of $\Sigma_{p/q}$ is within $C|\alpha-p/q|^{1/2}$ of $\Sigma$, and every point of $\Sigma$ is within $C|\alpha-p/q|^{1/2}$ of $\Sigma_{p/q}$;
\item for every $\varphi$, every point of $\Sigma_{p/q}$ is within $\Lip(v)/(2q)$ of $\sigma(H_{v,p/q,\varphi})$.
\end{enumerate}
\end{lemma}

\begin{proof}
1. Let $E\in\sigma(H_{v,p/q,\varphi})$ and take $u$ from Lemma~\ref{lem:window}, supported in $[c-M,c+M]$. Choose $\theta'$ with $\theta'+c\alpha=\varphi+cp/q$. For $|n-c|\le M$ the potentials of $H_{v,\alpha,\theta'}$ and $H_{v,p/q,\varphi}$ differ by at most $\Lip(v)M|\alpha-p/q|$. Hence $\dist(E,\Sigma)\le 8/M+\Lip(v)M|\alpha-p/q|$, and $M=\lceil|\alpha-p/q|^{-1/2}\rceil$ gives the first claim. The second claim is proved in the same way with the roles exchanged.

2. Translation by $k$ sites conjugates $H_{v,p/q,\varphi}$ to $H_{v,p/q,\varphi+kp/q}$, and $\{kp/q\}$ is the full grid $q^{-1}\Z/\Z$. Every phase $\varphi'$ is within $1/(2q)$ of a grid translate of $\varphi$, and the potential is $\Lip(v)$-Lipschitz in the phase.
\end{proof}

\subsection{Quasi-locality}

\begin{lemma}\label{lem:QL}
Let $H_0$ be a bounded Schr\"odinger operator on $\ell^2(\Z)$ with hopping one and real potential, let $R>0$, and let $g\in C^\infty(\R)$ be supported in $[-R,R]$. For every integer $M\ge0$ there is $C_{M,R}$, independent of $H_0$, such that
\[
\|\ad_X^M(g(H_0))\|\le C_{M,R}\|g\|_{C^{M+2}},\qquad |\langle\delta_j,g(H_0)\delta_k\rangle|\le C_{M,R}\|g\|_{C^{M+2}}\,|j-k|^{-M}\quad(j\ne k).
\]
\end{lemma}

\begin{proof}
Write $g(H_0)=(2\pi)^{-1}\int\hat g(t)e^{itH_0}\,dt$ and $H_u=e^{iuX}H_0e^{-iuX}$, which replaces $S^{\pm1}$ by $e^{\mp iu}S^{\pm1}$. Then $\ad_X^M(e^{itH_0})=(-i\partial_u)^Me^{itH_u}|_{u=0}$, and all $u$-derivatives of $H_u$ have norm at most $2$. Repeated Duhamel expansion gives $\|\ad_X^M(e^{itH_0})\|\le C_M(1+|t|)^M$. Hence $\|\ad^M_X(g(H_0))\|\le C_M\int(1+|t|)^M|\hat g(t)|\,dt\le C_{M,R}\|g\|_{C^{M+2}}$. The second bound follows from $\langle\delta_j,\ad_X^M(Y)\delta_k\rangle=(j-k)^M\langle\delta_j,Y\delta_k\rangle$.
\end{proof}

\section{From two estimates to the transport law}\label{sec:covering}

The upper bounds of this paper have one architecture. Two estimates about a single operator suffice, and the rest of the argument is a covering of all large times. This section states the two estimates and proves the covering. Section~\ref{sec:typeI} establishes them for supercritical type~I potentials, Section~\ref{sec:moments} for packets filtered in energy, and Section~\ref{sec:beyond} under a separation hypothesis.

Throughout this section $H=H_{v,\alpha,\theta}$ has a bounded real sampling function $v$, a fixed irrational frequency $\alpha$ with $\beta(\alpha)<\infty$, and a fixed phase $\theta$. The initial state $\psi$ is a fixed unit vector with finite moments, and we write $P=P_\psi$. Let $\ell>0$.

\begin{definition}\label{def:M}
We say that $H$ and $\psi$ satisfy:
\begin{itemize}
\item[(M1)] if there are $C_D\ge1$ and $\kappa>0$, and for every $M>0$ a constant $C_M$, such that for all sufficiently large $n$, every $T\ge1$ and every $N\ge C_Dq_n$,
\[
P(N,T)\le C_MN^{-M}+C_MT^3e^{-\kappa q_n};
\]
\item[(M2$_\ell$)] if for every fixed $c>0$, $0<d<B$, $\eps>0$ and $m>0$ there are $C$ and $n_0$ such that for every $n\ge n_0$ with $b_n\ge c$ and every $T$ with $e^{dq_n}\le T\le e^{Bq_n}$,
\[
P\bigl(R_n(T),T\bigr)\le CT^{-m},\qquad R_n(T)=e^{\eps q_n}\Bigl[1+\min\bigl(Te^{-\ell q_n},\,e^{(b_n-\ell)_+q_n}\bigr)\Bigr].
\]
\end{itemize}
\end{definition}

(M1) is a confinement estimate at every denominator. It confines the packet to distance $O(q_n)$ until times $e^{\kappa q_n/3}$, whatever the next denominator is. (M2$_\ell$) describes the motion at a scale with a definite Liouville ratio $b_n\ge c$. Up to the factor $e^{\eps q_n}$, the packet moves with speed $e^{-\ell q_n}$ and stops at the radius $e^{(b_n-\ell)_+q_n}$. In (M2$_\ell$), $n_0$ may depend on the phase. For $\psi=\delta_0$ the per-denominator estimate of Section~\ref{sec:typeI} has the stronger form $P(N,T)\le Ce^{-c_DN}+CT^3e^{-\kappa q_n}$, which implies (M1).

\subsection{Windows at scales with a definite ratio}

\begin{proposition}\label{prop:selected}
Assume (M2$_\ell$). Let $0<a<1$, $m>0$, $c>0$ and $0<d<B$, and let $B_0>\beta(\alpha)$ satisfy $f_0:=\max\{0,1-\ell/B_0\}<a$. There are $C$ and $n_0$ such that for every $n\ge n_0$ with $b_n\ge c$ and every $T$ with $e^{dq_n}\le T\le e^{Bq_n}$,
\[
P(T^a,T)\le CT^{-m}.
\]
\end{proposition}

\begin{proof}
Choose $\eps>0$ with $\eps/d<(a-f_0)/2$, and apply (M2$_\ell$) with these $c,d,B,\eps,m$. Since $\limsup b_n=\beta(\alpha)<B_0$, we may also assume $b_n\le B_0$ for $n\ge n_0$. Write $q=q_n$, $b=b_n$ and $\tau=\log T/q\in[d,B]$. Then
\[
\frac{\log R_n(T)}{\log T}\le\frac{\eps}{d}+\frac{\log2}{dq}+\frac{\max\{0,\min(\tau-\ell,\,(b-\ell)_+)\}}{\tau}.
\]
The last term is at most $f_0$. If $b\le\ell$ it vanishes. If $b>\ell$ it equals $0$ for $\tau\le\ell$, equals $1-\ell/\tau\le1-\ell/b$ for $\ell<\tau\le b$, and equals $(b-\ell)/\tau\le1-\ell/b$ for $\tau\ge b$; and $1-\ell/b\le1-\ell/B_0$. Hence $\log R_n(T)/\log T<(a+f_0)/2+o(1)<a$ for large $n$, so $R_n(T)\le T^a$ and $P(T^a,T)\le P(R_n(T),T)\le CT^{-m}$.
\end{proof}

\subsection{Covering every time}

\begin{theorem}[Covering]\label{thm:covering}
Assume (M1) and (M2$_\ell$), and $\beta(\alpha)<\infty$. For every $a>\max\{0,1-\ell/\beta(\alpha)\}$ and every $m>0$, $P(T^a,T)=O(T^{-m})$ as $T\to\infty$. Consequently
\[
\alpha^+_{u,\psi}(H)\le\max\Bigl\{0,\ 1-\frac{\ell}{\beta(\alpha)}\Bigr\}.
\]
If $\beta(\alpha)=0$, the same conclusion holds assuming (M1) alone.
\end{theorem}

\begin{proof}
It suffices to treat $a<1$; for $a\ge1$ use monotonicity of $P(\cdot,T)$ and some $a'\in(\max\{0,1-\ell/\beta\},1)$. Let $C_D,\kappa$ be as in (M1). Fix $K\ge1$ with $\kappa K>m+3$, and put $d=1/K$ and $c=a/(2K)$. Fix $B_0>\beta(\alpha)$ with $\max\{0,1-\ell/B_0\}<a$, and put $B=\max\{2B_0/a,2d\}$. Let $C$ and $n_0$ be given by Proposition~\ref{prop:selected}.

For large $T$ let $j$ be the least index with $q_j\ge K\log T$. Two cases exhaust all possibilities.

\emph{Case A: $q_j\le T^a/(2C_D)$.} Let $N=\lfloor T^a\rfloor$, so $N\ge T^a/2\ge C_Dq_j$ for large $T$. By (M1) with $M=2m/a$,
\[
P(T^a,T)\le P(N,T)\le C_M2^MT^{-aM}+C_MT^3e^{-\kappa q_j}\le C_M2^MT^{-2m}+C_MT^{3-\kappa K}=O(T^{-m}).
\]

\emph{Case B: $q_j>T^a/(2C_D)$.} Let $q=q_{j-1}$. By the choice of $j$, $q<K\log T$, so $\tau=\log T/q>d$. By~\eqref{eq:bn}, $q_j<D_{j-1}=e^{b_{j-1}q}$, so
\[
a\log T-\log(2C_D)<\log q_j<b_{j-1}q.
\]
Along a sequence of times in this case, $q_j\to\infty$, hence $j\to\infty$ and $q\to\infty$. Dividing by $q<K\log T$ gives $b_{j-1}>a/K-\log(2C_D)/q\ge a/(2K)=c$ for large $T$. Also $b_{j-1}\le B_0$ for large $j$, so $\tau<(B_0+\log(2C_D)/q)/a\le B$. Thus $T\in[e^{dq},e^{Bq}]$ and $b_{j-1}\ge c$, and Proposition~\ref{prop:selected} gives $P(T^a,T)\le CT^{-m}$ once $j-1\ge n_0$.

If $\beta(\alpha)=0$, then $b_n\to0$, so Case~B cannot occur for large $T$ and (M2$_\ell$) is never used. The bound on $\alpha^+_{u,\psi}$ follows from Lemma~\ref{lem:moments}(3), letting $a$ decrease to $\max\{0,1-\ell/\beta\}$.
\end{proof}

The covering uses no information about the scales between $q_{j-1}$ and $q_j$ other than the two displayed inequalities. In particular it allows $\liminf_n\beta_n=0$. Case~A uses no information about the phase.

\section{The lower bound}\label{sec:lower}

The lower bound holds for every phase and for general Lipschitz sampling functions. In this section $v:\T\to\R$ is Lipschitz and $H_\theta=H_{v,\alpha,\theta}$.

\begin{theorem}\label{thm:L}
Let $\alpha$ be irrational with $0<\beta(\alpha)<\infty$. Assume that $L(\cdot)$ is continuous on a neighborhood of a point $E_*\in\Sigma$ with $L(E_*)=\ell:=\min_\Sigma L$. Then for every phase $\theta$,
\[
\au(H_\theta)\ge\max\Bigl\{0,\ 1-\frac{\ell}{\beta(\alpha)}\Bigr\}.
\]
\end{theorem}

For real-analytic $v$ the Lyapunov exponent is continuous~\cite[Theorem~1]{BourgainJitomirskaya2002}, so its minimum over $\Sigma$ is attained and the hypothesis holds at every minimizer.

\subsection{Elliptic blocks}

For a finite sequence of energies $E-v_r$, $0\le r<q$, let $T_r=\begin{pmatrix}E-v_r&-1\\1&0\end{pmatrix}$, $C_0=I$, $C_r=T_{r-1}\cdots T_0$ and $\mathcal A=C_q$. Suppose $\mathcal A$ is real and elliptic, $|\tr\mathcal A|<2$. Choose $B\in SL(2,\R)$ with $B^{-1}\mathcal AB=R_\varrho$, the rotation by angle $\varrho$, and put $u_r=e_1^{\mathsf T}C_rB$, the first row of $C_rB$, and
\[
S_{\mathcal A}=\sum_{r=0}^{q-1}\|u_r\|^2,\qquad N_\kappa=\begin{pmatrix}0&0\\-1&0\end{pmatrix},\qquad J_0=\begin{pmatrix}0&-1\\1&0\end{pmatrix}.
\]

\begin{lemma}\label{lem:blocks}
\begin{enumerate}
\item $S_{\mathcal A}=|\partial_E\tr\mathcal A|/|\sin\varrho|$.
\item For every $z\in\C$ and $0\le r\le q$, with $\mathcal A(E+z)$ and $C_r(E+z)$ the products at the energy $E+z$,
\[
\|B^{-1}\mathcal A(E+z)B\|\le e^{|z|S_{\mathcal A}},\qquad\|C_r(E+z)B\|\le\|C_r(E)B\|\,e^{|z|S_{\mathcal A}}.
\]
\item If $|\tr\mathcal A|\le3/2$ and the entries $v_r$ depend on a parameter $\phi$ with $|\partial_\phi v_r|\le\mathrm{Lip}$, then $\|B^{-1}(\partial_\phi\mathcal A)B\|\le C\,\mathrm{Lip}\,|\partial_E\tr\mathcal A|$ and $\|B^{-1}(\partial_E\mathcal A)B\|\le C|\partial_E\tr\mathcal A|$, with an absolute constant $C$.
\item Let $z_{\mathcal A}$ be the fixed point of $\mathcal A$ in the upper half-plane. Any $B\in SL(2,\R)$ with $B(i)=z_{\mathcal A}$ conjugates $\mathcal A$ to a rotation, and $\|B\|^2=\|B^{-1}\|^2=e^{d(i,z_{\mathcal A})}$, where $d$ is the hyperbolic distance. Moreover $\|\mathcal A\|_{\mathrm{HS}}^2=2+4\sinh^2\bigl(d(i,z_{\mathcal A})\bigr)\sin^2\varrho$. If $|\tr\mathcal A|\le3/2$, then $e^{d(i,z_{\mathcal A})}\le C\|\mathcal A\|$, and $z_{\mathcal A}$ moves with hyperbolic speed at most $C\|B^{-1}\mathcal A'B\|$ under any parameter derivative $\mathcal A'$.
\item Let $w_0,\dots,w_{q-1}\in\C$, and let $C^w_r$ be the products $C_r$ with every $E-v_j$ replaced by $E-v_j+w_j$. Then
\[
C^w_rB=C_rB\prod_{j=r-1}^{0}(I+w_jQ_j),\qquad Q_j=-J_0u_j^{\mathsf T}u_j,
\]
with the factors ordered from $j=r-1$ on the left to $j=0$ on the right. This identity holds for every $B\in SL(2,\R)$, with $u_j=e_1^{\mathsf T}C_jB$. Each $Q_j$ satisfies $Q_j^2=0$ and $\|Q_j\|=\|u_j\|^2$. In particular $\bigl\|\prod_{j<r}(I+w_jQ_j)-I\bigr\|\le\exp\bigl(\sum_{j<r}|w_j|\,\|u_j\|^2\bigr)-1$.
\end{enumerate}
\end{lemma}

\begin{proof}
Since $T_r(E-v_r+w)=T_r(E-v_r)(I+wN_\kappa)$ for every $w\in\C$, induction on $r$ gives $C^w_r=C_r\prod_{j=r-1}^{0}(I+w_jC_j^{-1}N_\kappa C_j)$. Conjugating by $B$ gives part~5 with $Q_j=B^{-1}C_j^{-1}N_\kappa C_jB$. A direct computation, using $M^{-1}=-J_0M^{\mathsf T}J_0$ for $M=C_jB\in SL(2,\R)$, gives $Q_j=-J_0u_j^{\mathsf T}u_j$, a rank-one matrix with $\|Q_j\|=\|u_j\|^2$ and $Q_j^2=0$ because $u_jJ_0u_j^{\mathsf T}=0$. The product bound follows from $\|I+w_jQ_j\|\le1+|w_j|\|u_j\|^2$. Part~2 is the case $w_j=z$, together with $B^{-1}\mathcal AB=R_\varrho$. Differentiating at $z=0$, $B^{-1}(\partial_E\mathcal A)B=R_\varrho\sum_rQ_r$, and the trace of $R_\varrho(-J_0u_r^{\mathsf T}u_r)$ equals $\sin\varrho\,\|u_r\|^2$, because the symmetric part of $R_\varrho J_0$ is $-\sin\varrho\,I$ and the antisymmetric part has zero quadratic form. This proves part~1. A change of $v_r$ at rate $c_r$ gives in the same way $B^{-1}\mathcal A'B=R_\varrho\sum_rc_rQ_r$, so $\|B^{-1}\mathcal A'B\|\le\max_r|c_r|\,S_{\mathcal A}$. When $|\tr\mathcal A|\le3/2$, $|\sin\varrho|\ge\sqrt7/4$, and part~3 follows from part~1. For part~4, recall that every $M\in SL(2,\R)$ satisfies $\|M\|_{\mathrm{HS}}^2=2\cosh d(i,M(i))$; since its singular values are $\sigma$ and $\sigma^{-1}$, this gives $\|M\|^2=e^{d(i,M(i))}$. Apply this to $B$ and $B^{-1}$. If $B(i)=z_{\mathcal A}$, then $B^{-1}\mathcal AB$ fixes $i$ and is a rotation. Conjugating to that rotation and differentiating $\mathcal A(z_{\mathcal A})=z_{\mathcal A}$ shows that the fixed point moves at hyperbolic speed at most $|1-e^{-2i\varrho}|^{-1}\|B^{-1}\mathcal A'B\|$. Finally $\|\mathcal A\|_{\mathrm{HS}}^2=2+4\sinh^2\bigl(d(i,z_{\mathcal A})\bigr)\sin^2\varrho$, which with $|\sin\varrho|\ge\sqrt7/4$ gives $e^{d(i,z_{\mathcal A})}\le C\|\mathcal A\|$.
\end{proof}

\subsection{Proof of Theorem~\ref{thm:L}}

If $\beta(\alpha)\le\ell$ there is nothing to prove, so assume $\ell<\beta(\alpha)$. Fix $h$ with $\ell<h<\beta(\alpha)$.

\emph{Uniform real growth.} By continuity of $L$ near $E_*$ there is a compact neighborhood $U$ of $E_*$ on which $L<h$. Lemma~\ref{lem:growth} gives $C_h$ with $\|A_r(E,x)\|\le C_he^{hr}$ for $E\in U$, $x\in\T$ and $r\ge0$. By the product rule, with $D_q=Cqe^{hq}$,
\begin{equation}\label{eq:Dq}
|\partial_E\tr A_q(E,x)|\le D_q,\qquad|\partial_x\tr A_q(E,x)|\le\Lip(v)\,D_q\qquad(E\in U),
\end{equation}
the second almost everywhere in $x$.

\emph{A band near $E_*$.} Choose convergents with $\beta_n\to\beta(\alpha)$; write $q=q_n$, $\delta=\delta_n$. Fix a phase $x_0$ and let $H^{\mathrm{per}}$ be the $q$-periodic operator whose potential repeats $v(x_0+k\alpha)$, $0\le k<q$. Its monodromy is $A_q(E,x_0)$, and it differs in norm by at most $\Lip(v)|\delta|$ from the $p/q$-periodic operator at phase $x_0$. By Lemma~\ref{lem:ratspec}, the phase-union spectrum $\Sigma_{p/q}$ converges to $\Sigma$ and lies within $\Lip(v)/(2q)$ of each fixed-phase spectrum. So $\sigma(H^{\mathrm{per}})$ has a point converging to $E_*$. Every band of a $q$-periodic operator with hopping one has width at most $2\pi/q$: distributing the Bloch phase over the $q$ bonds makes the Floquet matrices $2/q$-Lipschitz in the Bloch phase. Hence, for large $n$, an entire band of $H^{\mathrm{per}}$ lies in $U$. On that band $\tr A_q(E,x_0)$ is monotone and takes all values in $[-2,2]$, so the set
\[
S_q=\{E\text{ in the band}:|\tr A_q(E,x_0)|\le1\}
\]
is an interval of length at least $2/D_q$.

\emph{Slowly drifting blocks.} The blocks of the actual operator starting at sites $mq$, for $\theta=x_0$, are $A^{(m)}(E)=A_q(E,x_0+mq\alpha)=A_q(E,x_0+m\delta)$, since $q\alpha=p+\delta$. Let
\[
K_q=\Bigl\lfloor\frac{c_0}{D_q|\delta|}\Bigr\rfloor,\qquad N_q=qK_q,\qquad T_q=C_0K_qD_q,
\]
with $c_0$ small and $C_0$ large fixed constants. Since $\log(1/|\delta|)=b_nq$ and $b_n\to\beta(\alpha)>h$, $K_q\to\infty$. For $E\in S_q$ and $0\le m<K_q$, \eqref{eq:Dq} gives $|\tr A^{(m)}(E)-\tr A_q(E,x_0)|\le\Lip(v)D_qK_q|\delta|\le\Lip(v)c_0\le1/2$, so every block satisfies $|\tr A^{(m)}|\le3/2$.

By Lemma~\ref{lem:blocks}(3)--(4), the fixed points $z_m$ of the blocks $A^{(m)}$ move a total hyperbolic distance $\sum_md(z_m,z_{m+1})\le C\,\Lip(v)D_qK_q|\delta|\le C$. Choose any $B_m\in SL(2,\R)$ with $B_m(i)=z_m$. Since $B_{m+1}^{-1}B_m$ maps $i$ to a point at distance $d(z_m,z_{m+1})$ from $i$, Lemma~\ref{lem:blocks}(4) gives $\|B_{m+1}^{-1}B_m\|=e^{d(z_m,z_{m+1})/2}$, so $\prod_m\|B_{m+1}^{-1}B_m\|\le e^{C}$. By Lemma~\ref{lem:blocks}(1)--(2), each block at the energy $E+i/T_q$ satisfies $\|B_m^{-1}A^{(m)}(E+i/T_q)B_m\|\le e^{CD_q/T_q}$, and the product of $K_q$ such factors is at most $e^{CK_qD_q/T_q}\le e^{C/C_0}$. The endpoint lifts cost at most $Ce^{hq}$ by Lemma~\ref{lem:blocks}(4). A partial last block is handled by the second inequality in Lemma~\ref{lem:blocks}(2). Hence, writing $\Phi(x,z)$ for the transfer matrix of $H_{x_0}$ from site $0$ to site $x$,
\begin{equation}\label{eq:L7}
\max_{0\le x\le N_q}\|\Phi(x,E+i/T_q)\|\le Ce^{2hq}\qquad(E\in S_q).
\end{equation}

\emph{Resolvent.} Let $G(n;z)=\langle\delta_n,(H_{x_0}-z)^{-1}\delta_0\rangle$. For $n\ge1$ the vector $(G(n+1;z),G(n;z))$ is obtained from $(G(1;z),G(0;z))$ by the transfer matrices starting at site $1$, which are bounded as in~\eqref{eq:L7} after a shift of the starting site that changes only constants. Since these matrices have determinant one and $|G(0;E+i/T)|\ge\operatorname{Im}G(0;E+i/T)\ge c/T$ for bounded $E$,
\[
|G(n;E+i/T_q)|^2+|G(n+1;E+i/T_q)|^2\ge cT_q^{-2}e^{-4hq}\qquad(E\in S_q,\ 0\le n\le N_q).
\]
By Lemma~\ref{lem:parseval}, $a(n,T)=(\pi T)^{-1}\int_\R|G(n;E+i/T)|^2\,dE$. Summing over the pairs of sites between $N_q/2$ and $N_q$ gives
\begin{equation}\label{eq:L8}
P(N_q/2,T_q)\ge c\,N_q\,|S_q|\,T_q^{-3}e^{-4hq}\ge T_q^{-C_*}
\end{equation}
for a finite constant $C_*$ depending only on $h$ and $\beta(\alpha)$, because $\log T_q=b_nq+O(\log q)$ and $|S_q|\ge2/D_q$.

\emph{The exponent.} From $\log N_q=(b_n-h)q+O(\log q)$ and $\log T_q=b_nq+O(\log q)$ we get $\log N_q/\log T_q\to1-h/\beta(\alpha)$. By~\eqref{eq:L8}, $M_p(T_q)\ge(N_q/2)^pP(N_q/2,T_q)\ge c(N_q/2)^pT_q^{-C_*}$, so $\beta^+(p)\ge1-h/\beta(\alpha)-C_*/p$. Letting $p\to\infty$ and then $h\downarrow\ell$ proves the theorem for the phase $x_0$. The phase was arbitrary.
\qed

The mass produced by this argument is polynomially small. Section~\ref{sec:moments} replaces it by a fixed fraction of the packet, which controls every moment.

\section{Supercritical type I potentials}\label{sec:typeI}

This section proves the two inputs (M1) and (M2$_{\ell_*}$) of the covering theorem (Definition~\ref{def:M}, with $\psi=\delta_0$) for real-analytic potentials, where $\ell_*$ is the smallest Lyapunov exponent on the spectrum. Theorem~\ref{thm:tI-M1} proves (M1), in a stronger exponential form, for every real-analytic potential whose Lyapunov exponent is positive on the spectrum. Theorem~\ref{thm:tI-window} is a single-scale window estimate under a geometric hypothesis on the rational discriminant near the phase of the initial site. Theorems~\ref{thm:tI-M2even} and~\ref{thm:tI-M2} verify that hypothesis, and hence prove (M2$_{\ell_*}$), for almost every phase when, in addition, the potential is of type I in the sense of Ge and Jitomirskaya. Theorem~\ref{thm:tI-law} combines them with the covering argument and with Theorem~\ref{thm:L}; it is Theorem~\ref{thm:A}. Theorem~\ref{thm:tI-NP} treats potentials of the form $w(sx)$; it is Theorem~\ref{thm:Cov}.

The proof of (M2) uses no conserved operator. Its single object is the resolvent vector $(H-E-i/T)^{-1}\delta_0$. Comparing $q$ consecutive transfer matrices with the periodic approximant at the local phase gives an exact $q$-step recursion for this vector (Lemma~\ref{lem:tI-Q}). The coefficient of the recursion is the rational discriminant at a phase that drifts by $\delta_n$ per site. Its error is controlled by the Lyapunov exponent at the energy $E$ itself. The type I hypothesis makes the discriminant, as a function of the phase, behave like a single cosine near the phase of the initial site. Where it is large, the recursion forces exponential decay. Where it is small, its imaginary part has a fixed sign, and a weighted estimate bounds the speed.

\subsection{Setting and type I potentials}\label{sec:tI-setting}

In this section only, $v:\T\to\R$ is real analytic. It extends holomorphically to a strip $|\operatorname{Im}x|<\delta_v$, and we fix $0<\delta_1<\delta_v$. We write $H_\theta=H_{v,\alpha,\theta}$ for the operator $(H_\theta u)(n)=u(n+1)+u(n-1)+v(\theta+n\alpha)u(n)$, and define $T(E,x)$, $A_r(E,x)$, $\Sigma$, $P(N,T)$ and $\au$ as in Sections~\ref{sec:intro} and~\ref{sec:prelim} with this $v$. Transfer matrices are used at complex energies $E$ and at complex phases $x$ with $|\operatorname{Im}x|<\delta_v$. We write $\|v'\|=\max_\T|v'|$, so that $\Lip(v)=\|v'\|$, and $\ell(E)=L(E)$ for the Lyapunov exponent of $v$ at $E\in\C$. By~\cite[Theorem~1]{BourgainJitomirskaya2002}, $\ell$ is continuous on $\R$. The operator is \emph{supercritical} if $\ell>0$ on $\Sigma$, and then $\ell_*=\min_\Sigma\ell>0$. The continued-fraction notation $p_n/q_n$, $\delta_n$, $D_n$, $b_n$ is that of Section~\ref{sec:prelim}. At a fixed scale we write $q,p,\delta,D,b$. Constants may depend on $v$, $\alpha$ and explicitly fixed parameters, never on the phase, the scale or the time. ``For large $q$'' means for every convergent denominator $q=q_n$ with $n$ beyond a threshold that depends only on the same data.

\emph{Complexification.} For real $E$ and $|y|<\delta_v$ let $L(E,y)$ be the Lyapunov exponent of the cocycle $x\mapsto T(E,x+iy)$ over the rotation by $\alpha$, so that $L(E,0)=\ell(E)$. Since $v$ is real, $A_r(E,x-iy)$ is the complex conjugate of $A_r(E,x+iy)$ for real $x$, so $L(E,\cdot)$ is even. It is convex~\cite[Section~1.3]{Avila2015}. The \emph{acceleration} is $\omega(E)=\lim_{y\downarrow0}(L(E,y)-\ell(E))/(2\pi y)$. We use three facts from Avila's global theory. Theorem numbers for~\cite{Avila2015} refer to its arXiv version~v1.
\begin{itemize}
\item[(A1)] \emph{Quantization}~\cite[Theorem~5]{Avila2015}. The acceleration of an analytic $SL(2,\C)$ cocycle over an irrational rotation is an integer. Applied to the cocycles at every height, this makes $L(E,\cdot)$ a convex piecewise affine function with slopes in $2\pi\Z$.
\item[(A2)] \emph{Regularity}~\cite[Theorem~6]{Avila2015}. A cocycle is \emph{regular} if its complexified Lyapunov exponent is affine for heights near zero. If its Lyapunov exponent is positive, it is regular if and only if it is uniformly hyperbolic.
\item[(A3)] \emph{Uniform hyperbolicity}~\cite[Section~1.3]{Avila2015}. A real energy $E$ lies outside $\Sigma$ if and only if its cocycle is uniformly hyperbolic. Uniform hyperbolicity is an open condition, uniformly hyperbolic cocycles have positive Lyapunov exponent, and the Lyapunov exponent is a smooth function on the open set of uniformly hyperbolic cocycles.
\end{itemize}

\begin{definition}\label{def:tI}
Ge and Jitomirskaya call an energy of \emph{type I} if its T-acceleration equals one~\cite[Definition~2.1]{GeJitomirskaya2024}. For supercritical energies the T-acceleration is the acceleration $\omega(E)$~\cite[Section~2.1]{GeJitomirskaya2024}. We say that $H_{v,\alpha}$ is \emph{supercritical and type I} if $\ell(E)>0$ and $\omega(E)=1$ for every $E\in\Sigma$.
\end{definition}

Here~\cite{GeJitomirskaya2024} is cited in its arXiv version~v2. The almost Mathieu operator with $\lambda>1$ is supercritical and type I: by~\cite[Theorem~10 and Corollary~11]{Avila2015}, $L(E,y)=\max\{\ell(E),\log\lambda+2\pi|y|\}$ and $\ell(E)=\log\lambda$ for $E\in\Sigma$, so $L(E,y)=\log\lambda+2\pi|y|$ on $\Sigma$. The operator with potential $2\lambda\cos(2\pi sx)$, $s\ge2$, is not of type I: its acceleration on the spectrum is $s$ (Remark~\ref{rem:tI-acc}).

\begin{lemma}[A common strip]\label{lem:tI-T}
Let $H_{v,\alpha}$ be supercritical. It is of type I if and only if there is $h\in(0,\delta_1)$ such that
\begin{equation}\label{eq:tI-strip}
L(E,y)=\ell(E)+2\pi|y|\qquad\text{for all }E\in\Sigma\text{ and }|y|\le h.
\end{equation}
\end{lemma}

\begin{proof}
If~\eqref{eq:tI-strip} holds, then $\omega=1$ on $\Sigma$. Conversely, let $E'\in\Sigma$. By (A1), $L(E',\cdot)$ is affine on $[0,y_1]$ for some $y_1>0$, with slope $2\pi\omega(E')\ge0$. If $\omega(E')=0$, evenness makes $L(E',\cdot)$ constant near $0$. The cocycle of $E'$ is then regular with positive Lyapunov exponent, hence uniformly hyperbolic by (A2), which contradicts $E'\in\Sigma$ by (A3). So $\omega(E')\ge1$ at every $E'\in\Sigma$.

Now assume type I and fix $E\in\Sigma$. Then $L(E,y)=\ell(E)+2\pi y$ on some $[0,y_1]$ with $0<y_1<\delta_1$. Put $y_0=y_1/2$. At height $y_0$ the cocycle is regular with Lyapunov exponent at least $\ell(E)>0$, so it is uniformly hyperbolic by (A2). By (A3) there are $e>0$ and $0<d<y_0$ such that the cocycles of all real $E'\in[E-e,E+e]$ at all heights in $[y_0-d,y_0+d]$ are uniformly hyperbolic. Each of them is regular by (A2), so $L(E',\cdot)$ is affine on $[y_0-d,y_0+d]$, with a slope $\varsigma(E')\in2\pi\Z$ by (A1). The slope $\varsigma(E')=(L(E',y_0+d)-L(E',y_0-d))/(2d)$ is continuous in $E'$ by (A3), so it equals $\varsigma(E)=2\pi$. If moreover $E'\in\Sigma$, the convex function $L(E',\cdot)$ has slope at least $2\pi$ just to the right of $0$ and slope $2\pi$ at $y_0$. Hence its slope is $2\pi$ on $(0,y_0]$, and $L(E',y)=\ell(E')+2\pi y$ for $0\le y\le y_0$. Finitely many intervals $[E-e,E+e]$ cover the compact set $\Sigma$. Let $h$ be the smallest of the corresponding $y_0$. Evenness gives~\eqref{eq:tI-strip}.
\end{proof}

Whenever type I is assumed, $h$ denotes the number of Lemma~\ref{lem:tI-T}.

\subsection{Rational discriminants}\label{sec:tI-disc}

Let $p/q$ be a convergent of $\alpha$. For $|\operatorname{Im}\varphi|<\delta_v$ let $\Delta_\varphi(E)$ be the trace of the rational product $A^{(p/q)}_q(E,\varphi)$ of Section~\ref{sec:prelim}. As noted before Lemma~\ref{lem:ratprod}, $\Delta_\varphi$ is invariant under $\varphi\mapsto\varphi+1/q$; the argument there applies to complex phases as well. So
\[
t_q(x,E)=\Delta_{x/q}(E)
\]
is well defined. It is $1$-periodic and analytic in $x$ on $|\operatorname{Im}x|<q\delta_v$, a monic polynomial of degree $q$ in $E$, and real for real $x$ and $E$. Write $t_q(x,E)=\sum_{k\in\Z}c_{q,k}(E)e^{2\pi ikx}$. For real $E$, $c_{q,-k}(E)=\overline{c_{q,k}(E)}$. If $v$ is even, then $t_q(-x,E)=t_q(x,E)$. Indeed, the potentials at the phase $-x/q$ are those at $x/q$ in reverse cyclic order, and the trace of a product of matrices $\bigl(\begin{smallmatrix}a&-1\\1&0\end{smallmatrix}\bigr)$ is invariant under reversing the order, since conjugation by $\operatorname{diag}(1,-1)$ maps each factor to its transpose. So for even $v$ every $c_{q,k}(E)$ is real.

For the operator $H_\theta$ the \emph{local phase} at site $j$ is
\[
x_j=q(\theta+j\alpha)\bmod1=x_0+j\delta,\qquad x_0=q\theta\bmod 1,
\]
because $q\alpha=p+\delta$.

\begin{lemma}[Band centres]\label{lem:tI-bc}
For real $x$ the roots of $t_q(x,\cdot)$ are $q$ simple real numbers $E_1(x)<\dots<E_q(x)$. Each $E_b$ is real analytic and $1$-periodic, and $|E_b'|\le\|v'\|/q$. If $x$ is a real zero of $t_q(\cdot,E)$ for a real $E$, then
\begin{equation}\label{eq:tI-HF}
|\partial_Et_q(x,E)|\ge\frac{q}{\|v'\|}\,|\partial_xt_q(x,E)|.
\end{equation}
\end{lemma}

\begin{proof}
The roots of $t_q(x,\cdot)$ are the eigenvalues, with multiplicity, of the Floquet matrix of the $p/q$-periodic operator at phase $x/q$ with Bloch phase $\pi/2$, since $2\cos(\pi/2)=0$. This is a $q\times q$ Hermitian matrix, so the roots are real. A double root would give two independent Bloch solutions with multiplier $i$, forcing the real unimodular monodromy matrix to equal $iI$, which is impossible. So the roots are simple, and the implicit function theorem gives analytic branches. By the Hellmann--Feynman formula the phase derivative of a simple eigenvalue is at most the norm of the phase derivative of the Floquet matrix, which is at most $\|v'\|$; since $x=q\varphi$, $|E_b'(x)|\le\|v'\|/q$. Finally, differentiating $t_q(x,E_b(x))=0$ gives $\partial_xt_q=-\partial_Et_q\cdot E_b'$, which is~\eqref{eq:tI-HF}.
\end{proof}

We call $E_b(x)$ the \emph{band centres} at phase $x$. Each lies in the spectrum of the $p/q$-periodic operator at phase $x/q$, and each band of that operator contains exactly one of them.

\begin{lemma}[Location]\label{lem:tI-loc}
There is $C$ such that for every convergent $p/q$:
\begin{enumerate}
\item every band centre lies within $C(|\delta|/q)^{1/2}$ of $\Sigma$;
\item for every real $x$, every real $E$ with $\dist(E,\Sigma)\le1/q$ lies within $C/q$ of $\{E_1(x),\dots,E_q(x)\}$.
\end{enumerate}
\end{lemma}

\begin{proof}
Part~1 of Lemma~\ref{lem:ratspec}, with $|\alpha-p/q|=|\delta|/q$, gives part~1. For part~2, $E$ is within $1/q$ of $\Sigma$, hence within $1/q+C(|\delta|/q)^{1/2}\le C/q$ of the phase-union rational spectrum, hence within a further $\|v'\|/(2q)$ of the spectrum at phase $x/q$ by Lemma~\ref{lem:ratspec}(2). Every band of a $q$-periodic operator with hopping one has width at most $2\pi/q$ (proof of Theorem~\ref{thm:L}), and it contains a band centre.
\end{proof}

\begin{lemma}[Local growth]\label{lem:tI-growth}
Let $a>0$.
\begin{enumerate}
\item There are $\rho_a>0$ and $C_a$ such that
\[
\|A_r(z,x)\|\le C_ae^{(\ell(E)+a)r}
\]
for every real $E$ with $\dist(E,\Sigma)\le\rho_a$, every $z\in\C$ with $|z-E|\le\rho_a$, every $x\in\T$ and every $r\ge0$.
\item If $H_{v,\alpha}$ is supercritical and type I, the same holds with $x$ replaced by $x+iy$, $|y|\le h$, and $\ell(E)$ replaced by $\ell(E)+2\pi|y|$.
\item For all large $q$ the rational products $A^{(p/q)}_r(z,\varphi)$, $r\le q$, obey the bounds of parts~1 and~2 with $C_a$ replaced by $C_a'=C_a\exp(C_a\|v'\|_h)$, where $\|v'\|_h=\max_{|\operatorname{Im}x|\le h}|v'(x)|$ and $h=0$ in part~1.
\end{enumerate}
\end{lemma}

\begin{proof}
We prove part~2; part~1 is the case $h=0$, which does not use type I. Fix $(E_0,y_0)\in\Sigma\times[-h,h]$. By~\eqref{eq:tI-strip}, $L(E_0,y_0)=\ell(E_0)+2\pi|y_0|$. Since the Lyapunov exponent is the infimum of the averages $r^{-1}\int_\T\log\|A_r\|\,dx$, there is $r_0$ with $r_0^{-1}\int_\T\log\|A_{r_0}(E_0,x+iy_0)\|\,dx<\ell(E_0)+2\pi|y_0|+a/4$. The integrand is jointly continuous in $(z,y,x)$. So the same holds with $a/2$ for $(z,y)$ in a compact neighborhood $V_0$ of $(E_0,y_0)$, with $z$ complex. As in the proof of Lemma~\ref{lem:growth}, uniform convergence of Birkhoff averages over the uniquely ergodic rotation by $r_0\alpha$ gives $\|A_r(z,x+iy)\|\le C_0e^{(\ell(E_0)+2\pi|y_0|+3a/4)r}$ on $V_0$. Shrinking $V_0$, continuity of $\ell$ on $\R$ gives $\ell(E_0)+2\pi|y_0|\le\ell(E)+2\pi|y|+a/4$ for real $E$ with $(E,y)\in V_0$. Finitely many interiors of such $V_0$ cover the compact set $\Sigma\times[-h,h]$, so there is $\rho_a>0$ such that every ball of radius $3\rho_a$ centred in $\Sigma\times[-h,h]$ lies in one $V_0$. If $\dist(E,\Sigma)\le\rho_a$ and $|z-E|\le\rho_a$, choose $E'\in\Sigma$ with $|E-E'|\le\rho_a$; then $(E,y)$ and $(z,y)$ lie in the ball of radius $3\rho_a$ about $(E',y)$, hence in one $V_0$.

For part~3, the rational factor $T(z,\varphi+iy+kp/q)$ differs from $T(z,\varphi+iy+k\alpha)$ by at most $d=\|v'\|_h|\delta|$ in norm, for $0\le k<q$. Expanding the product as in the proof of Lemma~\ref{lem:ratprod}(1), a term with $k$ errors has at most $k+1$ maximal irrational segments. The sum is at most $C_ae^{(\ell(E)+2\pi|y|+a)r}(1+C_ad)^r\le C_ae^{(\ell(E)+2\pi|y|+a)r}e^{C_adr}$, because $\ell(E)+a>0$; and $dr\le dq\le\|v'\|_h$.
\end{proof}

\begin{corollary}[Bounds for the discriminant]\label{cor:tI-bounds}
For every $\zeta>0$ and all large $q$ the following holds for every real $E$ with $\dist(E,\Sigma)\le\rho_\zeta/2$, every real $E'$ with $|E'-E|\le1/q$, every real $x$ and every $k\ge0$:
\[
|t_q(x,z)|\le e^{q(\ell(E)+\zeta)}\ \ (|z-E|\le 1/q),\qquad
|\partial_E^k\partial_x^jt_q(x,E')|\le k!\,(2q)^ke^{q(\ell(E)+\zeta)}\ \ (j=0,1).
\]
If $H_{v,\alpha}$ is supercritical and type I, then $|\partial_E^kt_q(x,E')|\le k!\,(2q)^ke^{q(\ell(E)+\zeta)+2\pi|\operatorname{Im}x|}$ for $|\operatorname{Im}x|\le qh$.
\end{corollary}

\begin{proof}
Apply Lemma~\ref{lem:tI-growth} with $a=\zeta/2$. The trace of a $2\times2$ matrix is at most twice its norm. Differentiating $\Delta_\varphi$ in $\varphi$ inserts the matrix $\bigl(\begin{smallmatrix}-v'&0\\0&0\end{smallmatrix}\bigr)$ at one of $q$ positions, between two rational segments whose lengths add up to $q-1$; since $\partial_x=q^{-1}\partial_\varphi$, this gives $|\partial_xt_q(x,z)|\le2\|v'\|C_a'^2e^{(\ell(E)+a)q}$. For large $q$ both bounds are at most $e^{q(\ell(E)+\zeta)}$ when $|z-E|\le2/q$. The derivatives in $E$ follow from Cauchy's estimate on the circle of radius $1/(2q)$ about $E'$. The last statement is proved in the same way from Lemma~\ref{lem:tI-growth}(2).
\end{proof}

Parts 1 and 3 of Lemma~\ref{lem:tI-growth} (with $h=0$) and the first statement of Corollary~\ref{cor:tI-bounds} use only the continuity of $\ell$ on $\R$, the compactness of $\Sigma$ and unique ergodicity. They hold, with the same proofs, when $\Sigma$ is replaced by any compact interval of real energies; Section~\ref{sec:beyond} uses them on such an interval.

\subsection{An energy-local \texorpdfstring{$q$}{q}-step identity}\label{sec:tI-Q}

\begin{lemma}[A $q$-step identity]\label{lem:tI-Q}
Fix $a>0$, and let $\rho_a$ and $C_a'$ be as in Lemma~\ref{lem:tI-growth}. For all large $q$ the following holds. Let $E$ be real with $\dist(E,\Sigma)\le\rho_a/2$, let $|z-E|\le1/q$, let $j\in\Z$, and let $\psi$ satisfy $((H_\theta-z)\psi)(k)=0$ for $j-q+1\le k\le j+q$. Then
\begin{equation}\label{eq:tI-Q}
\psi(j+q)+\psi(j-q)=t_q(x_j,z)\,\psi(j)+e_j,\qquad |e_j|\le\eps_{\mathrm{res}}(E)\bigl(|\psi(j)|+|\psi(j+1)|\bigr),
\end{equation}
where $\eps_{\mathrm{res}}(E)=2C_a'^2\,q\,\|v'\|\,|\delta|\,e^{(\ell(E)+a)q}$. In particular~\eqref{eq:tI-Q} holds for $\psi=(H_\theta-z)^{-1}\delta_0$ whenever $|j|\ge q+1$.
\end{lemma}

\begin{proof}
Write $w_k=(\psi(k+1),\psi(k))^{\mathsf T}$ and $T_k=T(z,\theta+k\alpha)$. The equation at site $k$ reads $w_k=T_kw_{k-1}$. Hence $w_{j+q}=A_+w_j$ and $w_j=A_0w_{j-q}$, with $A_+=T_{j+q}\cdots T_{j+1}$ and $A_0=T_j\cdots T_{j-q+1}$. Let $B_j$ be the $q$-periodic operator with potential $v(\theta+j\alpha+(k-j)p/q)$ at site $k$, and $T^B_k$ its transfer matrices. By periodicity both products $T^B_{j+q}\cdots T^B_{j+1}$ and $T^B_j\cdots T^B_{j-q+1}$ equal one matrix $M$, a period transfer matrix of the $p/q$-periodic operator. Its trace is $t_q(x_j,z)$, because the trace does not depend on which period one starts from and $q(\theta+j\alpha)=x_j\bmod1$. Since $\det M=1$, $M+M^{-1}=(\tr M)I$. So
\[
w_{j+q}+w_{j-q}=t_q(x_j,z)\,w_j+\bigl[(A_+-M)+(A_0^{-1}-M^{-1})\bigr]w_j,
\]
and the second component is~\eqref{eq:tI-Q} with $|e_j|\le(\|A_+-M\|+\|A_0^{-1}-M^{-1}\|)(|\psi(j)|+|\psi(j+1)|)$.

For $|k-j|\le q$ the potentials of $H_\theta$ and $B_j$ at site $k$ differ by at most $\|v'\|\,|k-j|\,|\alpha-p/q|\le\|v'\|\,|\delta|$. Telescoping,
\[
A_+-M=\sum_{i=1}^{q}T^B_{j+q}\cdots T^B_{j+i+1}\,(T_{j+i}-T^B_{j+i})\,T_{j+i-1}\cdots T_{j+1}.
\]
The term with index $i$ contains a rational segment of length $q-i$ and an irrational segment of length $i-1$. By Lemma~\ref{lem:tI-growth}, with each segment bounded according to its own length, it has norm at most $C_a'^2\|v'\|\,|\delta|\,e^{(\ell(E)+a)(q-1)}$. The same bound holds for $A_0^{-1}-M^{-1}$, telescoped through the inverse one-step matrices, because a product of inverses is the inverse of a segment and $\|B^{-1}\|=\|B\|$ in $SL(2,\C)$. Summing over $i$ gives the claim. For $\psi=(H_\theta-z)^{-1}\delta_0$ the equation is homogeneous at every site except $0$, which lies outside $[j-q+1,j+q]$ when $|j|\ge q+1$.
\end{proof}

The bound on $e_j$ uses the length of each segment. Bounding both segments by their maximal length $q$ would give the exponent $2\ell(E)$ in place of $\ell(E)$, which is too large for what follows.

\subsection{The per-denominator estimate}\label{sec:tI-M1}

\begin{lemma}\label{lem:tI-herman}
Let $H_{v,\alpha}$ be supercritical, and put $K_E=\max(4,\|v\|_\infty+3)$. There is $\gamma_v>0$ such that for all large $q$, every $x\in\T$ and every $z$ with $|\operatorname{Re}z|\le K_E$ and $0\le\operatorname{Im}z\le1$,
\[
\max_{1\le r\le2q}\|A_r(z,x)\|\ge e^{\gamma_vq}.
\]
\end{lemma}

\begin{proof}
\emph{Positivity.} The function $\ell$ is continuous on $\R$, positive on $\Sigma$ by assumption and positive off $\Sigma$ by (A3). Let $\gamma>0$ be its minimum on $[-K_E,K_E]$. By the Thouless formula $\ell(z)=\int\log|z-E'|\,dk(E')$, valid at every $z\in\C$~\cite{CraigSimon1983}, where $k$ is the integrated density of states, we have $\ell(E+i\eta)\ge\ell(E)$ for real $E$ and $\eta\ge0$. So $\ell\ge\gamma$ on the rectangle $\mathcal R=\{|\operatorname{Re}z|\le K_E,\,0\le\operatorname{Im}z\le1\}$. Since $\ell(z)=\inf_nn^{-1}\int_\T\log\|A_n(z,x)\|\,dx$, we get $\int_\T\log\|A_n(z,x)\|\,dx\ge n\gamma$ for every $n\ge1$ and $z\in\mathcal R$.

\emph{Bounds.} Let $\Gamma=\log\bigl(K_E+3+\max_{|\operatorname{Im}x|\le\delta_1}|v(x)|\bigr)$, so that $\|T(z,x)\|\le e^\Gamma$ for $z\in\mathcal R$ and $|\operatorname{Im}x|\le\delta_1$. Fix $z\in\mathcal R$ and $n\ge1$, and let $f(x)=\|A_n(z,x)\|_{\mathrm{HS}}^2$ for real $x$. It extends analytically to $|\operatorname{Im}x|<\delta_v$ as $\sum_{i,j}a_{ij}(x)\overline{a_{ij}(\bar x)}$, where $a_{ij}$ are the entries of $A_n(z,\cdot)$, and $|f|\le4e^{2n\Gamma}$ on $|\operatorname{Im}x|\le\delta_1$. On the real line, $\int_\T\log f\ge2n\gamma$ and $\log f\le2n\Gamma+\log2$.

\emph{A threshold.} Let $t_0=e^{n\gamma}$. For every real $u$, $t_0^{-1}\int_{t_0}^{2t_0}\log|u-t|\,dt\ge\log t_0-1-\log2$, the minimum over $u$ being attained at $u=3t_0/2$. By Fubini's theorem there is $t\in[t_0,2t_0]$ with $\int_\T\log|f-t|\ge n\gamma-2$.

\emph{Counting zeros.} Put $F=f-t$; then $|F|\le e^{C_1n}$ on $|\operatorname{Im}x|\le\delta_1$, with $C_1=2\Gamma+2$. Let $a_J=\min(\delta_1,1)/4$ and cover $\T$ by $\lceil1/a_J\rceil$ closed intervals $I_k$ of length $a_J$. Since $\int_\T\log|F|\ge n\gamma-2\ge0$ for $n\ge2/\gamma$ and $\int_{\T\setminus I_k}\log|F|\le C_1n$, we get $\int_{I_k}\log|F|\ge-C_1n$, so some $x_k\in I_k$ has $\log|F(x_k)|\ge-C_1n/a_J$. Jensen's formula on the disk of radius $2a_J$ about $x_k$, which lies in the strip $|\operatorname{Im}x|\le\delta_1$, bounds the number of zeros of $F$ in the disk of radius $a_J$, which contains $I_k$, by $(C_1n+C_1n/a_J)/\log2$. Hence $F$ has at most $C_2n$ real zeros in $\T$, with $C_2=\lceil1/a_J\rceil(1+1/a_J)C_1/\log2$. The centres were chosen where $F$ is not small, so no assumption on the location of zeros is needed.

\emph{A long interval.} Let $G=\{x\in\T:f(x)>t\}$. Off $G$, $\log f\le\log t\le n\gamma+\log2$. Splitting $\int_\T\log f$ gives $2n\gamma\le|G|(2n\Gamma+\log2)+(1-|G|)(n\gamma+\log2)$, so $|G|\ge(n\gamma-\log2)/(2n\Gamma)\ge\gamma/(4\Gamma)$ for $n\ge2\log2/\gamma$. The boundary points of the open set $G$ are zeros of $F$, so $G$ has at most $C_2n+1$ components, and one of them is an arc of length at least $c_4/n$, $c_4=\gamma/(4\Gamma(C_2+1))$.

\emph{Orbit hitting.} Take $n=\lfloor c_4q/8\rfloor$, so that this arc has length at least $8/q$, and $n\ge c_4q/16$ for large $q$. On it $\|A_n(z,\cdot)\|\ge2^{-1/2}\|A_n(z,\cdot)\|_{\mathrm{HS}}>2^{-1/2}t^{1/2}\ge2^{-1/2}e^{n\gamma/2}$. Let $p/q$ be the convergent. The points $x+j\alpha$, $0\le j<q$, are within $1/q$ of a translate of the grid $q^{-1}\Z$, because $|j\alpha-jp/q|<j/q^2<1/q$. Hence they are $3/(2q)$-dense, and some $j<q$ has $x+j\alpha$ in the arc, so $\|A_n(z,x+j\alpha)\|\ge2^{-1/2}e^{n\gamma/2}$.

\emph{Conclusion.} Since $A_{j+n}(z,x)=A_n(z,x+j\alpha)A_j(z,x)$ and $\|M^{-1}\|=\|M\|$ for a $2\times2$ matrix of determinant one,
\[
2^{-1/2}e^{n\gamma/2}\le\|A_{j+n}(z,x)\|\,\|A_j(z,x)\|.
\]
So one of $\|A_j(z,x)\|$ (if $j\ge1$) or $\|A_{j+n}(z,x)\|$ is at least $2^{-1/4}e^{n\gamma/4}$, with $1\le j+n\le2q$. Since $n\ge c_4q/16$, the lemma holds with $\gamma_v=c_4\gamma/128$.
\end{proof}

\begin{theorem}[(M1) for analytic potentials]\label{thm:tI-M1}
Let $v$ be real analytic and $H_{v,\alpha}$ supercritical. There are $C_D\ge1$ and $C,c_D,\kappa>0$ such that for every phase $\theta$, every sufficiently large $n$, every $T\ge1$ and every $N\ge C_Dq_n$,
\[
P(N,T)\le Ce^{-c_DN}+CT^3e^{-\kappa q_n}.
\]
\end{theorem}

\begin{proof}
We use the general transport bound of Damanik and Tcheremchantsev~\cite[Theorem~7]{DT2007}. For a Schr\"odinger operator with bounded potential, choose $K_E\ge4$ with $\sigma(H)\subset[-K_E+1,K_E-1]$; here $K_E=\max(4,\|v\|_\infty+3)$ works, since $\|H_\theta\|\le2+\|v\|_\infty$. Their theorem bounds the right outside probability, for the initial state $\delta_1$ and the same exponential time average, by
\[
C\exp(-cN)+CT^3\int_{-K_E}^{K_E}\Bigl(\max_{1\le k\le N}\bigl\|\Phi(k,E+i/T)\bigr\|^2\Bigr)^{-1}dE,
\]
where $\Phi(k,z)$ is the transfer matrix from site $1$ to site $k$, and the left outside probability by the analogous expression for negative $k$. The constants depend only on $K_E$ and $\|v\|_\infty$. For $T\ge1$ and $|E|\le K_E$, the point $E+i/T$ lies in the rectangle of Lemma~\ref{lem:tI-herman}. The transfer matrices from site $1$ are $A_k(z,\theta+\alpha)$, so the integrand of the right half-line is at most $e^{-2\gamma_vq_n}$ once $N\ge2q_n$. The left half-line is handled by reflection: $(Ru)(n)=u(-n)$ satisfies $RH_{v,\alpha,\theta}R=H_{\tilde v,\alpha,-\theta}$ with $\tilde v(x)=v(-x)$. The function $\tilde v$ is real analytic. Since $R\delta_0=\delta_0$, the two families have the same integrated density of states, hence by the Thouless formula the same Lyapunov exponent; so $H_{\tilde v,\alpha}$ is supercritical, and Lemma~\ref{lem:tI-herman} applies to it. If $Hu=zu$, then $Ru$ solves the reflected equation, so the left transfer matrices of $H_{v,\alpha,\theta}$ are right transfer matrices of $H_{\tilde v,\alpha,-\theta}$, up to conjugation by $\bigl(\begin{smallmatrix}0&1\\1&0\end{smallmatrix}\bigr)$ and one one-step factor of bounded norm. With $\kappa=2\min(\gamma_v,\gamma_{\tilde v})$ both integrands are at most $e^{-\kappa q_n}$ once $N\ge2q_n$. Changing the initial site from $\delta_1$ to $\delta_0$ shifts the phase by $\alpha$ and the radius by one, which is absorbed by enlarging $C_D$ and $C$. Finitely many small $n$ are absorbed into $C$, since $P\le1\le T^3$.
\end{proof}

\begin{corollary}\label{cor:tI-beta0}
Let $v$ be real analytic, $H_{v,\alpha}$ supercritical and $\beta(\alpha)=0$. Then for every phase, every $a>0$ and every $m>0$, $P(T^a,T)=O(T^{-m})$. In particular $\au(H_\theta)=0$ for every $\theta$.
\end{corollary}

\begin{proof}
By Theorem~\ref{thm:tI-M1}, (M1) holds for every phase. When $\beta(\alpha)=0$, Theorem~\ref{thm:covering} needs (M1) alone, and its Case~A uses no information about the phase. So $\au\le a$ for every $a>0$ and every phase.
\end{proof}

The conclusion is known. Damanik and Tcheremchantsev proved it for trigonometric polynomial potentials whose Lyapunov exponent is bounded below by a positive constant at every complex energy~\cite[Theorem~5]{DT2007}. Jitomirskaya and Mavi proved $\beta^+(p)=0$ for Diophantine frequencies and piecewise H\"older potentials whose Lyapunov exponent is continuous and positive on $\R$~\cite[Theorem~1.7]{JitomirskayaMavi2017}, and recorded that a simple modification of their proof gives the same for analytic potentials whenever $\log q_{n+1}/q_n\to0$, that is, $\beta(\alpha)=0$~\cite[Remark~1.12]{JitomirskayaMavi2017}. Their hypotheses hold here, since $\ell$ is continuous and positive on $\Sigma$ and, by (A3), off $\Sigma$. Their upper exponent is defined without time averaging, for the average of the packets started at $\delta_0$ and $\delta_1$, so it bounds ours. Lemma~\ref{lem:tI-herman} and Theorem~\ref{thm:tI-M1} give a quantitative route to this endpoint.

\subsection{The dominant phase coefficient}\label{sec:tI-C3}

We import the finite harmonic structure of rational discriminants from Jitomirskaya and Marx. Their $t_{p/q}(\theta,E)$ and $a_{q,k}(E)$ are our $t_q(q\theta,E)$ and $c_{q,k}(E)$, and their complex phase $\theta+i\eps$ corresponds to our $x+iq\eps$.

\begin{lemma}[{\cite[Lemma~3.5, (3.11)--(3.12)]{JitomirskayaMarx2012}}]\label{lem:tI-JM}
There are $K_0\in\N$ and $c_{\mathrm{JM}}>0$ such that, along the convergents of $\alpha$:
\begin{enumerate}
\item $\bigl|t_q(x,E)-\sum_{|k|\le K_0}c_{q,k}(E)e^{2\pi ikx}\bigr|\le e^{-c_{\mathrm{JM}}q}$ for all $E\in\Sigma$ and $|\operatorname{Im}x|\le q\delta_1$, for large $q$;
\item $L(E,y)=\max\bigl\{0,\max_{|k|\le K_0}\bigl(q^{-1}\log|c_{q,k}(E)|-2\pi ky\bigr)\bigr\}+o(1)$ as $q\to\infty$, uniformly in $E\in\Sigma$ and $|y|\le\delta_1$.
\end{enumerate}
\end{lemma}

Jitomirskaya and Marx prove these statements in their Appendix~A for an arbitrary analytic $SL(2,\C)$ cocycle; see (A.13)--(A.17) for part~1 and (A.31) for part~2. Their uniformity over $E\in\Sigma$ comes from the uniform bound (A.2) on the analytic norms of the cocycles, together with continuity of the Lyapunov exponent on the compact family. Proposition~3.1 of~\cite{JitomirskayaMarx2012} is a different statement, about energies with zero Lyapunov exponent, and is not used. We cite~\cite{JitomirskayaMarx2012} in its arXiv version~v2.

\begin{proposition}[Dominant phase coefficient]\label{prop:tI-C3}
Let $H_{v,\alpha}$ be supercritical and type I. Write $a_{q,k}(E)=q^{-1}\log|c_{q,k}(E)|$. Uniformly in $E\in\Sigma$, as $q\to\infty$ along the convergents,
\begin{enumerate}
\item $a_{q,\pm1}(E)=\ell(E)+o(1)$;
\item $a_{q,k}(E)\le\ell(E)-2\pi(|k|-1)h+o(1)$ for $2\le|k|\le K_0$, and $a_{q,0}(E)\le\ell(E)+o(1)$.
\end{enumerate}
Let $\eta_0=\min(\pi h,c_{\mathrm{JM}})/4$. For large $q$ and every $E\in\Sigma$:
\begin{enumerate}
\setcounter{enumi}{2}
\item on the real line $t_q(x,E)=c_{q,0}(E)+c_{q,1}(E)e^{2\pi ix}+c_{q,-1}(E)e^{-2\pi ix}+r_q(x,E)$ with $|r_q|+|\partial_xr_q|\le|c_{q,1}(E)|e^{-\eta_0q}$;
\item for $y\in\{h/4,h/2\}$ and real $u$, $t_q(u\pm iqy,E)=c_{q,\mp1}(E)\,e^{2\pi qy}e^{\mp2\pi iu}\bigl(1+\epsilon_\pm(u)\bigr)$ with $|\epsilon_\pm(u)|\le e^{-\eta_0q}$.
\end{enumerate}
\end{proposition}

\begin{proof}
For real $E$, $|c_{q,-k}|=|c_{q,k}|$. So for $y\ge0$ Lemma~\ref{lem:tI-JM}(2) reads $F_q(y)=\ell(E)+2\pi y+o(1)$ on $[0,h]$, uniformly on $\Sigma$, where $F_q(y)=\max\{0,\max_{0\le k\le K_0}(a_{q,k}+2\pi ky)\}$. Fix $\epsilon>0$ with $\epsilon<\pi h/8$ and $\epsilon<c_{\mathrm{JM}}/2$, and let $q$ be so large that the error is at most $\epsilon$. Since $F_q(h/2)\ge\ell(E)+\pi h-\epsilon>0$, some line $k_0\in\{0,\dots,K_0\}$ attains $F_q(h/2)$. Lines lie below $F_q$, so $F_q(h/4)\ge F_q(h/2)-\pi k_0h/2$ and $F_q(3h/4)\ge F_q(h/2)+\pi k_0h/2$. With the two-sided bounds on $F_q$ these give $k_0\ge1-4\epsilon/(\pi h)$ and $k_0\le1+4\epsilon/(\pi h)$, so $k_0=1$ and $|a_{q,1}-\ell(E)|\le\epsilon$. This is part~1. For $k\ge2$, $a_{q,k}+2\pi kh\le F_q(h)\le\ell(E)+2\pi h+\epsilon$, and $a_{q,0}\le F_q(0)\le\ell(E)+\epsilon$. This is part~2.

For part~3, the modes $2\le|k|\le K_0$ contribute at most $2K_0(1+2\pi K_0)e^{q(\ell(E)-2\pi h+\epsilon)}$ to $|r_q|+|\partial_xr_q|$. The tail beyond $K_0$ is at most $e^{-c_{\mathrm{JM}}q}$ on $|\operatorname{Im}x|\le q\delta_1$, so its $x$-derivative on the real line is at most $e^{-c_{\mathrm{JM}}q}$ by Cauchy's estimate on the unit disk. Divide by $|c_{q,1}(E)|\ge e^{q(\ell(E)-\epsilon)}$, which is part~1. The ratio is at most $2K_0(1+2\pi K_0)e^{q(-2\pi h+2\epsilon)}+2e^{q(-c_{\mathrm{JM}}-\ell(E)+\epsilon)}$, which is at most $e^{-\eta_0q}$ for large $q$, because $2\epsilon<\pi h/4$, $\epsilon<c_{\mathrm{JM}}/2$ and $\ell(E)>0$. For part~4, on the line $\operatorname{Im}x=qy$ the mode $k$ has modulus $|c_{q,k}|e^{-2\pi kqy}$. The mode $-1$ has size $e^{q(\ell(E)+2\pi y+o(1))}$. By parts~1--2, the mode $0$ and the modes $k\ge1$ are at most $e^{q(\ell(E)+o(1))}$, and the modes $k\le-2$ are at most $e^{q(\ell(E)+2\pi y-2\pi(|k|-1)(h-y)+o(1))}$. With $h/4\le y\le h/2$ every other mode is smaller by a factor $e^{-q(\pi h/2-o(1))}$, and so is the tail. The line $\operatorname{Im}x=-qy$ is symmetric.
\end{proof}

For $\xi>0$ write
\[
\mathcal N_q(\xi)=\{E\in\R:\dist(E,\Sigma)\le e^{-\xi q}\}.
\]

\begin{lemma}[Near the spectrum]\label{lem:tI-near}
Let $H_{v,\alpha}$ be supercritical and type I, $0<\xi<\eta_0$ and $\zeta>0$. For large $q$ and every $E\in\mathcal N_q(\xi)$:
\begin{enumerate}
\item $e^{q(\ell(E)-\zeta)}\le|c_{q,\pm1}(E)|\le e^{q(\ell(E)+\zeta)}$;
\item $t_q(x,E)=c_{q,0}(E)+c_{q,1}(E)e^{2\pi ix}+c_{q,-1}(E)e^{-2\pi ix}+r_q(x,E)$ on the real line, with $|r_q|+|\partial_xr_q|\le|c_{q,1}(E)|e^{-\xi q/2}$;
\item part~4 of Proposition~\ref{prop:tI-C3} holds at $E$ with $|\epsilon_\pm|\le e^{-\xi q/2}$.
\end{enumerate}
\end{lemma}

\begin{proof}
Pick $E'\in\Sigma$ with $|E-E'|\le e^{-\xi q}$. By Corollary~\ref{cor:tI-bounds} with $\zeta=\xi/8$, every point between $E'$ and $E$ satisfies $|\partial_E\partial_x^jt_q|\le2qe^{q(\ell(E')+\xi/8)+2\pi|\operatorname{Im}x|}$ for $j=0,1$ on the real line and for $j=0$ on $|\operatorname{Im}x|\le qh$. By Proposition~\ref{prop:tI-C3}(1), $|c_{q,\pm1}(E')|\ge e^{q(\ell(E')-\xi/8)}$. Hence on the real line
\[
\sup_x|\partial_x^j(t_q(x,E)-t_q(x,E'))|\le2qe^{-3\xi q/4}|c_{q,1}(E')|,\qquad j=0,1.
\]
Fourier coefficients move by at most the supremum, so $|c_{q,k}(E)-c_{q,k}(E')|\le2qe^{-3\xi q/4}|c_{q,1}(E')|$. Together with $|\ell(E)-\ell(E')|=o(1)$ this gives part~1. The remainder $r_q(\cdot,E)-r_q(\cdot,E')$ is a sum of four such differences, so it and its $x$-derivative are at most $C qe^{-3\xi q/4}|c_{q,1}(E')|$. With Proposition~\ref{prop:tI-C3}(3) and $\xi<\eta_0$ this gives part~2. On the lines $\operatorname{Im}x=\pm qy$ the difference $t_q(x,E)-t_q(x,E')$ is at most $2qe^{-3\xi q/4}|c_{q,\mp1}(E')|e^{2\pi qy}$, which gives part~3.
\end{proof}

\subsection{A window estimate under a geometric hypothesis}\label{sec:tI-window}

The dynamical argument uses the potential only through the following property of the discriminant near the phase $x_0=q\theta$ of the initial site. Distances on $\T$ are arc lengths.

\begin{definition}\label{def:tI-W}
Let $\xi,\zeta>0$, let $q$ be a convergent denominator, $x_0\in\T$ and $e^{-\zeta q/4}\le\sigma\le1$. We say that \emph{(W) holds at $(q,x_0,\sigma)$} if for every $E\in\mathcal N_q(\xi)$ the set
\[
Z_E=\bigl\{x\in\T:\ |x-x_0|\le\sigma/16,\ |t_q(x,E)|<\sigma^2e^{q(\ell(E)-\zeta)}\bigr\}
\]
is empty or lies within $\sigma/16$ of the centre of a closed arc $J_E$ of length $\sigma/4$ on which $t_q(\cdot,E)$ is strictly monotone with $|\partial_xt_q(x,E)|\ge\sigma e^{q(\ell(E)-\zeta)}$.
\end{definition}

\begin{theorem}[Window estimate]\label{thm:tI-window}
Let $H_{v,\alpha}$ be supercritical. Fix $c>0$, $0<d<B$, $\eps>0$ and $m>0$, and let
\[
0<\xi<\tfrac1{100}\min\{\eps,c,d,\ell_*,1\},\qquad0<\zeta<\xi/100.
\]
There are $C$ and $n_1$ with the following property. Let $n\ge n_1$ with $b_n\ge c$, and let $\theta$ be a phase such that (W) holds at $(q_n,q_n\theta,\sigma)$, with these $\xi,\zeta$, for some $\sigma\in[e^{-\zeta q_n/4},1]$. Let $T$ satisfy $e^{dq_n}\le T\le e^{Bq_n}$, and write $\psi_E=(H_\theta-E-i/T)^{-1}\delta_0$. Then:
\begin{enumerate}
\item $P(R_n(T),T)\le CT^{-m}$, where $R_n(T)=e^{\eps q_n}\bigl[1+\min\bigl(Te^{-\ell_*q_n},\,e^{(b_n-\ell_*)_+q_n}\bigr)\bigr]$;
\item (per energy) if $E$ is real with $\dist(E,\Sigma)\le e^{-\xi q_n}$, then $\|\un_{\{|j|>N\}}\psi_E\|\le4T^{-m}$ for every
\[
N\ge e^{11\xi q_n}\Bigl[1+\min\bigl(Te^{-\ell(E)q_n},\,e^{(b_n-\ell(E))_+q_n}\bigr)\Bigr];
\]
\item (far energies) if $E$ is real with $g:=\dist(E,\Sigma)\ge e^{-\xi q_n}$, then $\|\un_{\{|j|>N\}}\psi_E\|^2\le64g^{-2}g'^{-1}e^{-g'N/4}$ for every $N\ge0$, where $g'=\min(g,1)$.
\end{enumerate}
\end{theorem}

No upper bound on $b_n$ is needed. Since $q^{-3}\ge e^{-\zeta q/4}$ for large $q$, the range of $\sigma$ contains $[q_n^{-3},1]$ for large $n$. Section~\ref{sec:beyond} uses exponentially small $\sigma$.

\begin{proof}[Proof of Theorem~\ref{thm:tI-window}]
\emph{Notation.} Write $q=q_n$, $b=b_n$, $\delta=\delta_n$, $D=|\delta|^{-1}=e^{bq}$, $x_0=q\theta\bmod1$, $x_j=x_0+j\delta$, $\eta_T=1/T$ and $\mathcal N_q=\mathcal N_q(\xi)$. Put
\[
N_\Omega=\lfloor D\sigma/16\rfloor,\qquad\Omega=[-N_\Omega,N_\Omega]\cap\Z,
\]
so that $|x_j-x_0|\le\sigma/16$ for $j\in\Omega$. By Lemma~\ref{lem:parseval},
\begin{equation}\label{eq:tI-parseval}
P(N,T)=\frac1{\pi T}\int_\R\bigl\|\un_{\{|j|>N\}}\psi_E\bigr\|^2dE,\qquad\psi_E=(H_\theta-E-i\eta_T)^{-1}\delta_0,
\end{equation}
and $\|\psi_E\|\le T$.

\emph{Parameters.} Use Lemma~\ref{lem:tI-Q} with $a=\zeta/2$ and Corollary~\ref{cor:tI-bounds} with $\zeta$. For $E\in\mathcal N_q$ put $\eps_{\mathrm{res}}=\eps_{\mathrm{res}}(E)$ and
\[
\Lambda=e^{3\xi q}(2+2\eps_{\mathrm{res}}),\quad\varpi=\sigma e^{q(\ell(E)-\zeta)},\quad\mu_0=\frac{q\varpi}{\|v'\|},\quad\rho_E=\frac{\mu_0}{32q^2e^{q(\ell(E)+\zeta)}},\quad\eta_*=\frac{8\Lambda}{\mu_0}.
\]
Let $\ell^+$ be the maximum of $\ell$ on the $1$-neighborhood of $\Sigma$. For large $q$, uniformly in $E\in\mathcal N_q$, in the phase and in $b\ge c$:
\begin{itemize}
\item[(E1)] $\ell_*-\zeta\le\ell(E)\le\ell^+$;
\item[(E2)] $\eps_{\mathrm{res}}\le e^{(\ell(E)+\zeta-b)q}$, and $\vartheta=(2+\eps_{\mathrm{res}})/(\Lambda-\eps_{\mathrm{res}})\le e^{-3\xi q}$;
\item[(E3)] $1000\,q\Lambda\le\min\{\mu_0,\ \mu_0\rho_E,\ \sigma^2e^{q(\ell(E)-\zeta)}\}$, $\rho_E\le1/(4q)$ and $4q\eta_*\le1$.
\end{itemize}
(E1) holds since $\mathcal N_q$ shrinks to $\Sigma$. For (E2), $\Lambda-\eps_{\mathrm{res}}\ge e^{3\xi q}(2+\eps_{\mathrm{res}})$. For (E3), $\Lambda\le2e^{3\xi q}+2e^{(3\xi+\ell(E)+\zeta-b)q}$, while $\mu_0\rho_E=\sigma^2e^{q(\ell(E)-3\zeta)}/(32\|v'\|^2)$, $\mu_0\ge q\sigma^2e^{q(\ell(E)-\zeta)}/\|v'\|$ and $\sigma^2\ge e^{-\zeta q/2}$; the claim follows from $3\xi+5\zeta<\ell_*-\zeta$ and $3\xi+5\zeta<b$. Taylor's formula in $\eta$, with $|\partial_E^kt_q(x,E)|\le k!(2q)^ke^{q(\ell(E)+\zeta)}$ from Corollary~\ref{cor:tI-bounds}, gives for $0<\eta\le\eta_*$ and real $x$
\begin{equation}\label{eq:tI-taylor}
|\operatorname{Re}t_q(x,E+i\eta)-t_q(x,E)|\le\Lambda,\qquad|\operatorname{Im}t_q(x,E+i\eta)-\eta\,\partial_Et_q(x,E)|\le\eta\mu_0/4.
\end{equation}
Indeed the real part of the remainder is at most $2e^{q(\ell(E)+\zeta)}(2q\eta)^2\le16\Lambda^2/(\mu_0\rho_E)\le\Lambda$, and its imaginary part at most $2e^{q(\ell(E)+\zeta)}(2q\eta)^3\le\eta\cdot32q\Lambda^2/(\mu_0\rho_E)\le\eta\mu_0/4$, by (E3).

\emph{Far energies.} Let $E\notin\mathcal N_q$ and $g=\dist(E,\Sigma)$, so $g\ge e^{-\xi q}$. Put $g'=\min(g,1)$ and $\mu'=g'/8$. The operator $e^{\pm\mu'X}H_\theta e^{\mp\mu'X}$ differs from $H_\theta$ by at most $e^{\mu'}-e^{-\mu'}\le3g/8$ in norm, so its resolvent at $z=E+i\eta_T$ has norm at most $2/g$. This is the Combes--Thomas bound $|\langle\delta_j,\psi_E\rangle|\le(2/g)e^{-g'|j|/8}$. Summing, $\|\un_{\{|j|>N\}}\psi_E\|^2\le64g^{-2}g'^{-1}e^{-g'N/4}$. For $g\le1$ and $N\ge e^{2\xi q}$ this is at most $64e^{3\xi q}\exp(-e^{\xi q}/4)$, and these $E$ form a set of bounded length. For $g\ge1$ it is at most $64g^{-2}e^{-N/4}$, which is integrable in $E$. Since $R_n(T)\ge e^{\eps q}\ge e^{2\xi q}$, the far energies contribute at most $Ce^{-e^{\xi q}/8}$ to~\eqref{eq:tI-parseval}.

From now on $E\in\mathcal N_q$ is fixed, $z=E+i\eta_T$, $\psi=\psi_E$, and $d_j=t_q(x_j,z)$. Since $\eta_T\le e^{-dq}\le1/q$, Lemma~\ref{lem:tI-Q} gives
\begin{equation}\label{eq:tI-rec}
\psi(j+q)+\psi(j-q)=d_j\psi(j)+e_j,\qquad|e_j|\le\eps_{\mathrm{res}}(|\psi(j)|+|\psi(j+1)|),\qquad|j|\ge q+1.
\end{equation}
If $Z_E=\emptyset$, then $|t_q(x_j,E)|\ge\sigma^2e^{q(\ell(E)-\zeta)}\ge10\Lambda$ for every $j\in\Omega$ by (E3). Otherwise let $J_E$ be the arc of Definition~\ref{def:tI-W}.

\emph{Near-root lemma.} Let $x\in\T$ with $|x-x_0|\le\sigma/16$ and $|t_q(x,E)|<10\Lambda$. We claim that $|\partial_Et_q(x,E')|\ge\mu_0/2$, with a fixed sign, for $|E'-E|\le\rho_E$, and that $t_q(x,\cdot)$ has a root within $20\Lambda/\mu_0\le\rho_E/2$ of $E$. By (E3), $x\in Z_E$, so $x$ lies in the middle half of $J_E$, where $t_q(\cdot,E)$ is monotone with $|\partial_xt_q|\ge\varpi$. Moving from $x$ in the direction in which $|t_q|$ decreases, we meet a zero $x_*$ of $t_q(\cdot,E)$ with $|x-x_*|\le10\Lambda/\varpi$; this distance is less than $\sigma/16$ by (E3), so we stay in $J_E$. By~\eqref{eq:tI-HF} at $(x_*,E)$, $|\partial_Et_q(x_*,E)|\ge q\varpi/\|v'\|=\mu_0$. The mixed derivative is at most $2qe^{q(\ell(E)+\zeta)}$ by Corollary~\ref{cor:tI-bounds}, so
\[
|\partial_Et_q(x,E)-\partial_Et_q(x_*,E)|\le\frac{20q\Lambda e^{q(\ell(E)+\zeta)}}{\varpi}=\frac{\mu_0}{4}\cdot\frac{80\|v'\|\Lambda}{\sigma^2e^{q(\ell(E)-3\zeta)}}\le\frac{\mu_0}4,
\]
by (E3). Hence $|\partial_Et_q(x,E)|\ge3\mu_0/4$. The second energy derivative is at most $8q^2e^{q(\ell(E)+\zeta)}$, so $\partial_Et_q(x,\cdot)$ changes by at most $\mu_0/4$ on $|E'-E|\le\rho_E$. Since $|t_q(x,E)|<10\Lambda$, monotonicity puts a root within $20\Lambda/\mu_0$ of $E$, and $40\Lambda\le\mu_0\rho_E$ by (E3).

\emph{Forbidden decay.} Let $\mathcal A=\{j\in\Omega:|d_j|<\Lambda\}$, $q'=q+1$, and
\[
\mathcal B=\mathcal A\cup[-q',q']\cup\bigl(\Z\setminus[-N_\Omega+q,N_\Omega-q]\bigr).
\]
We claim that $|\psi(j)|\le\vartheta^kT\le e^{-3k\xi q}T$ whenever $\dist(j,\mathcal B)\ge kq$ for an integer $k\ge0$. Let $M(s)=\max\{|\psi(j)|:\dist(j,\mathcal B)\ge s\}$, a maximum over a finite set, with $M(s)=0$ if the set is empty; $M(0)\le\|\psi\|\le T$. Let $s\ge q$ and $\dist(j,\mathcal B)\ge s$. Then $|j|\ge q+1$ and $|d_j|\ge\Lambda$, and the sites $j\pm q$ and $j+1$ are at distance at least $s-q$ from $\mathcal B$. By~\eqref{eq:tI-rec}, $\Lambda|\psi(j)|\le2M(s-q)+\eps_{\mathrm{res}}(|\psi(j)|+M(s-q))$, so $M(s)\le\vartheta M(s-q)$. Iterate and use (E2).

\emph{Cutting.} For an integer $c_+\ge1$, the restriction $\psi_+$ of $\psi$ to $(c_+,\infty)$ solves $(H_+-z)\psi_+=-\psi(c_+)\delta_{c_++1}$, where $H_+$ is the restriction of $H_\theta$ to $\ell^2((c_+,\infty)\cap\Z)$. Hence $\|\un_{(c_+,\infty)}\psi\|\le T|\psi(c_+)|$, and similarly $\|\un_{(-\infty,-c_+)}\psi\|\le T|\psi(-c_+)|$. Put $M_0=q\lceil e^{\xi q}\rceil$ and $N_0=q+1+M_0$. By the forbidden decay, a cut at a site at distance at least $M_0$ from $\mathcal B$ costs at most $e^{-3\xi qe^{\xi q}}T^2$. Since $N_\Omega\ge e^{(b-\zeta/4)q}/16-1$ and $b\ge c>100\xi$, every site of modulus at most $e^{(b-\xi)q}$ is at distance more than $M_0$ from $\Z\setminus[-N_\Omega+q,N_\Omega-q]$.

\emph{Regime I: $\eta_T\ge\eta_*$.} We claim $|d_j|\ge4\Lambda$ for every $j\in\Omega$. The polynomial $t_q(x_j,\cdot)$ is monic with real roots $E_k$, so $|d_j|=\prod_k|z-E_k|$. If $|t_q(x_j,E)|\ge10\Lambda$, then $|d_j|\ge\prod_k|E-E_k|=|t_q(x_j,E)|\ge10\Lambda$. Otherwise the near-root lemma gives a root $E_i$ with $|E-E_i|\le\rho_E/2$ and $|\partial_Et_q(x_j,\cdot)|\ge\mu_0/2$ on $|E'-E_i|\le\rho_E/2$. By the mean value theorem, $g(E')=\prod_{k\ne i}(E'-E_k)=t_q(x_j,E')/(E'-E_i)$ satisfies $|g(E')|\ge\mu_0/2$ there. So $|d_j|\ge|z-E_i|\prod_{k\ne i}|E-E_k|\ge\eta_T\mu_0/2\ge4\Lambda$. Thus $\mathcal A=\emptyset$. The sites $\pm N_0$ are at distance $M_0$ from $[-q',q']$, and the cuts there give $\|\un_{\{|j|>N_0\}}\psi\|\le2e^{-3\xi qe^{\xi q}}T^2$.

\emph{Regime II: $\eta_T<\eta_*$; confinement.} By~\eqref{eq:tI-taylor}, $|d_j|<\Lambda$ forces $|t_q(x_j,E)|<2\Lambda$, so $\mathcal A\subset I$, where
\[
I=\{j\in\Omega:|t_q(x_j,E)|<2\Lambda\},\qquad U=\{j\in\Omega:|t_q(x_j,E)|\le9\Lambda\}.
\]
If $Z_E=\emptyset$ both are empty. Otherwise the points $x_j$, $j\in I\cup U$, lie in $Z_E$ by (E3), hence in the middle half of $J_E$, where $t_q(\cdot,E)$ is strictly monotone with slope at least $\varpi$. Since $j\mapsto x_j$ is affine with step $\delta$ and $N_\Omega|\delta|<1/2$, the sets $I$ and $U$ are intervals of integers, with
\[
\#I\le\frac{4\Lambda}{\varpi|\delta|}+1\le e^{4\xi q}\bigl(1+e^{(b-\ell(E))q}\bigr),\qquad\#U\le\frac{18\Lambda}{\varpi|\delta|}+1\le e^{5\xi q}\bigl(1+e^{(b-\ell(E))q}\bigr).
\]
For these bounds, $\Lambda D/\varpi\le2e^{(3\xi+2\zeta)q}(e^{(b-\ell(E))q}+e^{\zeta q})$ by (E2) and $\sigma\ge e^{-\zeta q/4}$. If $I=\emptyset$ or $\dist(0,I)>N_0+M_0$, put $N_s=N_0$. Otherwise put $N_s=\max\{N_0,\max_{j\in I}|j|\}+M_0$. In both cases $\pm N_s$ are at distance at least $M_0$ from $\mathcal B$, and
\[
N_s\le2N_0+2M_0+\#I\le e^{5\xi q}\bigl(1+e^{(b-\ell(E))_+q}\bigr)\le e^{(b-\xi)q},
\]
because $6\xi<\min(b,\ell(E))$. The cuts at $\pm N_s$ give $\|\un_{\{|j|>N_s\}}\psi\|\le2e^{-3\xi qe^{\xi q}}T^2$.

\emph{Regime II: velocity.} Now assume in addition that $b>\ell(E)+5\xi$, that $T\le T_b:=De^{-3\zeta q}$, and that $I\ne\emptyset$ with $\dist(0,I)\le N_0+M_0$.

\emph{Geometry.} By Corollary~\ref{cor:tI-bounds}, $t_q(x_j,E)$ changes by at most $|\delta|e^{q(\ell(E)+\zeta)}\le e^{-(5\xi-\zeta)q}$ per site. Over $4(N_0+M_0)\le16qe^{\xi q}$ sites this is less than $\Lambda$. So $|t_q(x_j,E)|<3\Lambda$ for $|j|\le3(N_0+M_0)$, and these sites lie in $U$. A site with $|t_q|<3\Lambda$ and a site of $\Omega\setminus U$ are at least $6\Lambda e^{bq}e^{-q(\ell(E)+\zeta)}\ge e^{(8\xi-\zeta)q}$ apart. Since $U$ is an interval of length at most $e^{5\xi q}(1+e^{(b-\ell(E))q})$ containing $0$, every site $j'\notin U$ within distance $q+1$ of $U$ lies in $\Omega$, at distance at least $M_0$ from $\mathcal B$, and $|\psi(j')|\le e^{-3\xi qe^{\xi q}}T$ by the forbidden decay.

\emph{Sign.} The set of real $x$ with $|x-x_0|\le\sigma/16$ and $|t_q(x,E)|\le9\Lambda$ is an interval, because it lies in the middle half of $J_E$, where $t_q(\cdot,E)$ is monotone. By the near-root lemma, $\partial_Et_q(x,E)\ne0$ on this interval, so its sign $\varsigma$ is constant there. The interval contains $x_j$ for every $j\in U$. By~\eqref{eq:tI-taylor},
\[
\varsigma\operatorname{Im}d_j\ge\eta_T\mu_0/2-\eta_T\mu_0/4=:\mu\qquad(j\in U).
\]

\emph{Weighted estimate.} Let $\nu=\min(\mu,1)/(20q)$, $N_*=\lceil C_M\log T/\nu\rceil$ with $C_M$ chosen below, $w(j)=e^{\nu\min(|j|,N_*)}$ and $\phi=\un_Uw\psi$. Define $(\mathcal Qu)(j)=d_ju(j)-u(j+q)-u(j-q)$. Since $\langle\phi,(S^q+S^{-q})\phi\rangle$ is real and $\phi$ is supported in $U$,
\[
\mu\|\phi\|^2\le\varsigma\operatorname{Im}\langle\phi,\mathcal Q\phi\rangle\le\|\phi\|\,\|\un_U\mathcal Q\phi\|.
\]
For $j\in U$ with $|j|\ge q+1$, by~\eqref{eq:tI-rec},
\[
(\mathcal Q\phi)(j)=-w(j)e_j+\sum_{\pm:\,j\pm q\in U}\bigl(w(j)-w(j\pm q)\bigr)\psi(j\pm q)+\sum_{\pm:\,j\pm q\notin U}w(j)\psi(j\pm q).
\]
Since $|\log w(j)-\log w(j')|\le\nu q\le1/20$ for $|j-j'|\le q$, the weight differences contribute at most $2(e^{\nu q}-1)\|\phi\|\le2.1\nu q\|\phi\|\le0.11\mu\|\phi\|$. The residual contributes at most $(1+e^\nu)\eps_{\mathrm{res}}\|\phi\|\le2.1\eps_{\mathrm{res}}\|\phi\|$, plus a boundary term at the right end of $U$. Now $\eps_{\mathrm{res}}/\mu=4T\eps_{\mathrm{res}}/\mu_0\le4T_b\eps_{\mathrm{res}}/\mu_0\le Ce^{-5\zeta q/4}$, using $\sigma\ge e^{-\zeta q/4}$ and $a=\zeta/2$. So the residual contributes at most $0.1\mu\|\phi\|$ for large $q$. The terms in which $j\pm q$ or $j+1$ lies outside $U$ are boundary terms. There are at most $6(q+1)$ of them, each at most $e^{2\nu}T^{C_M}\cdot e^{-3\xi qe^{\xi q}}T$, and their sum is at most $1$ for large $q$, since $T\le e^{Bq}$. For $|j|\le q$ we have $j\in U$, and the remaining source term $w(j)(d_j\psi(j)-\psi(j+q)-\psi(j-q))$ is at most $e^{\nu q}(e^{q(\ell(E)+\zeta)}+2)T$. Since $T\ge e^{dq}$, the source terms have norm at most $T^{C_S}$ with $C_S=3+(\ell^++1)/d$. Altogether $\mu\|\phi\|\le0.21\mu\|\phi\|+T^{C_S}+1$. Since $1/\mu=4T/\mu_0\le T/2$,
\[
\|\phi\|\le2T^{C_S+1}.
\]
Take $C_M=C_S+m+3$. For $j\in U$ with $|j|\ge N_*$, $|\psi(j)|\le\|\phi\|e^{-\nu N_*}\le2T^{C_S+1-C_M}$. Cut at $c_+=N_*$ if $N_*\in U$, and otherwise at $c_+=1+\max U\le N_*$, which is a boundary site; similarly on the left. The cutting estimate gives $\|\un_{\{|j|>N_*\}}\psi\|\le4T^{-m}$. Finally, $1/\nu=20q\max(1,4T/\mu_0)$ and $\mu_0\ge e^{q(\ell(E)-2\zeta)}/\|v'\|$, so
\[
N_*\le20C_MBq^2\bigl(1+4\|v'\|Te^{-q(\ell(E)-2\zeta)}\bigr)+1\le e^{\xi q}\bigl(1+Te^{-\ell(E)q}\bigr).
\]
The additive polynomial part of this radius is absorbed by the factor $e^{\xi q}$ and the term $1$.

\emph{Assembly.} In every case there is a radius $N(E)$ with $\|\un_{\{|j|>N(E)\}}\psi_E\|^2\le16T^{-2m}$ for large $q$, and
\[
N(E)\le e^{11\xi q}\Bigl[1+\min\bigl(Te^{-\ell(E)q},e^{(b-\ell(E))_+q}\bigr)\Bigr].
\]
In regime I, and in regime II when $I$ is empty or far from the source, $N(E)=N_0\le e^{2\xi q}$. In regime II with $b\le\ell(E)+5\xi$, $N(E)=N_s\le2e^{10\xi q}$. If $b>\ell(E)+5\xi$ and $T\ge T_b$, then $e^{(b-\ell(E))q}=De^{-\ell(E)q}\le e^{3\zeta q}Te^{-\ell(E)q}$, so $N(E)=N_s\le e^{(5\xi+3\zeta)q}[1+\min(Te^{-\ell(E)q},e^{(b-\ell(E))q})]$. If $b>\ell(E)+5\xi$ and $T<T_b$, take $N(E)=\min(N_s,N_*)$. By (E1), $\ell(E)\ge\ell_*-\zeta$, so $N(E)\le e^{(11\xi+\zeta)q}[1+\min(Te^{-\ell_*q},e^{(b-\ell_*)_+q})]\le R_n(T)$, because $11\xi+\zeta<\eps$. Inserting this and the far-energy bound into~\eqref{eq:tI-parseval}, where $\mathcal N_q$ has bounded length, gives $P(R_n(T),T)\le CT^{-m}$. This is part~1. Since $\|\un_{\{|j|>N\}}\psi_E\|$ is nonincreasing in $N$, the bound on $N(E)$ is part~2, and the far-energy step is part~3. All thresholds depend only on $v,\alpha,c,d,B,\eps,m,\xi,\zeta$.
\end{proof}

\subsection{Even potentials}\label{sec:tI-even}

\begin{lemma}\label{lem:tI-Weven}
Let $v$ be even and $H_{v,\alpha}$ supercritical and type I. Let $0<\xi<\eta_0$ and $\zeta>0$. For large $q$ and every $x_0$ with $s=|\sin2\pi x_0|\ge q^{-2}$, (W) holds at $(q,x_0,s)$.
\end{lemma}

\begin{proof}
Let $E\in\mathcal N_q(\xi)$ and $J_E=\{x:|x-x_0|\le s/8\}$. Since $v$ is even, $c_{q,-1}(E)=c_{q,1}(E)$ is real, and Lemma~\ref{lem:tI-near} (with $\zeta/2$) gives $t_q(x,E)=c_{q,0}+2c_{q,1}\cos2\pi x+r_q$ on the real line. On $J_E$, $|\sin2\pi x|\ge s(1-\pi/4)\ge s/5$, so $\sin2\pi x$ has a fixed sign, and $|4\pi c_{q,1}\sin2\pi x|\ge2.5|c_{q,1}|s$. Also $|\partial_xr_q|\le|c_{q,1}|e^{-\xi q/2}\le|c_{q,1}|s/2$. Hence $\partial_xt_q=-4\pi c_{q,1}\sin2\pi x+\partial_xr_q$ has a fixed sign on $J_E$ and $|\partial_xt_q|\ge2|c_{q,1}|s\ge se^{q(\ell(E)-\zeta)}$. The set $Z_E$ lies within $s/16$ of $x_0$, the centre of $J_E$.
\end{proof}

\begin{theorem}[(M2) for even type I potentials]\label{thm:tI-M2even}
Let $v$ be even and $H_{v,\alpha}$ supercritical and type I. Fix $c>0$, $0<d<B$, $\eps>0$ and $m>0$. There are $C$ and $n_1$ such that for every $n\ge n_1$ with $b_n\ge c$, every phase with $|\sin(2\pi q_n\theta)|\ge q_n^{-2}$ and every $T$ with $e^{dq_n}\le T\le e^{Bq_n}$,
\[
P(R_n(T),T)\le CT^{-m}.
\]
In particular (M2) holds for every $\theta\in\Ga$, with $n_0=\max(n_1,n_G(\theta))$.
\end{theorem}

\begin{proof}
Choose $\xi<\min\{\eps,c,d,\ell_*,1,\eta_0\}/100$ and $\zeta<\xi/100$. Lemma~\ref{lem:tI-Weven} gives (W) at $(q_n,q_n\theta,s)$ with $s\in[q_n^{-2},1]$, and Theorem~\ref{thm:tI-window} applies.
\end{proof}

For the almost Mathieu operator $v(x)=2\lambda\cos2\pi x$, $\lambda>1$, Chambers' formula~\cite{Chambers1965} gives $t_q(x,E)=P_{p/q}(E)-2\lambda^q\cos2\pi x$ for a monic polynomial $P_{p/q}$, so the expansion of Lemma~\ref{lem:tI-near} holds exactly, with $c_{q,\pm1}=-\lambda^q$ and $r_q=0$. Theorem~\ref{thm:tI-M2even} then gives (M2$_{\log\lambda}$) on $\Ga$ without the conserved polynomial used in the companion paper~\cite{DMamo}.

\subsection{Potentials that are not even}\label{sec:tI-noneven}

Without evenness, the leading coefficients $c_{q,\pm1}(E)$ are complex conjugates with an energy-dependent argument. The phase at which the discriminant turns therefore depends on the energy, and no single sine condition controls all energies. The substitute is a factorization of the discriminant in the complex phase variable.

\begin{lemma}[Two zeros in a strip]\label{lem:tI-H2}
Let $H_{v,\alpha}$ be supercritical and type I, $0<\xi<\eta_0$ and $\zeta>0$. For large $q$ and every $E\in\mathcal N_q(\xi)$:
\begin{enumerate}
\item $t_q(\cdot,E)$ has exactly two zeros $z_1,z_2$, counted with multiplicity, in each period of the strip $|\operatorname{Im}x|<qh/2$. Both satisfy $|\operatorname{Im}z_k|<qh/4$, and either both are real or $z_2=\bar z_1$.
\item $t_q(x,E)=g_E(x)\sin\pi(x-z_1)\sin\pi(x-z_2)$, where $g_E$ is $1$-periodic, analytic and zero-free on $|\operatorname{Im}x|\le qh/2$, and for real $x$
\[
e^{q(\ell(E)-\zeta)}\le|g_E(x)|\le e^{q(\ell(E)+\zeta)},\qquad|g_E'(x)|\le\zeta|g_E(x)|.
\]
\end{enumerate}
\end{lemma}

\begin{proof}
1. By Lemma~\ref{lem:tI-near}(3), on the line $\operatorname{Im}x=qy$, $y\in\{h/4,h/2\}$, the function $t_q(\cdot,E)$ is a nonzero multiple of $e^{-2\pi iu}$ times $1+\epsilon_+$ with $|\epsilon_+|<1$, so it has no zeros there and winds $-1$ times over a period. On $\operatorname{Im}x=-qy$ it winds $+1$ times. The vertical sides of a period rectangle cancel by periodicity, so the argument principle gives exactly two zeros in $|\operatorname{Im}x|<qy$, for both values of $y$. For real $E$, $t_q(\bar x,E)=\overline{t_q(x,E)}$, so non-real zeros come in conjugate pairs.

2. The product $\sin\pi(x-z_1)\sin\pi(x-z_2)$ is $1$-periodic, so $g_E$ is periodic, analytic and zero-free on the closed strip $|\operatorname{Im}x|\le qh/2$. On the line $\operatorname{Im}x=qh/2$, $\operatorname{Im}(x-z_k)\ge qh/4$, so $\log|\sin\pi(x-z_k)|=\pi(qh/2-\operatorname{Im}z_k)-\log2+O(e^{-\pi qh/2})$. Since $\operatorname{Im}z_1+\operatorname{Im}z_2=0$, the logarithm of the product is $\pi qh-2\log2+o(1)$. By Lemma~\ref{lem:tI-near}(1),(3), $\log|t_q|=\log|c_{q,-1}(E)|+\pi qh+o(1)=q(\ell(E)\pm\zeta/2)+\pi qh$. So $|\log|g_E|-q\ell(E)|\le\zeta q/2+O(1)$ on this line, and likewise on $\operatorname{Im}x=-qh/2$. The function $\log|g_E|-q\ell(E)$ is harmonic on the open strip, continuous up to the boundary and periodic, so by the maximum and minimum principles it is at most $\zeta q$ in modulus on the whole strip. If $u$ is harmonic on a closed disk $D$ of radius $R$ and centre $x_c$, each partial derivative $\partial_iu$ is harmonic, so the mean value property and the divergence theorem give $\partial_iu(x_c)=(\pi R^2)^{-1}\int_{\partial D}u\,\nu_i\,ds$, hence $|\partial_iu(x_c)|\le(2/R)\max_D|u|$. Disks of radius $R=qh/4$ about real points lie in the strip, so $|\nabla\log|g_E||\le8\sqrt2\,\zeta/h$ on the real line, and $|g_E'/g_E|=|\nabla\log|g_E||$ by the Cauchy--Riemann equations. Replacing $\zeta$ by $\zeta h/12$ at the start gives the claim.
\end{proof}

The same count holds on a fixed neighborhood of the spectrum, at every scale. This is needed to define the phase set at every convergent, including scales where $b_n$ is small.

\begin{lemma}[Two zeros at every scale]\label{lem:tI-fixed}
Let $H_{v,\alpha}$ be supercritical and type I. There are a closed neighborhood $U_*$ of $\Sigma$ in $\R$ and $q_*$ such that for every convergent denominator $q\ge q_*$ and every $E\in U_*$, $t_q(\cdot,E)$ has exactly two zeros, counted with multiplicity, in each period of the strip $|\operatorname{Im}x|<qh/4$.
\end{lemma}

\begin{proof}
\emph{Uniform hyperbolicity at positive heights.} For $E\in\Sigma$ and $y\in[h/8,3h/8]$, the function $L(E,\cdot)$ is affine near $y$ by~\eqref{eq:tI-strip}, and $L(E,y)\ge\ell_*>0$. So the cocycle at height $y$ is regular, hence uniformly hyperbolic by (A2). By (A3) and compactness there is $e_0>0$ such that the cocycles of all $E\in U_*=\{E\in\R:\dist(E,\Sigma)\le e_0\}$ at all heights $y\in[h/8,3h/8]$ are uniformly hyperbolic.

\emph{An affine profile off the spectrum.} For $E\in U_*$ each of these cocycles is regular by (A2), so $L(E,\cdot)$ is affine on $[h/8,3h/8]$, with a slope in $2\pi\Z$ by (A1). The slope is continuous in $E$ by (A3), hence constant on each component of $U_*$. Each component meets $\Sigma$, where the slope is $2\pi$. So $L(E,y)=a(E)+2\pi y$ on $[h/8,3h/8]$ for $E\in U_*$, with $a$ continuous and $a(E)+\pi h/4>0$. We do not claim $a(E)=\ell(E)$ off $\Sigma$.

\emph{The finite harmonic structure on $U_*$.} The proofs of (A.13)--(A.17) and (A.31) in~\cite[Appendix~A]{JitomirskayaMarx2012} are written for an arbitrary analytic $SL(2,\C)$ cocycle. Applied to the compact family of cocycles of $E\in U_*$, whose analytic norms are uniformly bounded and whose Lyapunov exponents at the heights $y\in[h/8,3h/8]$ depend continuously on $E$ by (A3), they give the two statements of Lemma~\ref{lem:tI-JM} with $\Sigma$ replaced by $U_*$ and $|y|\le\delta_1$ replaced by $h/8\le y\le3h/8$.

\emph{Dominance.} Let $\epsilon\in(0,\pi h/16)$ and let $q$ be so large that $F(y)=\max\{0,\max_{|k|\le K_0}\ell_k(y)\}$, with $\ell_k(y)=a_{q,k}-2\pi ky$, is within $\epsilon$ of $a(E)+2\pi y$ on $[h/8,3h/8]$. Each line $\ell_k$ lies below $F$, and has slope $-2\pi k$. Comparing at $h/8$ and $3h/8$,
\[
\ell_k(h/4)\le a(E)+\pi h/2+\epsilon-|2\pi+2\pi k|\,h/8.
\]
For $k\ne-1$ this is at most $a(E)+\pi h/2-\pi h/4+\epsilon$, while $F(h/4)\ge a(E)+\pi h/2-\epsilon>0$. So the mode $-1$ exceeds every other mode $|k|\le K_0$ on the line $\operatorname{Im}x=qh/4$ by a factor $e^{q(\pi h/4-2\epsilon)}$, and it exceeds the exponentially small tail as well. On $\operatorname{Im}x=-qh/4$ the mode $+1$ dominates in the same way, since $|c_{q,-k}|=|c_{q,k}|$. The argument principle concludes as in Lemma~\ref{lem:tI-H2}.
\end{proof}

\begin{lemma}[Band-centre curves]\label{lem:tI-curves}
Let $H_{v,\alpha}$ be supercritical and type I. For large $q$:
\begin{enumerate}
\item each band centre $E_b$ has exactly two critical points per period, a maximum and a minimum, and is strictly monotone on the two arcs between them. Let $\mathrm{Turn}_q$ be the set of these $2q$ points;
\item if $E_b(x)=E_b(x')$ with $x\ne x'$, each of the two arcs of $\T$ between $x$ and $x'$ contains a point of $\mathrm{Turn}_q$; in particular $|x-x'|\ge\dist(x,\mathrm{Turn}_q)$;
\item if $x$ is a real zero of $t_q(\cdot,E)$, $E$ real, then the other zero $x'$ in the strip is real, and either $x'=x\in\mathrm{Turn}_q$, or $x'\ne x$ and $E_b(x')=E_b(x)=E$ for the band $b$ with $E_b(x)=E$.
\end{enumerate}
Moreover, let $c>0$, $0<\xi<\min(c/3,\eta_0)$ and $\zeta>0$. For large $q$ with $b\ge c$, every band $b$ and every real $x$,
\begin{enumerate}
\setcounter{enumi}{3}
\item $|\partial_xt_q(x,E_b(x))|\ge2\pi e^{q(\ell(E_b(x))-\zeta)}\dist(x,\mathrm{Turn}_q)$ and $|E_b'(x)|\ge e^{-2\zeta q}\dist(x,\mathrm{Turn}_q)/q$.
\end{enumerate}
\end{lemma}

\begin{proof}
By Lemma~\ref{lem:tI-loc}(1), all band centres lie in $U_*$ for large $q$, so Lemma~\ref{lem:tI-fixed} applies to their values. Near a point $x$ with $E_b(x)=E$, $t_q(x',E)=\prod_{b'}(E-E_{b'}(x'))$ and the factors with $b'\ne b$ do not vanish. So the order of the zero of $t_q(\cdot,E)$ at $x$ equals the order of the zero of $E_b-E$ at $x$, and the total number of solutions of $E_{b'}(x)=E$ over all bands, counted with order, is at most two.

1. $E_b$ is not constant, since otherwise $t_q(\cdot,E)$ would vanish identically. Let $x_M$ and $x_m$ be points of maximum and minimum. Each value strictly between the two extreme values is taken on each of the two arcs between $x_M$ and $x_m$, so it is taken exactly once on each, with order one. Hence $E_b$ is strictly monotone on both arcs with nonzero derivative, and the extreme values are taken only at $x_M$ and $x_m$.

2. By part~1, $x$ and $x'$ lie in the interiors of different monotone arcs, so each arc between them contains $x_M$ or $x_m$.

3. If $x'$ were not real, $\bar x'$ would be a third zero. If $x'=x$, the zero is double, so $\partial_xt_q(x,E)=0$ while $\partial_Et_q(x,E)\ne0$ (Lemma~\ref{lem:tI-bc}); hence $E_b'(x)=0$ and $x\in\mathrm{Turn}_q$. If $x'\ne x$ and $E_{b'}(x')=E$ with $b'\ne b$, then $E$ is taken by $E_b$ either at a second point or with order two at $x$; with $x'$ this gives three zeros.

4. By Lemma~\ref{lem:tI-loc}(1), $E=E_b(x)$ is within $C(|\delta|/q)^{1/2}\le Ce^{-cq/2}$ of $\Sigma$, so $E\in\mathcal N_q(\xi)$. If $x'=x$, then $x\in\mathrm{Turn}_q$ and there is nothing to prove. Otherwise Lemma~\ref{lem:tI-H2} gives $\partial_xt_q(x,E)=\pi g_E(x)\sin\pi(x-x')$, and $|\sin\pi u|\ge2\dist(u,\Z)$, so part~2 gives the first bound. The second follows from $E_b'=-\partial_xt_q/\partial_Et_q$ and $|\partial_Et_q(x,E)|\le2qe^{q(\ell(E)+\zeta)}$.
\end{proof}

\begin{definition}\label{def:tI-good}
Let $\tau_q=q^{-3}$. A point $x_0\in\T$ is \emph{good at scale $q$} if $\dist(x_0,\mathrm{Turn}_q)\ge2\tau_q$. Let
\[
G^v_\alpha=\{\theta\in\T:\ q_n\theta\bmod1\text{ is good at scale }q_n\text{ for all sufficiently large }n\},
\]
and for $\theta\in G^v_\alpha$ let $n_G^v(\theta)$ be an index from which on this holds.
\end{definition}

\begin{lemma}\label{lem:tI-phaseset}
Let $H_{v,\alpha}$ be supercritical and type I. Then $G^v_\alpha$ has full Lebesgue measure. If $v$ is even, then $\Ga\subset G^v_\alpha$.
\end{lemma}

\begin{proof}
The set of points that are not good at scale $q$ is a union of at most $2q$ arcs of length $4\tau_q$, of measure at most $8q^{-2}$. The map $\theta\mapsto q\theta$ preserves Lebesgue measure, and $\sum_nq_n^{-2}<\infty$, so the Borel--Cantelli lemma gives full measure. If $v$ is even, $t_q(-x,E)=t_q(x,E)$, so $E_b(-x)=E_b(x)$ and, by periodicity, $E_b(\frac12+x)=E_b(\frac12-x)$. Hence $0$ and $\frac12$ are critical points of every $E_b$, and by Lemma~\ref{lem:tI-curves}(1) $\mathrm{Turn}_q=\{0,\frac12\}$. Since $|\sin2\pi x_0|\le2\pi\dist(x_0,\{0,\frac12\})$, the condition $|\sin2\pi q_n\theta|\ge q_n^{-2}$ implies goodness once $q_n\ge4\pi$.
\end{proof}

The only place where the phase set must be defined at every convergent is the Borel--Cantelli argument above. If one prefers not to use Lemma~\ref{lem:tI-fixed}, one may instead intersect, over $k\in\N$, the sets of phases that are good at all large $n$ with $b_n\ge1/k$. At those scales the band centres lie in $\mathcal N_q(\xi)$ for $\xi<1/(3k)$, so Lemmas~\ref{lem:tI-H2} and~\ref{lem:tI-curves} suffice. This also gives a full-measure set on which Theorem~\ref{thm:tI-M2} below applies at every selected scale, and it contains $G^v_\alpha$.

\begin{lemma}[Geometry at good phases]\label{lem:tI-NG}
Let $H_{v,\alpha}$ be supercritical and type I, $c>0$, $0<\xi<\min(c/3,\eta_0)$ and $\zeta'>0$. For large $q$ with $b\ge c$, every $x_0$ good at scale $q$ and every $E\in\mathcal N_q(\xi)$, with $\tau=\tau_q$ and $S_0=\{x:|x-x_0|\le\tau/8\}$, at least one of the following holds:
\begin{enumerate}
\item[(i)] $|t_q(x,E)|\ge\tau^2e^{q(\ell(E)-\zeta')}$ for every $x\in S_0$;
\item[(ii)] $t_q(\cdot,E)$ has a simple real zero $z_E$ with $|z_E-x_0|\le3\tau/8$. On $\{|x-z_E|\le\tau/8\}$ it is strictly monotone with $|\partial_xt_q(x,E)|\ge\tau e^{q(\ell(E)-\zeta'/10)}$, and $|t_q(x,E)|\ge2\tau e^{q(\ell(E)-\zeta'/10)}|x-z_E|$ for $x\in S_0$.
\end{enumerate}
\end{lemma}

\begin{proof}
Apply Lemmas~\ref{lem:tI-H2} and~\ref{lem:tI-curves} with $\zeta=\zeta'/20$; for large $q$ every band centre lies in $\mathcal N_q(\xi)$ (Lemma~\ref{lem:tI-loc}(1), $\xi<c/3$). Let $z_1,z_2$ be the zeros of Lemma~\ref{lem:tI-H2}, with $g=g_E$. Distances from complex points to $S_0$ are taken in the metric of the cylinder $(\R/\Z)\times\R$; for complex $w$, $|\sin\pi w|\ge2\dist(w,\Z)$, because $|\sin\pi(u+iy)|^2=\sin^2\pi u+\sinh^2\pi y$.

\emph{Case A: both zeros are at distance at least $\tau/4$ from $S_0$.} Then on $S_0$, $|t_q|\ge e^{q(\ell(E)-\zeta)}(\tau/2)^2$, which gives (i).

\emph{Case B: exactly one zero, $z_1$, is within $\tau/4$ of $S_0$, and $|z_1-z_2|\ge\tau$.} If $z_1$ were not real, $\bar z_1$ would be a zero at the same distance from the real set $S_0$; so $z_1$ is real, and then so is $z_2$. Also $|z_1-x_0|\le3\tau/8$. For real $x$ with $|x-z_1|\le\tau/8$, write $S_k=\sin\pi(x-z_k)$ and $C_k=\cos\pi(x-z_k)$. Then $|S_1|\le\pi\tau/8$, $|C_1|\ge0.99$ and $|S_2|\ge7\tau/4$, so
\[
|\partial_xt_q|=|g'S_1S_2+\pi g(C_1S_2+S_1C_2)|\ge|g|\bigl(5.4\tau-1.24\tau-0.01\tau\bigr)\ge4\tau|g|.
\]
This derivative is continuous and nonzero, hence of fixed sign. For $x\in S_0$, $|x-z_1|\le\tau/2$ and $|x-z_2|\ge\tau/2$, so $|t_q(x,E)|\ge|g|\cdot2|x-z_1|\cdot\tau$. With $|g|\ge e^{q(\ell(E)-\zeta)}$ this is (ii), with $z_E=z_1$.

\emph{Case C: some zero is within $\tau/4$ of $S_0$, and $|z_1-z_2|<\tau$.} Then one zero is within $3\tau/8$ of $x_0$ and the other within $11\tau/8$.

C1. Both zeros real. If $z_1=z_2$, then $z_1\in\mathrm{Turn}_q$ by Lemma~\ref{lem:tI-curves}(3). Otherwise they are two preimages of $E$ under one band centre, and the shorter arc between them, which lies within $11\tau/8$ of $x_0$, contains a point of $\mathrm{Turn}_q$ by Lemma~\ref{lem:tI-curves}(2). Both contradict $\dist(x_0,\mathrm{Turn}_q)\ge2\tau$.

C2. A conjugate pair $z,\bar z$ with $|\operatorname{Im}z|\ge\epsilon_I:=e^{-\zeta'q/4}\tau$. For real $x$, $|\sin\pi(x-z)|\ge2|\operatorname{Im}z|$, so $|t_q|\ge4e^{q(\ell(E)-\zeta)}\epsilon_I^2\ge\tau^2e^{q(\ell(E)-\zeta')}$ on $S_0$. This is (i).

C3. A conjugate pair with $|\operatorname{Im}z|<\epsilon_I$. We show this is impossible. Let $x_1=\operatorname{Re}z$, so $|x_1-x_0|\le3\tau/8$, and $\dist(x,\mathrm{Turn}_q)\ge\tau$ for $|x-x_1|\le\tau/2$. Let $p(u)=t_q(x_1,u)$. Its roots are the band centres $e_i=E_i(x_1)$; put $\ell_i=\ell(e_i)$. First,
\[
|p(E)|=|g(x_1)|\sinh^2(\pi|\operatorname{Im}z|)\le e^{q(\ell(E)+\zeta)}(2\pi\epsilon_I)^2\le4\pi^2e^{q(\ell(E)-0.45\zeta')}\tau^2.
\]
At each root, by~\eqref{eq:tI-HF} and Lemma~\ref{lem:tI-curves}(4),
\[
|p'(e_i)|\ge\mu_i:=\frac{q\tau}{\|v'\|}e^{q(\ell_i-\zeta)}.
\]
By Corollary~\ref{cor:tI-bounds}, $|p''|\le K_i:=8q^2e^{q(\ell_i+\zeta)}$ within $1/q$ of $e_i$. With $\rho_i=\mu_i/(4K_i)$, $|p'|\ge\mu_i/2$ with a fixed sign on $[e_i-\rho_i,e_i+\rho_i]$, so
\[
|p(e_i\pm\rho_i)|\ge\mu_i\rho_i/2=\frac{\tau^2e^{q(\ell_i-3\zeta)}}{64\|v'\|^2}.
\]
These neighborhoods are pairwise disjoint: overlapping neighborhoods of two consecutive roots would contain a critical point of $p$ between them, where $|p'|\ge\mu_i/2$ fails. Suppose $E$ lies outside all of them. If $E$ lies between consecutive roots $e_i<e_{i+1}$, then $\log|p|$ is concave on $[e_i+\rho_i,e_{i+1}-\rho_{i+1}]$, since $(\log|p|)''=-\sum_k(u-e_k)^{-2}$. Interpolate between the two endpoint bounds with their affine weights. By Lemma~\ref{lem:tI-loc}(2) some root is within $C/q$ of $E$, so one endpoint, say $e_i$, is within $C/q$ of $E$, and $\ell_i=\ell(E)+o(1)$. If $e_{i+1}$ is within $q^{-1/2}$ of $E$, then also $\ell_{i+1}=\ell(E)+o(1)$. Otherwise the weight of the far endpoint is $O(q^{-1/2})$, and since $\ell$ is bounded near $\Sigma$, it costs only $O(q^{1/2})$ in $\log|p(E)|$. In both cases
\[
\log|p(E)|\ge q(\ell(E)-4\zeta)+2\log\tau-O(q^{1/2}).
\]
If $E$ lies beyond the largest or below the smallest root, $|p|$ is monotone there and the nearest root, within $C/q$, gives the same bound. With $4\zeta=0.2\zeta'$ this contradicts the upper bound on $|p(E)|$ for large $q$. Hence $E$ lies in some $[e_i-\rho_i,e_i+\rho_i]$, and
\[
|E-e_i|\le\frac{2|p(E)|}{\mu_i}\le C\frac{\tau}{q}e^{-0.35\zeta'q}.
\]
On $\{|x-x_1|\le\tau/2\}$ the band centre $E_i$ is strictly monotone, since this arc contains no turning point, and $|E_i'|\ge e^{-2\zeta q}\tau/q$ by Lemma~\ref{lem:tI-curves}(4). So $E_i$ takes every value within $\tau^2e^{-0.1\zeta'q}/(2q)$ of $e_i$ on this arc. Since $C\tau e^{-0.35\zeta'q}/q$ is smaller for large $q$, $E=E_i(x)$ for some real $x$, a real zero of $t_q(\cdot,E)$. This contradicts the assumption that its two zeros form a non-real conjugate pair.

The cases are exhaustive. If both zeros are within $\tau/4$ of $S_0$, which has length $\tau/4$, they are less than $\tau$ apart, which is Case~C.
\end{proof}

\begin{lemma}\label{lem:tI-Wgood}
Let $H_{v,\alpha}$ be supercritical and type I, $c>0$, $0<\xi<\min(c/3,\eta_0)$ and $\zeta>0$. For large $q$ with $b\ge c$ and every $x_0$ good at scale $q$, (W) holds at $(q,x_0,\tau_q)$.
\end{lemma}

\begin{proof}
Apply Lemma~\ref{lem:tI-NG} with $\zeta'=\zeta$ and let $E\in\mathcal N_q(\xi)$. The arc $\{|x-x_0|\le\tau/16\}$ lies in $S_0$. In case (i), $Z_E=\emptyset$. In case (ii), let $J_E=\{|x-z_E|\le\tau/8\}$, on which $t_q(\cdot,E)$ is strictly monotone with $|\partial_xt_q|\ge\tau e^{q(\ell(E)-\zeta)}$. A point of $Z_E$ lies in $S_0$ and satisfies $2\tau e^{q(\ell(E)-\zeta/10)}|x-z_E|<\tau^2e^{q(\ell(E)-\zeta)}$, so $|x-z_E|<\tau/16$.
\end{proof}

\begin{theorem}[(M2) for type I potentials]\label{thm:tI-M2}
Let $H_{v,\alpha}$ be supercritical and type I. Fix $c>0$, $0<d<B$, $\eps>0$ and $m>0$. There are $C$ and $n_1$ such that for every $n\ge n_1$ with $b_n\ge c$, every phase $\theta$ for which $q_n\theta\bmod1$ is good at scale $q_n$, and every $T$ with $e^{dq_n}\le T\le e^{Bq_n}$,
\[
P(R_n(T),T)\le CT^{-m}.
\]
In particular (M2) holds for every $\theta\in G^v_\alpha$, with $n_0=\max(n_1,n_G^v(\theta))$.
\end{theorem}

\begin{proof}
Choose $\xi<\min\{\eps,c,d,\ell_*,1,\eta_0\}/100$ and $\zeta<\xi/100$. Lemma~\ref{lem:tI-Wgood} gives (W) at $(q_n,q_n\theta,\tau_{q_n})$, and Theorem~\ref{thm:tI-window} applies.
\end{proof}

\subsection{The transport law for type I potentials}\label{sec:tI-law}

\begin{theorem}\label{thm:tI-law}
Let $v$ be real analytic, $\alpha$ irrational, and $H_{v,\alpha}$ supercritical and type I. Let $\ell_*=\min_\Sigma\ell$.
\begin{enumerate}
\item If $\beta(\alpha)<\infty$, then for every $\theta\in G^v_\alpha$,
\[
\au(H_{v,\alpha,\theta})=\max\Bigl\{0,\ 1-\frac{\ell_*}{\beta(\alpha)}\Bigr\},
\]
read as $0$ when $\beta(\alpha)=0$. Moreover $P(T^a,T)=O(T^{-m})$ for every $a>\max\{0,1-\ell_*/\beta(\alpha)\}$ and $m>0$.
\item If $\beta(\alpha)=0$, part~1 holds for every phase. If $\beta(\alpha)=\infty$, then $\au(H_{v,\alpha,\theta})=1$ for every phase.
\item If $0<\beta(\alpha)<\infty$, then $\au(H_{v,\alpha,\theta})\ge1-\ell_*/\beta(\alpha)$ for every phase.
\end{enumerate}
If $v$ is even, $G^v_\alpha\supset\Ga$.
\end{theorem}

\begin{proof}
\emph{Upper bound.} Let $\theta\in G^v_\alpha$ and $\beta=\beta(\alpha)<\infty$. By Theorems~\ref{thm:tI-M1} and~\ref{thm:tI-M2}, $H_{v,\alpha,\theta}$ satisfies (M1) and (M2$_{\ell_*}$). Theorem~\ref{thm:covering} gives $P(T^a,T)=O(T^{-m})$ for every $a>\max\{0,1-\ell_*/\beta\}$ and every $m>0$, and $\au\le\max\{0,1-\ell_*/\beta\}$.

\emph{Lower bound.} The potential $v$ is Lipschitz, and $\ell$ is continuous by~\cite[Theorem~1]{BourgainJitomirskaya2002}, in particular near a point of $\Sigma$ where it attains $\ell_*$. So Theorem~\ref{thm:L} applies with $\ell=\ell_*$ and gives part~3, which with the upper bound gives part~1 when $\beta>\ell_*$. When $\beta\le\ell_*$ the claimed value is $0$.

\emph{Endpoints.} The case $\beta=0$ is Corollary~\ref{cor:tI-beta0}. For $\beta=\infty$, the Lyapunov exponent has no zeros, since it is positive on $\Sigma$ and, by (A3), off $\Sigma$; so the dynamics is quasiballistic, $\beta^+(p)=1$ for every $p>0$, for every phase, by~\cite[Theorem~3(2)]{JZ2022} (numbering of arXiv:1510.07086v2; there it is derived from Theorem~7, which holds for every phase). The moments there are those of~\eqref{eq:abel}, and $\au\le1$ always. The last sentence is Lemma~\ref{lem:tI-phaseset}.
\end{proof}

Theorem~\ref{thm:tI-law} is Theorem~\ref{thm:A}. For the almost Mathieu operator, which is even, supercritical and type~I with $\ell_*=\log\lambda$, it gives $\au=\max\{0,1-\log\lambda/\beta(\alpha)\}$ on the phase set $\Ga$. The companion paper~\cite{DMamo} proves this case independently, through a conserved polynomial.

\subsection{Covariance}\label{sec:tI-cov}

A potential may have a fundamental period shorter than one. If $v(x)=w(sx)$ for an integer $s\ge2$, then $H_{v,\alpha,\theta}=H_{w,s\alpha,s\theta}$, and the relevant arithmetic is that of $s\alpha$.

\begin{proposition}[Chebyshev identity]\label{prop:tI-cheb}
Let $v(x)=w(mx)$ with $w$ real analytic and $m\ge1$ an integer. Let $p/q$ be a convergent of $\alpha$ with $m\mid q$, and put $q'=q/m$. Then $p/q'$ is a convergent of $m\alpha$, with $q'(m\alpha)-p=q\alpha-p$, and for all $x,E$
\[
t^{(v)}_q(x,E)=2\,T_m\bigl(\tfrac12\,t^{(w)}_{q'}(x,E)\bigr),
\]
where $T_m$ is the Chebyshev polynomial, $t^{(v)}_q$ is the discriminant of $v$ at $p/q$, and $t^{(w)}_{q'}$ that of $w$ at $p/q'$, both in the variable $x$ of Section~\ref{sec:tI-disc}.
\end{proposition}

\begin{proof}
Since $\gcd(p,q)=1$, also $\gcd(p,q')=1$. Moreover $|m\alpha-p/q'|=|q\alpha-p|/q'<1/(q'q_{n+1})\le1/(2q'^2)$ when $m\ge2$, because $q_{n+1}>q=mq'$; by Legendre's theorem~\cite[Chapter~II]{Khinchin1964}, $p/q'$ is a convergent of $m\alpha$. For the identity, $v(\varphi+kp/q)=w(m\varphi+kp/q')$, a $q'$-periodic sequence in $k$. So the $q$-step transfer matrix of $v$ at phase $\varphi$ is $M^m$, where $M$ is the $q'$-step transfer matrix of $w$ at phase $m\varphi$. For $M\in SL(2,\C)$, $\tr M^m=2T_m(\tr M/2)$. Finally $q'(m\varphi)=q\varphi=x$.
\end{proof}

At a zero of $t^{(v)}_q(\cdot,E)$, $\frac12t^{(w)}_{q'}$ equals one of the numbers $\cos((2j-1)\pi/(2m))$, and $\partial_xt^{(v)}_q=T_m'(\frac12t^{(w)}_{q'})\,\partial_xt^{(w)}_{q'}$ with $|T_m'|\le m^2$ on $[-1,1]$. By Corollary~\ref{cor:tI-bounds} applied to $w$ at frequency $m\alpha$, whose Lyapunov exponent is that of $v$ at $\alpha$, the phase derivative of $t^{(v)}_q$ at its real zeros near $\Sigma$ is at most $m^2e^{q(\ell(E)+\zeta)/m}$, not of size $e^{q\ell(E)}$. So at scales divisible by $m$ the geometry of Lemma~\ref{lem:tI-NG} fails for $v$. This is how the frequency $s\alpha$, rather than $\alpha$, enters in the next theorem. We do not use, and do not claim, any statement about the individual distances between the clustered zeros.

\begin{lemma}\label{lem:tI-betas}
For every irrational $\alpha$ and integer $s\ge1$, $\beta(\alpha)\le\beta(s\alpha)\le s\beta(\alpha)$.
\end{lemma}

\begin{proof}
For irrational $y$, $\beta(y)=\limsup_{k\to\infty}k^{-1}\log(1/\|ky\|)$, where $\|\cdot\|$ is the distance to $\Z$. Indeed, if $q_n\le k<q_{n+1}$ are the convergent denominators of $y$, the best approximation property gives $\|ky\|\ge\|q_ny\|>1/(q_{n+1}+q_n)\ge1/(2q_{n+1})$~\cite[Chapter~II]{Khinchin1964}, so $k^{-1}\log(1/\|ky\|)\le q_n^{-1}\log(2q_{n+1})$; and $k=q_n$ gives $q_n^{-1}\log(1/\|q_ny\|)>q_n^{-1}\log q_{n+1}$. Now $\|ks\alpha\|\le s\|k\alpha\|$ gives $\beta(s\alpha)\ge\beta(\alpha)$, and $k^{-1}\log(1/\|k(s\alpha)\|)=s\cdot(sk)^{-1}\log(1/\|(sk)\alpha\|)$ gives $\beta(s\alpha)\le s\beta(\alpha)$.
\end{proof}

\begin{theorem}[Covariance]\label{thm:tI-NP}
Let $v$ be a nonconstant real-analytic function on $\T$, let $s\ge1$ be the greatest common divisor of the frequencies $k\ne0$ with $\hat v(k)\ne0$, and let $w$ be the real-analytic function with $v(x)=w(sx)$. Then $H_{v,\alpha,\theta}=H_{w,s\alpha,s\theta}$ for every $\theta$, and the two families have the same spectrum and the same Lyapunov exponent. Assume that $H_{w,s\alpha}$ is supercritical and type I, and let $\ell_*$ be the minimum of its Lyapunov exponent on its spectrum. Then:
\begin{enumerate}
\item if $\beta(s\alpha)<\infty$, then for every $\theta$ with $s\theta\in G^w_{s\alpha}$, a set of full measure,
\[
\au(H_{v,\alpha,\theta})=\max\Bigl\{0,\ 1-\frac{\ell_*}{\beta(s\alpha)}\Bigr\},
\]
read as $0$ when $\beta(s\alpha)=0$;
\item if $\beta(s\alpha)=0$ this holds for every $\theta$, and if $\beta(s\alpha)=\infty$, then $\au(H_{v,\alpha,\theta})=1$ for every $\theta$.
\end{enumerate}
Moreover $\beta(\alpha)\le\beta(s\alpha)\le s\beta(\alpha)$; in particular $\beta(s\alpha)$ is $0$ or $\infty$ exactly when $\beta(\alpha)$ is.
\end{theorem}

\begin{proof}
The function $w(y)=\sum_k\hat v(sk)e^{2\pi iky}$ is real analytic, and $v(\theta+n\alpha)=w(s\theta+ns\alpha)$ for every $n$. The spectra agree because the operators agree. The transfer matrices satisfy $A^{(v)}_r(E,x)=A^{(w,s\alpha)}_r(E,sx)$, and $\int_\T f(sx)\,dx=\int_\T f$, so the Lyapunov exponents agree. The map $\theta\mapsto s\theta$ preserves Lebesgue measure, so the preimage of the full-measure set $G^w_{s\alpha}$ (Lemma~\ref{lem:tI-phaseset}) has full measure. Apply Theorem~\ref{thm:tI-law} to $(w,s\alpha,s\theta)$, and Lemma~\ref{lem:tI-betas}.
\end{proof}

\begin{remark}\label{rem:tI-acc}
In the same way $L_v(E,y)=L_w(E,sy)$, where $L_w$ is computed at frequency $s\alpha$. So the acceleration of $v$ is $s$ times that of $w$, and for $s\ge2$ the operator $H_{v,\alpha}$ is not of type I even when $H_{w,s\alpha}$ is. For example $v(x)=2\lambda\cos(2\pi sx)$ with $\lambda>1$ has acceleration $s$ on its spectrum; Theorem~\ref{thm:tI-NP} gives $\au=\max\{0,1-\log\lambda/\beta(s\alpha)\}$ for almost every phase, and $\beta(s\alpha)$ can exceed $\beta(\alpha)$. For potentials with $s=1$ whose acceleration exceeds one somewhere on the spectrum, Theorem~\ref{thm:tI-M1} and Corollary~\ref{cor:tI-beta0} still apply, but the exact law at finite positive $\beta(\alpha)$ is not proved here; Section~\ref{sec:beyond} proves it under a separation hypothesis.
\end{remark}

\section{Every moment and local dimensions}\label{sec:moments}

Theorem~\ref{thm:A} concerns the farthest polynomial tails. The construction of Section~\ref{sec:lower} moves only a polynomially small part of the packet, which says nothing about a fixed moment. This section shows that a fixed fraction of the packet reaches the critical distance at resonant times, which controls every moment, and it resolves both the transport and the spectral measure energy by energy. Three ingredients enter. An arcsine law says that, at resonant scales, the rational discriminant evaluated along the spectral measure is distributed like the discriminant of a periodic operator. An elliptic-run estimate says that, at an energy where the local transfer matrices are elliptic with a margin, the Weyl solution keeps its size over a long run of periods and carries a fixed fraction of its mass beyond it. A Taylor expansion of energy filters transfers both to packets localized in energy.

\subsection{Notation}\label{sec:mom-notation}

Let $\sigma$ be the arcsine law on $[-2,2]$, with density $\pi^{-1}(4-x^2)^{-1/2}$; its moments are $m_k=\binom{k}{k/2}$ for even $k$ and $m_k=0$ for odd $k$. For $0<\beta<\infty$ and $\ell\in\R$ put $F_\beta(\ell)=\max\{0,1-\ell/\beta\}$; put $F_0\equiv0$ and $F_\infty\equiv1$. For a phase $\theta$ and a convergent $p_n/q_n$ let
\[
\Delta_n(E)=\Delta_\theta(E)=\tr A^{(p_n/q_n)}_{q_n}(E,\theta),
\]
the discriminant of the $p_n/q_n$-periodic operator $H_{v,p_n/q_n,\theta}$. For a finite Borel measure $\mu$ on $\R$ and $E\in\R$, the upper local dimension is $d^+_\mu(E)=\limsup_{r\downarrow0}\log\mu([E-r,E+r])/\log r$. The upper and lower packing dimensions of $\mu$ are the $\mu$-essential supremum and infimum of $d^+_\mu$, written $\dim_P^+(\mu)$ and $\dim_P^-(\mu)$. Recall that $\mu_\theta$ is the spectral measure of $\delta_0$ for $H_{v,\alpha,\theta}$.

\begin{lemma}[Periodic magic formula~\cite{DKS2010}]\label{lem:magic}
Let $H_{\mathrm{per}}$ be a $q$-periodic Schr\"odinger operator on $\ell^2(\Z)$ with hopping one and discriminant $\Delta$. Then $\Delta(H_{\mathrm{per}})=S^q+S^{-q}$.
\end{lemma}

This is the elementary direction of~\cite[Theorem~3.1]{DKS2010}, proved there by Floquet decomposition (equations (3.5)--(3.6)). The hypotheses hold: the operator is bounded, two-sided, periodic and has positive off-diagonal entries. Since $\Delta_n$ is the discriminant of $H_{v,p_n/q_n,\theta}$, the law of $\Delta_n$ under the spectral measure of $\delta_0$ for that periodic operator is exactly $\sigma$: its moments are $\langle\delta_0,(S^q+S^{-q})^k\delta_0\rangle=m_k$. The next two lemmas transfer this to the irrational operator.

\subsection{The arcsine law}\label{sec:arcsine}

\begin{lemma}[Arcsine law at strongly resonant scales]\label{lem:arcsine}
Let $v$ be Lipschitz, $\theta\in\T$, $L_{\max}=\max_\Sigma L$, and let $\mathcal N$ be an infinite set of indices with $b_n\ge c>L_{\max}$ for $n\in\mathcal N$. For every real $g\in C_c^\infty(\R)$ and every integer $k\ge0$,
\[
\int g\,\Delta_n^k\,d\mu_\theta\to m_k\int g\,d\mu_\theta\qquad(n\to\infty,\ n\in\mathcal N).
\]
Consequently, if $g\ge0$, then $\int g\,\phi(\Delta_n)\,d\mu_\theta\to\sigma(\phi)\int g\,d\mu_\theta$ for every bounded continuous $\phi$.
\end{lemma}

\begin{proof}
Fix $h\in(L_{\max},c)$. Let $U\supset\Sigma$ and $C$ be given by Lemma~\ref{lem:growth} with $K=\Sigma$, and $C'$ by Lemma~\ref{lem:ratprod}. Write $q=q_n$, $p=p_n$, $\delta=\delta_n$, $F=S^q+S^{-q}$ and $Q=\Delta_n(H_{v,\alpha,\theta})$. Since $\Delta_n$ is a real polynomial of degree $q$, $Q$ is a real symmetric operator with $\langle\delta_i,Q\delta_j\rangle=0$ for $|i-j|>q$.

\emph{Rows.} Fix $j\in\Z$ with $|j|\le(k+1)q$. Let $\varphi_j=\theta+j(\alpha-p/q)$ and $B_j=H_{v,p/q,\varphi_j}$, whose potential at site $i$ is $v(\theta+j\alpha+(i-j)p/q)$. It agrees with the potential of $H=H_{v,\alpha,\theta}$ at $i=j$ and differs from it by at most $\Lip(v)|i-j||\alpha-p/q|\le\Lip(v)|\delta|$ for $|i-j|\le q$. Let $C_j$ be $B_j$ with its potential on $[j-q,j+q]$ replaced by that of $H$, so $\|B_j-C_j\|\le\Lip(v)|\delta|$. A path of length at most $q$ starting at $j$ stays in $[j-q,j+q]$, so the $j$-th rows of $\Delta_\theta(C_j)$ and $Q$ coincide. The discriminant of $B_j$ is $\Delta_{\varphi_j}$, so by Lemma~\ref{lem:magic} the $j$-th row of $\Delta_{\varphi_j}(B_j)$ is that of $F$. Hence the $j$-th row of $Q-F$ is the $j$-th row of
\[
\bigl[\Delta_\theta(C_j)-\Delta_\theta(B_j)\bigr]+\bigl[\Delta_\theta(B_j)-\Delta_{\varphi_j}(B_j)\bigr].
\]
By Lemma~\ref{lem:ratspec}, $\sigma(B_j)\subset\Sigma_{p/q}$ lies within $C|\delta/q|^{1/2}$ of $\Sigma$, and $\sigma(C_j)$ lies within $\Lip(v)|\delta|$ of $\sigma(B_j)$. Fix an open $U'$ with $\Sigma\subset U'\subset\overline{U'}\subset U$. For large $n$ both spectra lie in $U'$, uniformly in $j$. Since $|\varphi_j-\theta|=|j||\delta|/q$, Lemma~\ref{lem:ratprod}(2) bounds the norm of the second bracket by $2C'^2\Lip(v)e^{hq}|j||\delta|\le2(k+1)C'^2\Lip(v)qe^{hq}|\delta|$. For the first bracket fix $\chi\in C_c^\infty(U)$ with $\chi=1$ on $U'$ and put $f=\chi\Delta_\theta$; then $f(B_j)=\Delta_\theta(B_j)$ and $f(C_j)=\Delta_\theta(C_j)$. For bounded self-adjoint $B,C$, Fourier inversion and $\|e^{itB}-e^{itC}\|\le|t|\,\|B-C\|$ give
\[
\|f(B)-f(C)\|\le\frac{\|B-C\|}{2\pi}\int_\R|t|\,|\hat f(t)|\,dt\le c\,\|B-C\|\,\|f\|_{H^2(\R)},
\]
by the Cauchy--Schwarz inequality and $\int t^2(1+t^2)^{-2}\,dt<\infty$. By Lemma~\ref{lem:ratprod}(2), $\|f\|_{H^2}\le C_\chi q^2e^{hq}$. Altogether each row of $Q-F$ with index $|j|\le(k+1)q$ has $\ell^2$ norm at most $r_n:=C_kq^2e^{hq}|\delta|=C_kq^2e^{(h-b_n)q}$, which tends to zero exponentially.

\emph{Restricted norm.} Let $P_k=\un_{[-kq,kq]}$. The operator $(Q-F)P_k$ has nonzero rows only in $[-(k+1)q,(k+1)q]$, and every row and column of $Q-F$ has at most $2q+1$ nonzero entries. Since $Q-F$ is symmetric, the $\ell^1$ norms of the rows and columns of $(Q-F)P_k$ are at most $\sqrt{2q+1}\,r_n$, and the Schur test gives $\|(Q-F)P_k\|\le\eps_n:=\sqrt{2q+1}\,r_n\to0$.

\emph{Moments.} For $i\le k$ the vectors $Q^i\delta_0$ and $F^i\delta_0$ are supported in $[-iq,iq]$. Let $e_i=\|Q^i\delta_0-F^i\delta_0\|$. For $i<k$, $\|Q\phi\|\le(2+\eps_n)\|\phi\|$ for $\phi$ supported in $[-kq,kq]$, so
\[
e_{i+1}\le\|Q(Q^i\delta_0-F^i\delta_0)\|+\|(Q-F)F^i\delta_0\|\le(2+\eps_n)e_i+\eps_n2^i,
\]
and $e_k\le(2+\eps_n)^k-2^k\to0$. Now $F^k\delta_0=\sum_i\binom{k}{(k+i)/2}\delta_{iq}$, the sum over $|i|\le k$ with $i\equiv k\bmod 2$. Hence
\[
\int g\,\Delta_n^k\,d\mu_\theta=\langle g(H)\delta_0,Q^k\delta_0\rangle=m_k\int g\,d\mu_\theta+\sum_{i\ne0}\binom{k}{(k+i)/2}\langle g(H)\delta_0,\delta_{iq}\rangle+O(\|g\|_\infty e_k).
\]
By Lemma~\ref{lem:QL}, $|\langle\delta_{iq},g(H)\delta_0\rangle|\le C_g|iq|^{-1}\to0$ for $i\ne0$. This proves the first claim.

For $g\ge0$ with $\int g\,d\mu_\theta>0$, the probability measures $(\Delta_n)_*(g\mu_\theta)/\int g\,d\mu_\theta$ have moments converging to those of $\sigma$. The law $\sigma$ has compact support, so it is determined by its moments, and convergence of all moments to those of a moment-determinate law implies weak convergence~\cite[Theorems~30.1 and~30.2]{Billingsley1995}. If $\int g\,d\mu_\theta=0$ there is nothing to prove.
\end{proof}

The hypothesis $b_n>L_{\max}$ enters only through the global comparison of $\Delta_n(H)$ with $F$. For type~I potentials the window estimate of Section~\ref{sec:typeI} makes the comparison local in energy, and the hypothesis can be weakened to $b_n>L(E_0)$ near a given energy $E_0$.

\begin{lemma}[Localized arcsine law]\label{lem:arcsineloc}
Let $v$ be real analytic with $H_{v,\alpha}$ supercritical and type~I, and let $\theta\in G^v_\alpha$. Let $E_0\in\R$ and $\gamma>0$, and let $\mathcal N$ be an infinite set of indices with $b_n\ge\ell(E_0)+30\gamma$ for $n\in\mathcal N$. For every real $g\in C_c^\infty(\R)$ supported in $\{E:|\ell(E)-\ell(E_0)|<\gamma\}$ and every $\phi\in C_c^1(\R)$,
\[
\int g\,\phi(\Delta_n)\,d\mu_\theta-\sigma(\phi)\int g\,d\mu_\theta\to0\qquad(n\to\infty,\ n\in\mathcal N).
\]
If $g\ge0$, the same holds for every bounded continuous $\phi$.
\end{lemma}

\begin{proof}
Write $q=q_n$, $p=p_n$, $\delta=\delta_n$, $b=b_n$, $\ell_0=\ell(E_0)$, $J_g=\supp g$, $H=H_{v,\alpha,\theta}$, and let $Y=H_{v,p/q,\theta}$, the periodic operator whose discriminant is $\Delta_n$. Their potentials satisfy $|V_H(k)-V_Y(k)|\le\Lip(v)|k||\delta|/q$. Put $T=T_n=e^{(\ell_0+2\gamma)q}$, $z=E+i/T$, $G_H=(H-z)^{-1}\delta_0$ and $G_Y=(Y-z)^{-1}\delta_0$. Write $\mu_Y$ for the spectral measure of $\delta_0$ and $Y$.

\emph{Step 1: confinement on $J_g$.} Apply Theorem~\ref{thm:tI-window} with $c=\ell_0+30\gamma$, $d=\ell_0+\gamma$, $B=\ell_0+3\gamma$, $\eps=\gamma/2$, $m=1$, and with $\xi<\min\{\gamma/2,c,d,\ell_*,1,\eta_0\}/100$ and $\zeta<\xi/100$. Since $\theta\in G^v_\alpha$, Lemma~\ref{lem:tI-Wgood} gives (W) at $(q_n,q_n\theta,\tau_{q_n})$ for all large $n\in\mathcal N$. Let $R=e^{4\gamma q}$. If $E\in J_g$ and $\dist(E,\Sigma)\le e^{-\xi q}$, then $\ell(E)>\ell_0-\gamma$, so $e^{11\xi q}[1+Te^{-\ell(E)q}]\le e^{11\xi q}(1+e^{3\gamma q})\le R$ for large $q$, and part~2 of Theorem~\ref{thm:tI-window} gives $\|\un_{\{|k|>R\}}G_H\|\le4T^{-1}$. If $E\in J_g$ and $\dist(E,\Sigma)\ge e^{-\xi q}$, part~3 gives $\|\un_{\{|k|>R\}}G_H\|^2\le64e^{3\xi q}\exp(-e^{(4\gamma-\xi)q}/4)\le16T^{-2}$ for large $q$. So $\|\un_{\{|k|>R\}}G_H\|\le4T^{-1}$ for every $E\in J_g$ and all large $n\in\mathcal N$.

\emph{Step 2: comparison with the periodic operator.} The second resolvent identity gives $G_H-G_Y=-(Y-z)^{-1}(V_H-V_Y)G_H$, and $\|(Y-z)^{-1}\|\le T$. Splitting the multiplication by $V_H-V_Y$ into $|k|\le R$ and $|k|>R$,
\[
\|G_H-G_Y\|\le T\,\Lip(v)R|\delta|\,\|G_H\|+8\|v\|_\infty\qquad(E\in J_g).
\]
Put $D^2=(\pi T)^{-1}\int_{J_g}\|G_H-G_Y\|^2\,dE$. Since $(\pi T)^{-1}\int_\R\|G_H\|^2\,dE=1$ by Lemma~\ref{lem:parseval}, Minkowski's inequality gives
\[
D\le T\,\Lip(v)R|\delta|+8\|v\|_\infty\bigl(|J_g|/(\pi T)\bigr)^{1/2}\le Ce^{(\ell_0+6\gamma-b)q}+CT^{-1/2}\to0,
\]
because $b\ge\ell_0+30\gamma$. For every bounded Borel $f$ supported in $J_g$, writing $\|G_H\|^2-\|G_Y\|^2=\langle G_H-G_Y,G_H\rangle+\langle G_Y,G_H-G_Y\rangle$ and using the Cauchy--Schwarz inequality and the normalization for $H$ and $Y$,
\[
\Bigl|\frac1{\pi T}\int f(E)\bigl(\|G_H\|^2-\|G_Y\|^2\bigr)\,dE\Bigr|\le2\|f\|_\infty D.
\]

\emph{Step 3: smoothing.} By Lemma~\ref{lem:parseval}, $(\pi T)^{-1}\int f\|G_H\|^2dE=\int\mathcal P_{1/T}*f\,d\mu_\theta$, and the same holds for $Y$ and $\mu_Y$. For a Lipschitz function $f$ and $\epsilon\in(0,1]$,
\[
|\mathcal P_\epsilon*f-f|\le\Lip(f)\int_{|u|\le1}|u|\mathcal P_\epsilon(u)\,du+2\|f\|_\infty\int_{|u|>1}\mathcal P_\epsilon\le\Lip(f)\,\epsilon\log(1+\epsilon^{-2})+2\|f\|_\infty\epsilon.
\]
Apply this to $f_n=g\,\phi(\Delta_n)$, whose derivative is supported in $J_g$. Since $\ell<\ell_0+\gamma$ on $J_g$, Lemmas~\ref{lem:growth} and~\ref{lem:ratprod}(2) with $K=J_g$ and $h=\ell_0+3\gamma/2$ give $\sup_{J_g}|\Delta_n'|\le Cqe^{(\ell_0+3\gamma/2)q}$, so $\Lip(f_n)\le Cqe^{(\ell_0+3\gamma/2)q}$ and $\Lip(f_n)T^{-1}\log T\to0$. Together with Step~2,
\[
\int f_n\,d\mu_\theta-\int f_n\,d\mu_Y\to0,\qquad\int g\,d\mu_\theta-\int g\,d\mu_Y\to0.
\]

\emph{Step 4: the periodic operator.} By Lemma~\ref{lem:magic}, $\Delta_n(Y)=S^q+S^{-q}$, and $\sigma(Y)=\{|\Delta_n|\le2\}$. For a polynomial $P$,
\[
\int g\,P(\Delta_n)\,d\mu_Y=\langle g(Y)\delta_0,P(S^q+S^{-q})\delta_0\rangle=\sigma(P)\int g\,d\mu_Y+o(1),
\]
because $(S^q+S^{-q})^k\delta_0$ is a combination of the vectors $\delta_{iq}$, and $|\langle\delta_{iq},g(Y)\delta_0\rangle|\le C_g|iq|^{-1}$ for $i\ne0$ by Lemma~\ref{lem:QL}, uniformly in the periodic operator. The signed measure $(\Delta_n)_*(g\mu_Y)$ is supported in $[-2,2]$ and has total variation at most $\|g\|_\infty$. Approximating $\phi$ uniformly on $[-2,2]$ by polynomials gives $\int g\,\phi(\Delta_n)\,d\mu_Y-\sigma(\phi)\int g\,d\mu_Y\to0$. Combined with Step~3, this proves the first claim.

If $g\ge0$, the measures $(\Delta_n)_*(g\mu_\theta)$ have the fixed total mass $\int g\,d\mu_\theta$, and they converge vaguely to $\sigma\int g\,d\mu_\theta$ because $C_c^1$ is dense in $C_0$. With conservation of mass this is weak convergence.
\end{proof}

\subsection{The elliptic run}\label{sec:run}

In this subsection $v$ is Lipschitz, $\theta\in\T$, $p/q$ is a convergent of $\alpha$ with $\delta=q\alpha-p$, and $z=E+i\eps$ with $E$ real and $\eps>0$. Let $u\in\ell^2(\Z_{\ge0})$ satisfy
\[
u(k+1)+u(k-1)+v(\theta+k\alpha)u(k)=zu(k)\qquad(k\ge1).
\]
For $m\ge0$ put $x_m=\theta+(mq+1)\alpha$, $A_m(w)=A_q(w,x_m)$ and $\Psi_m=(u(mq+1),u(mq))^{\mathsf T}$, so that $\Psi_{m+1}=A_m(z)\Psi_m$ and $u(mq+1+r)=e_1^{\mathsf T}A_r(z,x_m)\Psi_m$ for $0\le r<q$. The block masses are $M_m=\sum_{r<q}|u(mq+1+r)|^2$. When $A_m(E)$ is elliptic, a \emph{frame} is a matrix $B_m\in SL(2,\R)$ with $B_m(i)$ the fixed point of $A_m(E)$ in the upper half-plane, so that $B_m^{-1}A_m(E)B_m=R_{\varrho_m}$ is a rotation (Lemma~\ref{lem:blocks}(4)). Given a frame we put
\[
u^{(m)}_r=e_1^{\mathsf T}A_r(E,x_m)B_m,\qquad S_m=\sum_{r<q}\|u^{(m)}_r\|^2,\qquad G_m=\sum_{r<q}(u^{(m)}_r)^{\mathsf T}u^{(m)}_r,\qquad\Phi_m=B_m^{-1}\Psi_m .
\]
Then $G_m$ is real, symmetric and positive semidefinite with trace $S_m$, and $S_m=S_{A_m(E)}$ in the notation of Section~\ref{sec:lower}.

\begin{lemma}[Elliptic run]\label{lem:run}
Let $0<\eta\le1$ and suppose $|\tr A_0(E)|\le2-\eta$. Let $B_0$ be a frame for $A_0(E)$, put $\kappa=(\Lip(v)+1)\max(\eps,|\delta|)$ and $K=\lfloor\eta^{3/2}/(2^{10}\kappa S_0)\rfloor$, and assume $K\ge8$. Let $m_0=\lceil K/2\rceil$. Then
\[
\sum_{k\ge m_0q+1}|u(k)|^2\ge\frac{\eta^{5/2}}{2^{22}}\,\frac{\|\Phi_0\|^2}{\kappa},\qquad\sum_{k\ge1}|u(k)|^2\le\Bigl(\frac1\kappa+\frac2\eps\Bigr)\|\Phi_0\|^2,\qquad|u(0)|^2\le\|B_0\|^2\|\Phi_0\|^2.
\]
\end{lemma}

The run has about $Kq$ sites, of order $q/(\max(\eps,|\delta|)S_0)$ up to factors depending on $\eta$ and $\Lip(v)$. The lemma controls the size of the conjugated vectors $\Phi_m$ and the masses of pairs of adjacent blocks; it says nothing about individual sites, where a standing wave may have nodes.

\begin{proof}
\emph{Margin.} If $M\in SL(2,\R)$ and $|\tr M|\le2-\eta/2$, then $M$ is elliptic, conjugate to a rotation $R_\varrho$ with $|\cos\varrho|\le1-\eta/4$, and $|\sin\varrho|\ge s_0:=\sqrt\eta/2$, since $\sin^2\varrho\ge\eta/2-\eta^2/16\ge\eta/4$.

\emph{Step 1: from one block to the next.} Since $x_{m+1}=x_m+\delta\bmod 1$, the potentials of block $m+1$ and block $m$ at the same position $r$ are $v(x_m+\delta+r\alpha)$ and $v(x_m+r\alpha)$; they differ by at most $\Lip(v)|\delta|$. Suppose $A_m(E)$ has a frame $B_m$. Lemma~\ref{lem:blocks}(5) gives, for $0\le r\le q$,
\[
A_r(E,x_{m+1})B_m=A_r(E,x_m)B_m\Pi_r,\qquad\|\Pi_r-I\|\le y_m:=e^{\Lip(v)|\delta|S_m}-1,
\]
and in the same way $A_r(z,x_m)B_m=A_r(E,x_m)B_m\Pi_r(i\eps)$ with $\|\Pi_r(i\eps)-I\|\le e^{\eps S_m}-1$. In particular $B_m^{-1}A_{m+1}(E)B_m=R_{\varrho_m}\Pi_q$ and $|\tr A_{m+1}(E)-\tr A_m(E)|\le2y_m$.

\emph{Step 2: frame update.} Let $M'=R_\varrho(I+Y)\in SL(2,\R)$ with $|\sin\varrho|\ge s_0$ and $\|Y\|\le y\le s_0/8$. Conjugation by the unitary Cayley matrix $2^{-1/2}\bigl(\begin{smallmatrix}1&-i\\1&i\end{smallmatrix}\bigr)$ maps $SL(2,\R)$ acting on the upper half-plane to matrices $\bigl(\begin{smallmatrix}a&b\\\bar b&\bar a\end{smallmatrix}\bigr)$ with $|a|^2-|b|^2=1$ acting on the disc, maps $i$ to $0$, maps $R_\varrho$ to $\operatorname{diag}(e^{-i\varrho},e^{i\varrho})$ and preserves operator norms. So $|b|\le y$ and $|a-e^{-i\varrho}|\le y$, hence $|\operatorname{Im}a|\ge s_0-y\ge s_0/2>|b|$. Then $M'$ is elliptic, and its fixed point in the disc has modulus $|b|/(|\operatorname{Im}a|+((\operatorname{Im}a)^2-|b|^2)^{1/2})\le2y/s_0\le1/4$. Since $\log\frac{1+r}{1-r}\le3r$ for $0\le r\le1/2$, the fixed point $w'$ of $M'$ in the upper half-plane satisfies $d(i,w')\le6y/s_0$. Apply this to $M'=B_m^{-1}A_{m+1}(E)B_m=R_{\varrho_m}\Pi_q$, with $y=y_m$, whenever $|\sin\varrho_m|\ge s_0$ and $y_m\le s_0/8$. Let $N_m\in SL(2,\R)$ be the positive symmetric matrix with $N_m(i)=w'$; then $\|N_m\|^2=\|N_m^{-1}\|^2=e^{d(i,w')}\le e^{d_m}$ with $d_m:=12y_m/s_0$. The matrix $B_{m+1}:=B_mN_m$ is a frame for $A_{m+1}(E)$, because $B_{m+1}(i)=B_m(w')$ is the fixed point of $A_{m+1}(E)=B_mM'B_m^{-1}$.

\emph{Step 3: bootstrap.} We show by induction on $m\le K$ that frames are defined and
\begin{equation}\label{eq:run-a}
|\tr A_m(E)|\le2-\eta/2,\qquad S_m\in[S_0/2,2S_0],\qquad\|\Phi_m\|\in[\|\Phi_0\|/2,2\|\Phi_0\|].
\end{equation}
This holds at $m=0$. Assume it for all $j<m$. Since $K\ge8$, $\kappa S_0\le\eta^{3/2}/2^{13}$ and $K\kappa S_0\le\eta^{3/2}/2^{10}$. Using $\Lip(v)|\delta|\le\kappa$ and $e^x-1\le2x$ for $0\le x\le1$, for $j<m$,
\[
y_j\le2\Lip(v)|\delta|S_j\le4\kappa S_0\le\eta^{3/2}/2^{11}\le s_0/8 .
\]
By Step~2 the frames $B_{j+1}$ are defined. By Step~1, $|\tr A_m(E)-\tr A_0(E)|\le2\sum_{j<m}y_j\le8K\kappa S_0\le\eta/2$, which gives the first claim in~\eqref{eq:run-a}. Next, with $\Pi_r$ the factors of Step~1 for the block $j$, $u^{(j+1)}_r=e_1^{\mathsf T}A_r(E,x_{j+1})B_jN_j=u^{(j)}_r\Pi_rN_j$, so
\[
S_j(1-y_j)^2e^{-d_j}\le S_{j+1}\le S_j(1+y_j)^2e^{d_j},
\]
and $|\log(S_m/S_0)|\le\sum_{j<m}(3y_j+d_j)\le\sum_{j<m}27y_j/\sqrt\eta\le108K\kappa S_0/\sqrt\eta\le108\eta/2^{10}<0.11$. Finally $\Phi_{j+1}=N_j^{-1}\bigl[B_j^{-1}A_j(z)B_j\bigr]\Phi_j$, and $B_j^{-1}A_j(z)B_j=R_{\varrho_j}\Pi_q(i\eps)$ and its inverse have norm at most $e^{\eps S_j}$, since a $2\times2$ matrix of determinant one and its inverse have the same norm. Since $\eps\le\kappa$,
\[
\bigl|\log(\|\Phi_m\|/\|\Phi_0\|)\bigr|\le\sum_{j<m}\bigl(\eps S_j+d_j/2\bigr)\le2K\kappa S_0+48K\kappa S_0/\sqrt\eta<0.05 .
\]
This proves~\eqref{eq:run-a} for $m$.

\emph{Step 4: pairs of blocks.} Let $G$ be a real symmetric positive semidefinite $2\times2$ matrix with trace $S$, and $R=R_\varrho$. Write $G=(S/2)I+G_0$ with $G_0$ traceless and $\|G_0\|=g\le S/2$. Conjugation by $R$ rotates $G_0$ by the angle $2\varrho$ in the plane of traceless symmetric matrices, so $\|G_0+R^{\mathsf T}G_0R\|=2g|\cos\varrho|$ and
\[
G+R^{\mathsf T}GR\ge S(1-|\cos\varrho|)\,I .
\]
This is an inequality between Hermitian matrices, so $\Phi^*(G+R^{\mathsf T}GR)\Phi\ge S(1-|\cos\varrho|)\|\Phi\|^2$ for every $\Phi\in\C^2$. Now fix $m<K$. By Step~1,
\[
M_m=\sum_r\bigl|u^{(m)}_r\Pi_r(i\eps)\Phi_m\bigr|^2,\qquad M_{m+1}=\sum_r\bigl|u^{(m)}_r\Pi_r\Pi'_r(i\eps)R_{\varrho_m}\Pi_q(i\eps)\Phi_m\bigr|^2,
\]
where $\Pi'_r(i\eps)$ is the factor of Lemma~\ref{lem:blocks}(5) for the complex shift of block $m+1$, computed with the rows $u^{(m)}_r\Pi_r$ in the frame $B_m$, whose squared norms sum to at most $(1+y_m)^2S_m\le2S_m$. Each of the factors $\Pi_r$, $\Pi'_r(i\eps)$, $\Pi_r(i\eps)$, $\Pi_q(i\eps)$ is within $\eps_2:=e^{2(\eps+\Lip(v)|\delta|)S_m}-1\le e^{4\kappa S_0}-1\le8\kappa S_0\le\eta^{3/2}/2^{10}$ of the identity. So the vectors acted on by $u^{(m)}_r$ differ from $\Phi_m$, respectively $R_{\varrho_m}\Phi_m$, by at most $\eps_2\|\Phi_m\|$, respectively $((1+\eps_2)^3-1)\|\Phi_m\|\le4\eps_2\|\Phi_m\|$. Since $\sum_r|u^{(m)}_rw|^2\le S_m\|w\|^2$, the triangle inequality in $\ell^2(r)$ and $(x-y)^2\ge x^2/2-y^2$ give
\[
M_m+M_{m+1}\ge\tfrac12\Phi_m^*\bigl(G_m+R_{\varrho_m}^{\mathsf T}G_mR_{\varrho_m}\bigr)\Phi_m-17\eps_2^2S_m\|\Phi_m\|^2\ge\Bigl(\frac\eta8-17\eps_2^2\Bigr)S_m\|\Phi_m\|^2\ge\frac\eta{16}S_m\|\Phi_m\|^2 .
\]
By~\eqref{eq:run-a}, $M_m+M_{m+1}\ge(\eta/2^7)S_0\|\Phi_0\|^2$. Also $M_m\le(1+\eps_2)^2S_m\|\Phi_m\|^2\le9S_0\|\Phi_0\|^2$ for $m<K$.

\emph{Step 5: the tail beyond the run.} Multiply the equation by $\overline{u(k)}$, sum over $k>n\ge0$ and take imaginary parts. Since $u\in\ell^2$ and the potential is real, the sum telescopes:
\[
\eps\sum_{k>n}|u(k)|^2=\operatorname{Im}\bigl(\overline{u(n+1)}\,u(n)\bigr).
\]
For $\Psi=(u(n+1),u(n))^{\mathsf T}=B\Phi$ with $B$ real of determinant one, $\overline{\Psi_1}\Psi_2-\Psi_1\overline{\Psi_2}=\det(\bar\Psi,\Psi)=\det(\bar\Phi,\Phi)$, so $|\operatorname{Im}(\overline{u(n+1)}u(n))|\le|\Phi_1||\Phi_2|\le\|\Phi\|^2/2$. At $n=Kq$, with $\Psi=\Psi_K=B_K\Phi_K$ and~\eqref{eq:run-a},
\[
\sum_{k\ge Kq+1}|u(k)|^2\le\frac{\|\Phi_K\|^2}{2\eps}\le\frac{2\|\Phi_0\|^2}{\eps}.
\]

\emph{Step 6: conclusion.} Let $X=\eta^{3/2}/(2^{10}\kappa S_0)$, so $K=\lfloor X\rfloor\ge8$ and $K\ge X/2$. The blocks $m_0\le m\le K$ contain at least $K/8$ disjoint adjacent pairs $(m,m+1)$ with $m<K$, and they consist of sites $k\ge m_0q+1$. By Step~4,
\[
\sum_{k\ge m_0q+1}|u(k)|^2\ge\frac K8\cdot\frac{\eta}{2^7}S_0\|\Phi_0\|^2\ge\frac{X\eta S_0}{2^{11}}\|\Phi_0\|^2=\frac{\eta^{5/2}}{2^{21}}\frac{\|\Phi_0\|^2}{\kappa}.
\]
The sites $1\le k\le Kq$ form the blocks $m<K$, whose total mass is at most $9KS_0\|\Phi_0\|^2\le9XS_0\|\Phi_0\|^2\le\|\Phi_0\|^2/\kappa$. With Step~5 this gives the second bound. Finally $u(0)=e_2^{\mathsf T}\Psi_0=e_2^{\mathsf T}B_0\Phi_0$.
\end{proof}

No turning-point analysis is needed: Step~5 controls everything beyond the run through the flux, at the price of a damping $\eps\ge|\delta|/C$ when the run is to carry a fixed fraction. No arithmetic condition on the phase is used.

\subsection{A fixed fraction at resonant times}\label{sec:fraction}

\begin{proposition}\label{prop:fraction}
Let $v$ be Lipschitz and $\theta\in\T$. Let $a,\gamma>0$, $0<\eta\le1/4$, and let $K\subset\{E\in\R:L(E)<a\}$ be a compact interval. Let $\mathcal N$ be an infinite set of indices with $b_n\ge a+\gamma$ for $n\in\mathcal N$, and put
\[
T_n=|\delta_n|^{-1},\qquad N_n=\bigl\lfloor e^{(b_n-a-\gamma)q_n}\bigr\rfloor,\qquad I_n=\{E\in K:|\Delta_n(E)|\le2-2\eta\},
\]
and $f_0=\eta^{5/2}/(2^{25}(\Lip(v)+1))$.
There is $C_\eta$ such that for all large $n\in\mathcal N$ and every $E\in I_n$, with $G(k)=\langle\delta_k,(H_{v,\alpha,\theta}-E-i/T_n)^{-1}\delta_0\rangle$:
\begin{enumerate}
\item $\sum_{|k|>N_n}|G(k)|^2\ge2f_0\sum_{k\in\Z}|G(k)|^2$;
\item $\mu_\theta\bigl([E-|\delta_n|,E+|\delta_n|]\bigr)\le C_\eta e^{(a+\gamma/2)q_n}|\delta_n|$.
\end{enumerate}
If moreover $K'\subset K$ is a compact interval in the interior of $K$, then for all large $n\in\mathcal N$,
\[
\frac1{\pi T_n}\int_{I_n}\bigl\|\un_{\{|k|>N_n\}}(H_{v,\alpha,\theta}-E-i/T_n)^{-1}\delta_0\bigr\|^2\,dE\ \ge\ f_0\,\mu_\theta\bigl(\{E\in K':|\Delta_n(E)|\le2-3\eta\}\bigr).
\]
In particular $P(N_n,T_n)$ is at least the right-hand side.
\end{proposition}

\begin{proof}
Since $L$ is upper semicontinuous, $\max_KL<a$. Put $h_0=a+\gamma/2$, and let $U\supset K$ and $C$ be given by Lemma~\ref{lem:growth} with $h=h_0$, and $C'$ by Lemma~\ref{lem:ratprod}. Write $q=q_n$, $\delta=\delta_n$, $b=b_n$, $\eps=|\delta|=1/T_n$, $\kappa=(\Lip(v)+1)|\delta|$, and fix $E\in I_n$.

\emph{The right half-line.} The sequence $u(k)=G(k)$, $k\ge0$, is square summable and solves the equation of Section~\ref{sec:run} with $z=E+i\eps$ for $k\ge1$. Its first block is $A_0(E)=A_q(E,\theta+\alpha)$. Since $\Delta_\varphi$ is $1/q$-periodic and $\alpha\equiv\delta/q\bmod 1/q$, Lemma~\ref{lem:ratprod}(2),(3) give
\[
|\tr A_q(E,\theta+\alpha)-\Delta_n(E)|\le|\tr A_q(E,\theta+\alpha)-\Delta_{\theta+\alpha}(E)|+|\Delta_{\theta+\delta/q}(E)-\Delta_\theta(E)|\le4C'^2\Lip(v)qe^{h_0q}|\delta|,
\]
which is at most $\eta$ for large $n$ because $b>h_0$. So $|\tr A_0(E)|\le2-\eta$. By Lemma~\ref{lem:blocks}(1), the product rule and $|\sin\varrho_0|\ge\sqrt\eta/2$, $S_0=|\partial_E\tr A_0(E)|/|\sin\varrho_0|\le(4/\sqrt\eta)C^2qe^{h_0q}$. Hence $X=\eta^{3/2}/(2^{10}\kappa S_0)\ge c_\eta e^{(b-h_0)q}/q$ with $c_\eta>0$, so the run length $\lfloor X\rfloor$ of Lemma~\ref{lem:run} is at least $8$, and its $m_0$ satisfies $m_0q\ge(X-1)q/2\ge e^{(b-a-\gamma)q}\ge N_n$ for large $n$, uniformly in $E\in I_n$. Lemma~\ref{lem:run} gives, with $\Phi^+_0$ and $B^+_0$ its vector and frame,
\[
\mathrm{far}_+:=\sum_{k>N_n}|G(k)|^2\ge\frac{\eta^{5/2}}{2^{22}}\frac{\|\Phi^+_0\|^2}{\kappa},\qquad\sum_{k\ge1}|G(k)|^2\le\Bigl(\frac1\kappa+\frac2{|\delta|}\Bigr)\|\Phi^+_0\|^2\le\frac{3(\Lip(v)+1)}{\kappa}\|\Phi^+_0\|^2 .
\]

\emph{The left half-line.} The sequence $\tilde u(k)=G(-k)$, $k\ge0$, solves the same equation for $k\ge1$ with $v$ replaced by $\tilde v(x)=v(-x)$ and $\theta$ by $-\theta$. The function $\tilde v$ has the same Lipschitz constant. Conjugation by $\operatorname{diag}(1,-1)$ maps each one-step matrix to its transpose, so a product of one-step matrices and the product in reversed order have the same norm and the same trace. Hence $\|A^{\tilde v}_r(E,x)\|=\|A_r(E,-x-(r-1)\alpha)\|$, the same holds for rational products, and Lemmas~\ref{lem:growth} and~\ref{lem:ratprod} hold for $\tilde v$ with the same $U$, $C$ and $C'$. The rational discriminant of $\tilde v$ at the phase $-\theta$ is a trace of the factors of $\Delta_\theta$ in reversed cyclic order, so it equals $\Delta_n$. The argument of the right half-line therefore applies verbatim and gives the same bounds for $\mathrm{far}_-:=\sum_{k<-N_n}|G(k)|^2$ and $\sum_{k\le-1}|G(k)|^2$, with a vector $\Phi^-_0$.

\emph{Part~1.} By Lemma~\ref{lem:blocks}(4), $\|A_0\|_{\mathrm{HS}}^2=2+4\sinh^2(d)\sin^2\varrho_0$ with $d=d(i,B^+_0(i))$, so $\|B^+_0\|^2=e^d\le1+2\sinh d\le1+2\sqrt2\|A_0\|/\sqrt\eta\le C_\eta e^{h_0q}$. So
\[
|G(0)|^2\le\|B^+_0\|^2\|\Phi^+_0\|^2\le C_\eta e^{h_0q}\kappa\,\frac{2^{22}}{\eta^{5/2}}\,\mathrm{far}_+=o(1)\,\mathrm{far}_+,
\]
since $\kappa e^{h_0q}\le Ce^{(h_0-b)q}$. Adding the two half-lines and the origin,
\[
\sum_k|G(k)|^2\le\Bigl(\frac{3(\Lip(v)+1)2^{22}}{\eta^{5/2}}+o(1)\Bigr)(\mathrm{far}_++\mathrm{far}_-)\le\frac{1}{2f_0}\sum_{|k|>N_n}|G(k)|^2 .
\]

\emph{Part~2.} Since $\operatorname{Im}G(0)=\eps\sum_k|G(k)|^2\ge\eps\,\mathrm{far}_+\ge\eta^{5/2}\|\Phi^+_0\|^2/(2^{22}(\Lip(v)+1))$ and $\operatorname{Im}G(0)\le|G(0)|\le\|B^+_0\|\|\Phi^+_0\|$, we get $\|\Phi^+_0\|\le C_\eta'\|B^+_0\|$ and $|G(0)|\le C'_\eta\|B^+_0\|^2\le C_\eta e^{h_0q}$. Finally $\operatorname{Im}G(0)=\int\eps\,((x-E)^2+\eps^2)^{-1}d\mu_\theta(x)\ge\mu_\theta([E-\eps,E+\eps])/(2\eps)$.

\emph{The integrated bound.} By Lemma~\ref{lem:parseval} and part~1,
\[
\frac1{\pi T_n}\int_{I_n}\sum_{|k|>N_n}|G(k)|^2\,dE\ge\frac{2f_0}{\pi T_n}\int_{I_n}\sum_k|G(k)|^2\,dE=2f_0\int\!\!\int_{I_n}\mathcal P_\eps(E-x)\,dE\,d\mu_\theta(x).
\]
Let $r_n=\min\{\dist(K',\R\setminus K),\ \eta/(2C'^2qe^{h_0q})\}$. If $x\in K'$ and $|\Delta_n(x)|\le2-3\eta$, then $[x-r_n,x+r_n]\subset I_n$ by Lemma~\ref{lem:ratprod}(2), and $\int_{|u|\le r_n}\mathcal P_\eps(u)\,du\ge1-2\eps/(\pi r_n)$. Since $\eps/r_n\le C_\eta qe^{(h_0-b)q}\to0$, the right side is at least $f_0\,\mu_\theta(\{x\in K':|\Delta_n(x)|\le2-3\eta\})$ for large $n$. Restricting the energy integral of Lemma~\ref{lem:parseval} to $I_n$ shows that $P(N_n,T_n)$ is at least the left side.
\end{proof}

\subsection{Strongly Liouville frequencies and Lipschitz potentials}\label{sec:U1}

When $\beta(\alpha)$ exceeds the largest Lyapunov exponent on the spectrum, the global arcsine law applies, and no type~I hypothesis, phase condition or confinement estimate is needed.

\begin{theorem}\label{thm:U1}
Let $v$ be Lipschitz with $L$ continuous on $\Sigma$, let $L_{\max}=\max_\Sigma L$, and let $\alpha$ be irrational with $\beta=\beta(\alpha)>L_{\max}$, possibly $\beta=\infty$. For every phase $\theta$ put $\ell_\theta=\min\{L(E):E\in\supp\mu_\theta\}$. Then:
\begin{enumerate}
\item $\beta^+(p)\ge F_\beta(\ell_\theta)$ for every $p>0$;
\item $d^+_{\mu_\theta}(E)\ge F_\beta(L(E))$ for $\mu_\theta$-almost every $E$; consequently $\dim_P^+(\mu_\theta)\ge F_\beta(\ell_\theta)$ and $\dim_P^-(\mu_\theta)\ge F_\beta(L_{\max})$.
\end{enumerate}
For every $\theta$, $\ell_\theta=\ell_*:=\min_\Sigma L$ holds at the phase $\theta$ or at the phase $\theta+\alpha$.
\end{theorem}

\begin{proof}
Fix $\eta=1/4$, and fix $c\in(L_{\max},\beta)$.

1. We may assume $\ell_\theta<\beta$. Let $E_0\in\supp\mu_\theta$ with $L(E_0)=\ell_\theta$, and let $\gamma>0$ with $\ell_\theta+2\gamma<\beta$. Put $a=\ell_\theta+\gamma$. Since $L$ is upper semicontinuous, $\{L<a\}$ is an open set containing $E_0$. Choose compact intervals $K'\subset\operatorname{int}K\subset K\subset\{L<a\}$ with $E_0$ in the interior of $K'$, and $g\in C_c^\infty$ with $0\le g\le1$, $\supp g\subset K'$ and $g>0$ near $E_0$, so that $\int g\,d\mu_\theta>0$. Choose an infinite set $\mathcal N$ with $b_n\to\beta$ and $b_n\ge\max(c,a+\gamma)$ for $n\in\mathcal N$. Since $0\le g\le\un_{K'}$, Lemma~\ref{lem:arcsine} and the portmanteau theorem applied to the open set $O_\eta=(-2+3\eta,2-3\eta)$ give
\[
\liminf_{n\in\mathcal N}\mu_\theta\bigl(\{E\in K':|\Delta_n(E)|\le2-3\eta\}\bigr)\ge\liminf_{n\in\mathcal N}\int g\,\un_{O_\eta}(\Delta_n)\,d\mu_\theta\ge\sigma(O_\eta)\int g\,d\mu_\theta=:2c_1>0 .
\]
Proposition~\ref{prop:fraction} gives $P(N_n,T_n)\ge f_0c_1$ for large $n\in\mathcal N$, and hence $M_p(T_n)\ge N_n^pP(N_n,T_n)\ge f_0c_1N_n^p$. Since $\log N_n\ge(b_n-a-\gamma)q_n-\log2$ and $\log T_n=b_nq_n$, we get $\beta^+(p)\ge1-(\ell_\theta+2\gamma)/\beta$, or $\beta^+(p)\ge1$ if $\beta=\infty$. Let $\gamma\to0$.

2. Fix rationals $a>0$ and $\gamma>0$ with $a+\gamma<\beta$, and choose an infinite set $\mathcal N$ with $b_n\to\beta$ and $b_n\ge\max(c,a+\gamma)$. The open set $\{L<a\}$ is a countable union of the interiors of compact intervals $K'_i$ with $K'_i\subset\operatorname{int}K_i$ and $K_i\subset\{L<a\}$. Fix $i$, and let $g\ge0$ be smooth, positive exactly on $\operatorname{int}K'_i$. For $0<\eta\le1/4$ let $\Lambda_\eta$ be the set of $E\in K'_i$ with $|\Delta_n(E)|\le2-2\eta$ for infinitely many $n\in\mathcal N$. For the finite measure $g\mu_\theta$, Fatou's lemma for sets and the portmanteau theorem for closed sets give
\[
(g\mu_\theta)(K'_i\setminus\Lambda_\eta)\le\liminf_{n\in\mathcal N}(g\mu_\theta)\bigl(\{|\Delta_n|>2-2\eta\}\bigr)\le\sigma\bigl(\{|x|\ge2-2\eta\}\bigr)\int g\,d\mu_\theta,
\]
which tends to zero with $\eta$. So $g\mu_\theta$-almost every $E\in K'_i$ lies in some $\Lambda_{1/j}$. For such $E$, Proposition~\ref{prop:fraction}(2) holds at infinitely many $n\in\mathcal N$, so
\[
\frac{\log\mu_\theta([E-|\delta_n|,E+|\delta_n|])}{\log|\delta_n|}\ge1-\frac{a+\gamma/2}{b_n}-\frac{\log C_\eta}{b_nq_n}\to1-\frac{a+\gamma/2}{\beta}
\]
along these $n$, and $d^+_{\mu_\theta}(E)\ge1-(a+\gamma/2)/\beta$, read as $1$ if $\beta=\infty$. Since $g>0$ on $\operatorname{int}K'_i$, this holds for $\mu_\theta$-almost every $E\in\operatorname{int}K'_i$, hence for $\mu_\theta$-almost every $E$ with $L(E)<a$. Intersecting over the countably many rational $a,\gamma$ gives $d^+_{\mu_\theta}(E)\ge F_\beta(L(E))$ for $\mu_\theta$-almost every $E$ with $L(E)<\beta$; for $L(E)\ge\beta$ there is nothing to prove.

The essential infimum of $L$ with respect to $\mu_\theta$ is $\ell_\theta$, because $L$ is continuous on $\Sigma\supset\supp\mu_\theta$, and $L\le L_{\max}$ on $\Sigma$. Since $F_\beta$ is nonincreasing, this gives the bounds on $\dim_P^\pm$.

The last statement holds because $\delta_0$ and $\delta_1$ form a cyclic pair for $H$, so $\supp\mu_{\delta_0}\cup\supp\mu_{\delta_1}=\Sigma$, and the translation $(Uu)(n)=u(n+1)$ satisfies $UH_{v,\alpha,\theta}U^{-1}=H_{v,\alpha,\theta+\alpha}$ and $U\delta_1=\delta_0$, so the spectral measure of $\delta_1$ at the phase $\theta$ is $\mu_{\theta+\alpha}$.
\end{proof}

\subsection{Type I potentials: every moment, energy by energy}\label{sec:typeI-moments}

In this subsection $v$ is real analytic, $H_{v,\alpha}$ is supercritical and of type~I, $\ell=L$ is continuous on $\R$, and $\theta\in G^v_\alpha$. For such $\theta$ put
\[
\ell_\theta=\min\{\ell(E):E\in\supp\mu_\theta\},\qquad\ell^\theta=\max\{\ell(E):E\in\supp\mu_\theta\}.
\]
Then $\ell_*\le\ell_\theta\le\ell^\theta$, and by continuity these are the $\mu_\theta$-essential infimum and supremum of $\ell$.

\begin{theorem}[Every moment]\label{thm:U2}
Let $0<\beta=\beta(\alpha)<\infty$. For every $p>0$,
\[
F_\beta(\ell_\theta)\le\beta^+(p)\le F_\beta(\ell_*).
\]
Moreover, for every $\gamma>0$ there is $c>0$ such that
\[
P\bigl(e^{(b_n-\ell_\theta-\gamma)q_n},\,|\delta_n|^{-1}\bigr)\ge c\qquad\text{for all large }n\text{ with }b_n\ge\ell_\theta+15\gamma.
\]
\end{theorem}

\begin{proof}
The upper bound follows from $\beta^+(p)\le\au$ and Theorem~\ref{thm:tI-law}. For the rest let $\gamma>0$, put $\gamma'=\gamma/2$, and let $E_0\in\supp\mu_\theta$ with $\ell(E_0)=\ell_\theta$. Choose compact intervals $K'\subset\operatorname{int}K$, with $E_0$ in the interior of $K'$, such that $|\ell-\ell_\theta|<\gamma'$ on $K$. Choose $g\in C_c^\infty$ with $0\le g\le1$, $\supp g\subset K'$ and $g>0$ near $E_0$, so that $\int g\,d\mu_\theta>0$. Let $\mathcal N=\{n:b_n\ge\ell_\theta+30\gamma'\}$. If $\mathcal N$ is finite there is nothing to prove; otherwise Lemma~\ref{lem:arcsineloc} with $E_0$, $\gamma'$ and $g$, and the portmanteau theorem, give
\[
\liminf_{n\in\mathcal N}\mu_\theta\bigl(\{E\in K':|\Delta_n(E)|\le2-3\eta\}\bigr)\ge\sigma\bigl((-2+3\eta,2-3\eta)\bigr)\int g\,d\mu_\theta=:2c_1>0
\]
with $\eta=1/4$. Proposition~\ref{prop:fraction} with $a=\ell_\theta+\gamma'$ and $\gamma'$ in place of $\gamma$ applies, because $K\subset\{\ell<a\}$ and $b_n\ge a+\gamma'$ on $\mathcal N$; its radius is $N_n=\lfloor e^{(b_n-\ell_\theta-\gamma)q_n}\rfloor$. It gives $P(N_n,|\delta_n|^{-1})\ge f_0c_1$ for large $n\in\mathcal N$, which is the fixed-fraction statement. If $\ell_\theta<\beta$, choose $\gamma$ with $\ell_\theta+15\gamma<\beta$ and $n\in\mathcal N$ with $b_n\to\beta$; as in the proof of Theorem~\ref{thm:U1}, $\beta^+(p)\ge1-(\ell_\theta+\gamma)/\beta$. Let $\gamma\to0$.
\end{proof}

Corollary~\ref{cor:allmoments} below shows that the two bounds of Theorem~\ref{thm:U2} coincide.

\begin{remark}\label{rem:U2-shift}
For $\beta(\alpha)=0$ all moment exponents vanish by Theorem~\ref{thm:tI-law}, and for $\beta(\alpha)=\infty$ they all equal one, for every phase~\cite[Theorem~3(2)]{JZ2022}.
\end{remark}

To resolve the dynamics in energy we filter the initial state. For real $g\in C_c^\infty(\R)$ let $\psi_g=g(H_{v,\alpha,\theta})\delta_0$. By Lemma~\ref{lem:QL}, $\psi_g$ has finite moments, and its spectral measure is $\mu_{\psi_g}=g^2\mu_\theta$.

\begin{lemma}[Filtering]\label{lem:filter}
Let $H$ be a Schr\"odinger operator with hopping one and bounded real potential, let $g\in C_c^\infty(\R)$ be real, and let $J,M\ge1$ be integers. For $T\ge1$ and real $E$ put $z=E+i/T$, $G_z=(H-z)^{-1}\delta_0$ and
\[
c_J(E,T)=\sum_{j<J}\frac{g^{(j)}(E)}{j!}\Bigl(\frac iT\Bigr)^j .
\]
There is $C$, depending on $g$, $J$, $M$ and a bound for $\|H\|$, such that for every $T\ge1$ and every integer $N\ge2J$,
\[
\Bigl|P_{\psi_g}(N,T)^{1/2}-\Bigl(\frac1{\pi T}\int_\R|c_J(E,T)|^2\,\bigl\|\un_{\{|n|>N\}}G_z\bigr\|^2\,dE\Bigr)^{1/2}\Bigr|\le CN^{-M}+CT^{-J},
\]
where $\psi_g=g(H)\delta_0$. Moreover $c_J(\cdot,T)$ vanishes outside $\supp g$ and $|c_J(E,T)-g(E)|\le C/T$.
\end{lemma}

\begin{proof}
Taylor's formula at $E$ gives $g(x)=\sum_{j<J}g^{(j)}(E)(x-E)^j/j!+(x-E)^Jr_E(x)$ with $r_E(x)=\int_0^1(1-s)^{J-1}g^{(J)}(E+s(x-E))\,ds/(J-1)!$, so $|r_E(x)|\le\|g^{(J)}\|_\infty/J!$ for all $E,x$. Put $a=i/T$. The identity $(H-E)G_z=\delta_0+aG_z$ and induction on $j$ give
\[
(H-E)^jG_z=\sum_{r=0}^{j-1}a^r(H-E)^{j-1-r}\delta_0+a^jG_z .
\]
Applying the functional calculus,
\[
g(H)G_z=c_J(E,T)\,G_z+w_E+a^Jr_E(H)G_z,
\]
where $w_E$ is a sum of the vectors $(H-E)^k\delta_0$, $k<J$, with coefficients bounded by $C$, and of the vectors $a^rr_E(H)(H-E)^{J-1-r}\delta_0$, $0\le r<J$. The former are supported in $[-J,J]$ and vanish under $\un_{\{|n|>N\}}$. For the latter, let $\sigma(H)\subset[-R_0,R_0]$ and fix $\chi\in C_c^\infty$ equal to $1$ on $[-R_0-1,R_0+1]$; then $r_E(H)=(\chi r_E)(H)$. The functions $\chi r_E$ are bounded in $C^{M+3}$ for $E$ in a compact set. If $E$ lies outside $\supp g$ and $|E|\ge2R_0+4$, all derivatives of $g$ vanish at $E$, so $r_E(x)=g(x)/(x-E)^J$ on $\supp\chi$ and $\|\chi r_E\|_{C^{M+3}}\le C(1+|E|)^{-J}$. Lemma~\ref{lem:QL} with $M+1$ in place of $M$ and the Schur test give $\|\un_{\{|n|>N\}}(\chi r_E)(H)\un_{[-J,J]}\|\le C(1+|E|)^{-J}N^{-M}$ for $N\ge2J$, and $\|(H-E)^{J-1-r}\delta_0\|\le(R_0+|E|)^{J-1}$. Hence $\|\un_{\{|n|>N\}}w_E\|\le C(1+|E|)^{-1}N^{-M}$. The last term has norm at most $CT^{-J}\|G_z\|$.

By Lemma~\ref{lem:parseval} with $\psi=\psi_g$, $P_{\psi_g}(N,T)=(\pi T)^{-1}\int\|\un_{\{|n|>N\}}g(H)G_z\|^2\,dE$, since $(H-z)^{-1}\psi_g=g(H)G_z$. Minkowski's inequality in $L^2(\R,dE/(\pi T);\ell^2(\Z))$, $(\pi T)^{-1}\int\|G_z\|^2dE=1$ and $\int(1+|E|)^{-2}dE<\infty$ give the claim. The last two statements are immediate from the definition of $c_J$.
\end{proof}

\begin{theorem}[Energy by energy]\label{thm:U3}
Let $\beta=\beta(\alpha)<\infty$, $E_0\in\Sigma$ and $0<\gamma<\ell_*/2$. Let $g\in C_c^\infty(\R)$ be real, supported in $\{E:|\ell(E)-\ell(E_0)|<\gamma\}$, with $\psi_g\ne0$. Then for every $p>0$,
\[
F_\beta\bigl(\ell(E_0)+\gamma\bigr)\le\beta^+_{\psi_g}(p)\le F_\beta\bigl(\ell(E_0)-\gamma\bigr).
\]
\end{theorem}

Letting the support of $g$ shrink to an energy $E$, a packet filtered near $E$ spreads, in every moment, with exponent $F_\beta(\ell(E))$. The statement concerns a family of filters of shrinking width; a fixed filter of positive width need not have a single exponent of this form.

\begin{proof}
We normalize $\psi=\psi_g/\|\psi_g\|$; this changes no exponent.

\emph{Upper bound.} Put $\ell=\ell(E_0)-\gamma$, which is positive. We verify (M1) and (M2$_\ell$) for $H=H_{v,\alpha,\theta}$ and $\psi$, and apply Theorem~\ref{thm:covering}. Since $e^{-itH}\psi_g=g(H)e^{-itH}\delta_0$, splitting $e^{-itH}\delta_0$ at $|n|=N/2$ gives
\[
P_{\psi_g}(N,T)\le2\bigl\|\un_{\{|n|>N\}}g(H)\un_{\{|n|\le N/2\}}\bigr\|^2+2\|g\|_\infty^2P(N/2,T)\le C_MN^{-M}+2\|g\|_\infty^2P(N/2,T)
\]
for every $M$, by Lemma~\ref{lem:QL} and the Schur test. Theorem~\ref{thm:tI-M1} gives (M1). For (M2$_\ell$) fix $c,d,B,\eps,m$. Apply Theorem~\ref{thm:tI-window} with $\eps/2$ in place of $\eps$, with $m+1$ in place of $m$, and with $\xi<\min\{\eps/2,c,d,\ell_*,1,\eta_0\}/100$ and $\zeta<\xi/100$; Lemma~\ref{lem:tI-Wgood} provides (W) at $(q_n,q_n\theta,\tau_{q_n})$ for large $n$ with $b_n\ge c$. Let
\[
R^g(T)=e^{11\xi q_n}\Bigl[1+\min\bigl(Te^{-\ell q_n},\,e^{(b_n-\ell)_+q_n}\bigr)\Bigr]\le R_n(T).
\]
For $E\in\supp g$ we have $\ell(E)>\ell$. If moreover $\dist(E,\Sigma)\le e^{-\xi q_n}$, part~2 of Theorem~\ref{thm:tI-window} gives $\|\un_{\{|n|>R^g(T)\}}G_z\|\le4T^{-m-1}$, since the radius there is nonincreasing in $\ell(E)$. If $\dist(E,\Sigma)\ge e^{-\xi q_n}$, part~3 gives the same bound for large $n$, since $R^g(T)\ge e^{11\xi q_n}$. Lemma~\ref{lem:filter} with an integer $J\ge\max\{1,\lceil m\rceil\}$ and an integer $M\ge mB/(11\xi)$, together with $|c_J|\le C$ and $c_J=0$ off $\supp g$, gives
\[
P_{\psi_g}\bigl(R^g(T),T\bigr)\le2\,\frac{C|\supp g|}{\pi T}\,16T^{-2m-2}+2\bigl(CR^g(T)^{-M}+CT^{-J}\bigr)^2\le CT^{-m}
\]
for $e^{dq_n}\le T\le e^{Bq_n}$, because $R^g(T)\ge T^{11\xi/B}$. This is (M2$_\ell$). Theorem~\ref{thm:covering} gives $\alpha^+_{u,\psi}\le F_\beta(\ell)$, also when $\beta=0$, and $\beta^+_\psi(p)\le\alpha^+_{u,\psi}$.

\emph{Lower bound.} We may assume $\ell(E_0)+\gamma<\beta$. Since $\int g^2d\mu_\theta=\|\psi_g\|^2>0$, there is $E_1\in\supp\mu_\theta$ with $g(E_1)\ne0$, and $\ell(E_1)<\ell(E_0)+\gamma$. Choose $\gamma_1>0$ with $\ell(E_1)+2\gamma_1\le\ell(E_0)+\gamma$ and $\ell(E_1)+30\gamma_1<\beta$. Choose compact intervals $K'\subset\operatorname{int}K$, with $E_1$ in the interior of $K'$, such that $|\ell-\ell(E_1)|<\gamma_1$ and $|g|\ge g_1:=|g(E_1)|/2$ on $K$. Choose $\tilde g\in C_c^\infty$ with $0\le\tilde g\le1$, $\supp\tilde g\subset K'$ and $\tilde g>0$ near $E_1$. Let $\mathcal N$ be an infinite set with $b_n\to\beta$ and $b_n\ge\ell(E_1)+30\gamma_1$. As in the proof of Theorem~\ref{thm:U2}, Lemma~\ref{lem:arcsineloc} and Proposition~\ref{prop:fraction}, with $a=\ell(E_1)+\gamma_1$, give $c_1>0$ with
\[
\frac1{\pi T_n}\int_{I_n}\bigl\|\un_{\{|k|>N_n\}}G_z\bigr\|^2\,dE\ge f_0c_1,\qquad N_n=\bigl\lfloor e^{(b_n-\ell(E_1)-2\gamma_1)q_n}\bigr\rfloor,\quad T_n=|\delta_n|^{-1},
\]
for large $n\in\mathcal N$, where $I_n\subset K$. On $K$, $|c_J(E,T_n)|\ge g_1-C/T_n\ge g_1/2$ for large $n$. Lemma~\ref{lem:filter} with $J=M=1$ gives
\[
P_{\psi_g}(N_n,T_n)^{1/2}\ge\frac{g_1}2(f_0c_1)^{1/2}-CN_n^{-1}-CT_n^{-1}\ge c_2>0
\]
for large $n\in\mathcal N$, because $N_n\to\infty$. Hence $M_{p,\psi_g}(T_n)\ge c_2^2N_n^p$, and as before $\beta^+_{\psi_g}(p)\ge1-(\ell(E_1)+2\gamma_1)/\beta\ge1-(\ell(E_0)+\gamma)/\beta$.
\end{proof}

A finite partition of unity in energy turns the filtered upper bounds into a bound for the unfiltered packet, which closes the gap in Theorem~\ref{thm:U2}.

\begin{corollary}[The complete moment law]\label{cor:allmoments}
Let $0<\beta=\beta(\alpha)<\infty$ and $\theta\in G^v_\alpha$. Then
\[
\beta^+(p)=F_\beta(\ell_\theta)=F_\beta(\ell_*)\qquad\text{for every }p>0.
\]
In particular, if $\ell_*<\beta$, then $\ell_\theta=\ell_*$: the Lyapunov exponent attains its minimum over the spectrum on the support of $\mu_\theta$.
\end{corollary}

\begin{proof}
Let $S=\supp\mu_\theta$, a compact subset of $\Sigma$, and let $0<r<\ell_*/2$. Since $\ell$ is continuous, every $E\in S$ has an open interval around it on which $|\ell-\ell(E)|<r$. Finitely many such intervals, centred at points $E_1,\dots,E_k\in S$, cover $S$. Take real functions $g_i\in C_c^\infty(\R)$, each supported in the $i$-th interval, with $\sum_ig_i=1$ on $S$. Then $\supp g_i\subset\{|\ell-\ell(E_i)|<r\}$, and $\sum_ig_i(H)\delta_0=\delta_0$, because $\|(\sum_ig_i-1)(H)\delta_0\|^2=\int|\sum_ig_i-1|^2\,d\mu_\theta=0$. Put $\psi_i=g_i(H)\delta_0$ and discard the indices with $\psi_i=0$. By Theorem~\ref{thm:U3} and $\ell(E_i)\ge\ell_\theta$,
\[
\beta^+_{\psi_i}(p)\le F_\beta\bigl(\ell(E_i)-r\bigr)\le F_\beta(\ell_\theta-r).
\]
At every time, the Cauchy--Schwarz inequality for the finite sum $\langle\delta_j,e^{-itH}\delta_0\rangle=\sum_i\langle\delta_j,e^{-itH}\psi_i\rangle$ gives $|\langle\delta_j,e^{-itH}\delta_0\rangle|^2\le k\sum_i|\langle\delta_j,e^{-itH}\psi_i\rangle|^2$. Multiplying by $|j|^p$, summing over $j$ and averaging in time,
\[
M_p(T)\le k\sum_iM_{p,\psi_i}(T),
\]
where $M_{p,\psi_i}$ are the moments of the unnormalized vectors. Since the sum is finite, $\beta^+(p)\le\max_i\beta^+_{\psi_i}(p)\le F_\beta(\ell_\theta-r)$. Letting $r\to0$ and using the continuity of $F_\beta$ gives $\beta^+(p)\le F_\beta(\ell_\theta)$, and Theorem~\ref{thm:U2} gives equality. Hence $\au=\sup_p\beta^+(p)=F_\beta(\ell_\theta)$, while $\au=F_\beta(\ell_*)$ by Theorem~\ref{thm:tI-law}. If $\ell_*<\beta$, then $F_\beta(\ell_*)>0$ and $F_\beta$ is strictly decreasing on $(-\infty,\beta]$, so $F_\beta(\ell_\theta)=F_\beta(\ell_*)$ forces $\ell_\theta=\ell_*$. If $\ell_*\ge\beta$, both values are zero.
\end{proof}

So at every phase of $G^v_\alpha$ the fastest-tail exponent, every moment exponent, and the exponent of the fixed fraction in Theorem~\ref{thm:U2} are all governed by the same number $F_\beta(\ell_*)$.

\begin{theorem}[Local dimensions]\label{thm:localdim}
Let $0<\beta=\beta(\alpha)<\infty$. Then
\[
d^+_{\mu_\theta}(E)=F_\beta\bigl(\ell(E)\bigr)\qquad\text{for }\mu_\theta\text{-almost every }E.
\]
Consequently $\dim_P^+(\mu_\theta)=F_\beta(\ell_\theta)$ and $\dim_P^-(\mu_\theta)=F_\beta(\ell^\theta)$. In particular $\dim_P^-(\mu_\theta)=\dim_P^+(\mu_\theta)$ if and only if $F_\beta\circ\ell$ is constant $\mu_\theta$-almost everywhere.
\end{theorem}

\begin{proof}
\emph{Lower bound.} Fix $\gamma>0$. For a rational $E_0$ with $\ell(E_0)+30\gamma<\beta$ and a rational $\rho>0$ such that $|\ell-\ell(E_0)|<\gamma$ on $K=[E_0-2\rho,E_0+2\rho]$, put $K'=[E_0-\rho,E_0+\rho]$. By continuity of $\ell$, the interiors of these countably many intervals $K'$ cover $\{E:\ell(E)+31\gamma<\beta\}$. Fix one of them, let $g\ge0$ be smooth and positive exactly on $\operatorname{int}K'$, and let $\mathcal N$ be infinite with $b_n\to\beta$ and $b_n\ge\ell(E_0)+30\gamma$. Lemma~\ref{lem:arcsineloc} gives the weak convergence of $(\Delta_n)_*(g\mu_\theta)$ to $\sigma\int g\,d\mu_\theta$. The argument of Theorem~\ref{thm:U1}(2), with Proposition~\ref{prop:fraction}(2) for $a=\ell(E_0)+\gamma$, gives $d^+_{\mu_\theta}(E)\ge1-(\ell(E_0)+3\gamma/2)/\beta\ge1-(\ell(E)+3\gamma)/\beta$ for $\mu_\theta$-almost every $E\in\operatorname{int}K'$. Letting $\gamma\to0$ along a sequence gives $d^+_{\mu_\theta}\ge F_\beta\circ\ell$ almost everywhere.

\emph{Upper bound.} We use the theorem of Barbaroux, Germinet and Tcheremchantsev~\cite[Theorem~1 and Proposition~1(iii)]{BGT2000}: for a bounded operator on $\ell^2(\Z)$ and a normalized initial vector $\psi$ with compactly supported spectral measure, the exponents of the Ces\`aro-averaged moments satisfy $\beta^+_\psi(p)\ge D^+_{\mu_\psi}(1/(1+p))\ge\dim_P^+(\mu_\psi)$, where $D^+$ is the upper generalized R\'enyi dimension. No finite-support hypothesis on $\psi$ is made. Since $(2/T)e^{-2t/T}\ge(2e^{-2}/T)\un_{[0,T]}(t)$, our Abel-averaged moments dominate $2e^{-2}$ times the Ces\`aro-averaged ones, so the same inequality holds for our $\beta^+_\psi(p)$.

Fix $\gamma\in(0,\ell_*/2)$, a countable dense set $\mathcal D\subset\Sigma$, and for $E_0\in\mathcal D$ and rational $\rho>0$ with $|\ell-\ell(E_0)|<\gamma$ on $K=[E_0-2\rho,E_0+2\rho]$ put $K''=[E_0-\rho,E_0+\rho]$. By continuity of $\ell$ the interiors of these intervals cover $\Sigma$. Fix one, and let $g\in C_c^\infty$ be real with $\supp g\subset\operatorname{int}K$ and $g=1$ on $K''$. If $\psi_g=0$, then $\mu_\theta(K'')=0$. Otherwise Theorem~\ref{thm:U3} and the inequality above give $\dim_P^+(g^2\mu_\theta)\le F_\beta(\ell(E_0)-\gamma)$. For $E\in\operatorname{int}K''$ the measures $g^2\mu_\theta$ and $\mu_\theta$ agree near $E$, so they have the same upper local dimension at $E$. Hence $d^+_{\mu_\theta}(E)\le F_\beta(\ell(E_0)-\gamma)\le F_\beta(\ell(E)-2\gamma)$ for $\mu_\theta$-almost every $E\in\operatorname{int}K''$, and so for $\mu_\theta$-almost every $E$. Let $\gamma\to0$ and use the continuity of $F_\beta$.

The consequences follow from $\dim_P^+=\operatorname{ess\,sup}d^+$, $\dim_P^-=\operatorname{ess\,inf}d^+$, the monotonicity of $F_\beta$, and the identification of $\ell_\theta,\ell^\theta$ with the essential extrema of $\ell$.
\end{proof}

For $\beta(\alpha)=0$ every spectral measure $\mu_\theta$ has $d^+_{\mu_\theta}=0$ almost everywhere, at every phase, by Theorem~\ref{thm:tI-law} and the inequality of~\cite{BGT2000}. The equality extends to all spectral measures.

\begin{corollary}\label{cor:allmeasures}
Let $0<\beta(\alpha)<\infty$ and $\theta\in G^v_\alpha\cap(G^v_\alpha-\alpha)$. For every $\varphi\in\ell^2(\Z)$ with $\mu_\varphi\ne0$, $d^+_{\mu_\varphi}(E)=F_\beta(\ell(E))$ for $\mu_\varphi$-almost every $E$. In particular $\dim_P^+(\mu_\varphi)\le F_\beta(\ell_*)$.
\end{corollary}

\begin{proof}
Let $\mu_0$ and $\mu_1$ be the spectral measures of $\delta_0$ and $\delta_1$ at the phase $\theta$; as in the proof of Theorem~\ref{thm:U1}, $\mu_1=\mu_{\theta+\alpha}$. Theorem~\ref{thm:localdim} at the phases $\theta$ and $\theta+\alpha$ gives $d^+_{\mu_i}=F_\beta\circ\ell$ $\mu_i$-almost everywhere, for $i=0,1$. Let $\nu=\mu_0+\mu_1$, and recall the differentiation theorem for Radon measures~\cite[Theorem~2.12 and Corollary~2.14]{Mattila1995}: if $\mu\ll\nu$, then at $\mu$-almost every $E$ the ratio $\mu([E-r,E+r])/\nu([E-r,E+r])$ converges as $r\to0$ to the density $d\mu/d\nu(E)\in(0,\infty)$, and at such $E$ taking logarithms gives $d^+_\mu(E)=d^+_\nu(E)$. Applied to $\mu=\mu_i$, this gives $d^+_\nu=d^+_{\mu_i}=F_\beta\circ\ell$ on a set $A_i$ of full $\mu_i$-measure. The complement of $A_0\cup A_1$ is null for $\mu_0$ and for $\mu_1$, hence for $\nu$, so $d^+_\nu=F_\beta\circ\ell$ $\nu$-almost everywhere. The pair $\delta_0,\delta_1$ is cyclic, so $\mu_\varphi\ll\nu$, and the differentiation theorem applied to $\mu=\mu_\varphi$ gives $d^+_{\mu_\varphi}=d^+_\nu=F_\beta\circ\ell$ $\mu_\varphi$-almost everywhere. The last claim follows from $\ell\ge\ell_*$ on $\Sigma$.
\end{proof}

For the almost Mathieu operator $\ell=\log\lambda$ on $\Sigma$, so Corollary~\ref{cor:allmeasures} shows that every nonzero spectral measure is exactly packing-dimensional for $\theta\in\Ga\cap(\Ga-\alpha)$, in agreement with the companion paper~\cite{DMamo}. For a type~I potential with $\ell_\theta<\ell^\theta$ and $\ell_\theta<\beta(\alpha)$, Theorem~\ref{thm:localdim} gives $\dim_P^-(\mu_\theta)<\dim_P^+(\mu_\theta)$: the upper local dimension is a nonconstant explicit function of the energy.

\section{Beyond type I: a conditional complete law}\label{sec:beyond}

For a type~I potential the rational discriminant $t_q(\cdot,E)$ has exactly two zeros in a strip of width proportional to $q$, and near the drifting phase it behaves like a single cosine (Section~\ref{sec:typeI}). When the acceleration is at least two somewhere on the spectrum, $t_q(\cdot,E)$ has more zeros near the real axis, and they may cluster. The dynamical argument of Theorem~\ref{thm:tI-window} needs only one geometric property: near the drifting phase, the zero of $t_q(\cdot,E)$ that the phase crosses is not exponentially close to another zero at the same energy. This section makes that property precise, shows that it implies the transport law for every real-analytic potential with positive Lyapunov exponent on its spectrum, and reduces it, for an explicit two-harmonic family, to a large-deviation estimate for one polynomial.

Throughout this section $v$ is real analytic and $H_{v,\alpha}$ is supercritical, but no type, evenness, primitivity or acceleration hypothesis is made. We use the notation of Section~\ref{sec:typeI}: $t_q(x,E)$ is the rational discriminant in the variable $x=q\theta$, $\ell=L$ is the Lyapunov exponent, $\ell_*=\min_\Sigma\ell$, and $\mathcal N_q(\xi)=\{E\in\R:\dist(E,\Sigma)\le e^{-\xi q}\}$. Lemmas~\ref{lem:tI-bc} and~\ref{lem:tI-Q}, Corollary~\ref{cor:tI-bounds} and Theorem~\ref{thm:tI-window} do not use the type~I hypothesis. Put
\[
J=[-\|v\|_\infty-3,\ \|v\|_\infty+3],\qquad J_q=\{E\in\R:\dist(E,\Sigma)<1/q\}.
\]
Then $\Sigma\subset J$, and $J_q\subset J$ and $\mathcal N_q(\xi)\subset J_q$ for large $q$. The Lyapunov exponent is positive on $\Sigma$ by hypothesis and off $\Sigma$ by (A3), and it is continuous, so $\gamma_J:=\min_J\ell>0$. By the remark after Corollary~\ref{cor:tI-bounds}, for every $\zeta>0$ and all large $q$,
\begin{equation}\label{eq:UB}
|\partial_E^k\partial_x^jt_q(x,E)|\le k!\,(2q)^ke^{q(\ell(E)+\zeta)}\qquad(E\in J,\ x\in\R,\ k\ge0,\ j=0,1).
\end{equation}
Distances between points of $\C/\Z$ are taken in the metric of the cylinder.

\subsection{Clusters of zeros and weak separation}\label{sec:WXS}

\begin{definition}\label{def:WXS}
Let $q$ be a convergent denominator and $\eps>0$, and put $\tau=e^{-\eps q}$. For $E\in\R$ let $Z_q(E)$ be the multiset of zeros of $t_q(\cdot,E)$ in $\{x\in\C/\Z:|\operatorname{Im}x|<1\}$, counted with multiplicity. Let
\[
B_q(\eps)=\bigl\{x_0\in\T:\ \text{for some }E\in J_q,\ Z_q(E)\text{ has two entries }z,z'\text{ with }|z-x_0|<\tau\text{ and }|z-z'|<\tau\bigr\}.
\]
We say that \emph{WXS$(v,\alpha)$ holds} if for every $\eps>0$ and every $c>0$, almost every $\theta\in\T$ has an index $n_0$ such that $q_n\theta\bmod1\notin B_{q_n}(\eps)$ for all $n\ge n_0$ with $b_n\ge c$.
\end{definition}

The set $B_q(\eps)$ is open, since for fixed witnesses $E,z,z'$ the condition on $x_0$ is open, and it decreases as $\eps$ increases. Since $\theta\mapsto q_n\theta$ preserves Lebesgue measure, the Borel--Cantelli lemma shows that WXS$(v,\alpha)$ holds if $\sum_{n:\,b_n\ge c}|B_{q_n}(\eps)|<\infty$ for all $\eps,c>0$. If $\beta(\alpha)=0$, only finitely many $n$ have $b_n\ge c$, and WXS$(v,\alpha)$ holds trivially.

WXS is a \emph{horizontal} condition: it concerns the zeros of $t_q(\cdot,E)$ as functions of the phase at one fixed energy, real or complex. It excludes a real double zero (a turning point of a band centre) near $x_0$, two band centres crossing the energy $E$ at nearby phases, and a pair of complex zeros close to the real axis. It is different from a \emph{vertical} separation of band centres $E_b(x)$ at a fixed phase; Remark~\ref{rem:bandproduct} shows how the two are related.

\subsection{Factorization}\label{sec:EF}

\begin{lemma}\label{lem:trigpot}
Let $v(x)=\sum_{|k|\le d}v_ke^{2\pi ikx}$ be a real trigonometric polynomial with $v_d\ne0$, and let $p/q$ be a convergent of $\alpha$. Then $t_q(x,E)=\sum_{|k|\le d}c_{q,k}(E)e^{2\pi ikx}$, and
\[
c_{q,\pm d}(E)=(-v_{\pm d})^q\,e^{\pm\pi idp(q-1)}\qquad\text{for every }E\in\C .
\]
In particular the extreme coefficients do not depend on $E$, and $|c_{q,\pm d}|=|v_d|^q$.
\end{lemma}

The energy independence of the extreme coefficients is due to Shamis~\cite[Propositions~3.1 and~3.2]{Shamis2011}; we include the short computation in our normalization.

\begin{proof}
Each one-step matrix has entries of degree at most $d$ in $e^{2\pi i\varphi}$, so $\Delta_\varphi(E)$ is a trigonometric polynomial of degree at most $qd$ in $\varphi$. It is invariant under $\varphi\mapsto\varphi+1/q$, so only the harmonics $e^{2\pi imq\varphi}=e^{2\pi imx}$ with $|m|\le d$ survive. In the expansion of the trace as a sum over closed index paths, degree $qd$ in $e^{2\pi i\varphi}$ forces the $(1,1)$ entry and its top harmonic $-v_de^{2\pi id(\varphi+jp/q)}$ at every step $j=0,\dots,q-1$. Their product is $(-v_d)^qe^{2\pi idq\varphi}e^{2\pi idp\,q(q-1)/(2q)}=(-v_d)^qe^{2\pi idx}e^{\pi idp(q-1)}$. The coefficient of $e^{-2\pi idx}$ is computed in the same way.
\end{proof}

For a real trigonometric polynomial $T(x)=\sum_{|k|\le K'}c_ke^{2\pi ikx}$ of exact degree $K'\ge1$ (so $c_{-k}=\overline{c_k}$ and $c_{K'}\ne0$) write $\zeta_1,\dots,\zeta_{2K'}$ for its zeros in $\C/\Z$, counted with multiplicity, and $m(T)=\int_0^1\log|T(y)|\,dy$. Put
\[
\mathsf g(u)=|\sin\pi u|\,e^{-\pi|\operatorname{Im}u|}\qquad(u\in\C/\Z),\qquad c_0=2e^{-\pi}.
\]

\begin{lemma}[Factorization]\label{lem:EF}
Let $T$ be as above.
\begin{enumerate}
\item $\sum_k\operatorname{Im}\zeta_k=0$, and with suitable representatives $\zeta_k\in\C$, $T(x)=c_{K'}(2i)^{2K'}e^{\pi i\sum_k\zeta_k}\prod_k\sin\pi(x-\zeta_k)$.
\item $c_0\min\{\dist(u,\Z),1\}\le\mathsf g(u)\le1$, and $m(T)=\log|c_{K'}|+\pi\sum_k|\operatorname{Im}\zeta_k|$.
\item For every $x\in\C$, $|T(x)|\ge4^{K'}e^{m(T)-2\pi K'|\operatorname{Im}x|}\prod_k\mathsf g(x-\zeta_k)$, with equality for real $x$.
\item At a simple real zero $\zeta_1$, $|T'(\zeta_1)|=\pi4^{K'}e^{m(T)-\Phi}$ with $\Phi=-\sum_{k\ge2}\log\mathsf g(\zeta_1-\zeta_k)\ge0$.
\item $m(T)\ge\log\max_{x\in\T}|T(x)|-K'\log4$.
\end{enumerate}
\end{lemma}

\begin{proof}
1. With $w=e^{2\pi ix}$, $w^{K'}T$ is a polynomial of degree $2K'$ in $w$ with nonzero constant coefficient $c_{-K'}$, so its roots $w_k=e^{2\pi i\zeta_k}$ are nonzero, and $\prod_kw_k=c_{-K'}/c_{K'}$ has modulus one. Hence $\sum_k\log|w_k|=0$, which is $\sum_k\operatorname{Im}\zeta_k=0$. The factorization follows from $e^{2\pi ix}-e^{2\pi i\zeta}=e^{\pi i(x+\zeta)}\,2i\sin\pi(x-\zeta)$.

2. Write $u=a+ib$. Since $|\sin\pi u|^2=\sin^2\pi a+\sinh^2\pi b\le\cosh^2\pi b$, $\mathsf g\le1$. If $|b|\ge1$ then $|\sin\pi u|\ge\sinh\pi|b|\ge0.49e^{\pi|b|}$, so $\mathsf g(u)\ge0.49\ge c_0$. If $|b|<1$ then $|\sin\pi u|^2\ge4\dist(a,\Z)^2+\pi^2b^2\ge4\dist(u,\Z)^2$, so $\mathsf g(u)\ge2e^{-\pi}\dist(u,\Z)$. The classical integral $\int_0^1\log|\sin\pi(y-\zeta)|\,dy=\pi|\operatorname{Im}\zeta|-\log2$ and part~1, where $|e^{\pi i\sum\zeta_k}|=1$, give the formula for $m(T)$.

3. By part~1, $|T(x)|=|c_{K'}|4^{K'}\prod_k|\sin\pi(x-\zeta_k)|$, and $|\sin\pi(x-\zeta)|=\mathsf g(x-\zeta)e^{\pi|\operatorname{Im}(x-\zeta)|}\ge\mathsf g(x-\zeta)e^{\pi|\operatorname{Im}\zeta|-\pi|\operatorname{Im}x|}$, with equality for real $x$. Use part~2.

4. Differentiate the product of part~1 at $\zeta_1$, and use part~3 for the remaining factors, with $\operatorname{Im}\zeta_1=0$.

5. Evaluate part~3 at a real point where $|T|$ is maximal and use $\mathsf g\le1$.
\end{proof}

\begin{remark}[Band products]\label{rem:bandproduct}
Let $v$ be a trigonometric polynomial of degree $d$, $E$ real, and let $x$ be a simple real zero of $t_q(\cdot,E)$, so that $E=E_b(x)$ for one band centre. Since $t_q(x,\cdot)$ is monic with the simple roots $E_{b'}(x)$, $\partial_Et_q(x,E)=\prod_{b'\ne b}(E_b(x)-E_{b'}(x))$, and $\partial_xt_q=-\partial_Et_q\,E_b'$ by Lemma~\ref{lem:tI-bc}. Lemma~\ref{lem:EF}(4) therefore gives the exact finite-period identity
\[
\log|E_b'(x)|+\sum_{b'\ne b}\log|E_b(x)-E_{b'}(x)|=\log(\pi4^d)+m\bigl(t_q(\cdot,E)\bigr)-\Phi(x,E),
\]
where $\Phi(x,E)$ is the quantity of Lemma~\ref{lem:EF}(4) for $t_q(\cdot,E)$. The horizontal loss $\Phi$ involves the distances from $x$ to all other phase zeros at the same energy. Vertical gaps $|E_b(x)-E_{b'}(x)|$ enter only through the full product over all other bands, together with the slope $E_b'(x)$.
\end{remark}

\subsection{Truncation and a lower bound}\label{sec:LB}

Recall that $v$ extends holomorphically to $|\operatorname{Im}x|<\delta_v$ and that $0<\delta_1<\delta_v$ is fixed.

\begin{lemma}[Truncation]\label{lem:TR}
Let $\Gamma=\max\{\log\|T(E,x)\|:E\in J,\ |\operatorname{Im}x|\le\delta_1\}$ and $K=\lceil(\Gamma+2)/(2\pi\delta_1)\rceil$. For $E\in J$ write $t_q(x,E)=\sum_kc_{q,k}(E)e^{2\pi ikx}$ and $T_{q,E}(x)=\sum_{|k|\le K}c_{q,k}(E)e^{2\pi ikx}$. For all large $q$, every $E\in J$ and every $x$ with $|\operatorname{Im}x|\le1$,
\[
|t_q(x,E)-T_{q,E}(x)|\le e^{-q},
\]
and $|\partial_x(t_q-T_{q,E})(x)|\le2e^{-q}$ for $|\operatorname{Im}x|\le1/2$. If $v$ is a trigonometric polynomial of degree $d$, we take $K=d$ and $T_{q,E}=t_q(\cdot,E)$.
\end{lemma}

This is the truncation argument of Avila~\cite[Section~2, equations (10)--(11)]{Avila2015}, reproduced by Shamis~\cite[Lemma~2.1]{Shamis2011}.

\begin{proof}
For $|\operatorname{Im}x|\le q\delta_1$ the $q$ factors of the rational product have phases with imaginary part at most $\delta_1$, so $|t_q(x,E)|\le2e^{\Gamma q}$. Shifting the contour of the Fourier integral to $\operatorname{Im}x=\mp q\delta_1$ gives $|c_{q,k}(E)|\le2e^{\Gamma q-2\pi|k|q\delta_1}$. For $|\operatorname{Im}x|\le1$,
\[
\sum_{|k|>K}|c_{q,k}(E)|e^{2\pi|k|}\le4\sum_{k\ge K+1}e^{\Gamma q-2\pi k(q\delta_1-1)}\le8e^{2\pi(K+1)}e^{(\Gamma-2\pi(K+1)\delta_1)q}\le8e^{2\pi(K+1)}e^{-2q}
\]
for large $q$, since $2\pi K\delta_1\ge\Gamma+2$. This is at most $e^{-q}$ for large $q$. The derivative bound follows from Cauchy's estimate on disks of radius $1/2$. For a trigonometric polynomial of degree $d$, Lemma~\ref{lem:trigpot} shows that $t_q(\cdot,E)$ has no harmonics beyond $d$.
\end{proof}

For real $E$, $T_{q,E}$ is a real trigonometric polynomial of some exact degree $K'\le K$; Lemma~\ref{lem:LB} shows that it is not identically zero.

\begin{lemma}[Lower bound for the mean]\label{lem:LB}
For every $\eta>0$ there is $q_0$ such that for every convergent denominator $q\ge q_0$ and every $E\in J$,
\[
m(T_{q,E})\ge q\bigl(\ell(E)-\eta\bigr).
\]
\end{lemma}

\begin{proof}
For a rational $p/q$ let $L(p/q,E)=q^{-1}\int_\T\log\rho\bigl(A^{(p/q)}_q(E,\varphi)\bigr)\,d\varphi$, where $\rho$ is the spectral radius; this is the Lyapunov exponent of the frequency $p/q$, averaged over the phase. Bourgain and Jitomirskaya proved that $(\omega,E)\mapsto L(\omega,E)$, with this definition at rational $\omega$, is jointly continuous at $(\alpha,E)$ for every irrational $\alpha$ and real $E$~\cite[Theorem~1]{BourgainJitomirskaya2002}. By compactness of $J$, $L(p_n/q_n,E)\to\ell(E)$ uniformly on $J$. For $M\in SL(2,\R)$, $\rho(M)\le\max\{1,|\tr M|\}$: the spectral radius is one if $M$ is elliptic or parabolic, and smaller than $|\tr M|$ if $M$ is hyperbolic. Since $\Delta_\varphi$ is $1/q$-periodic, $qL(p/q,E)\le\log\max\{1,\max_x|t_q(x,E)|\}$. It suffices to prove the lemma for $\eta_1=\min\{\eta,\gamma_J\}$, since its conclusion for $\eta_1$ implies the one for $\eta$. For large $q$, $L(p/q,E)\ge\ell(E)-\eta_1/2\ge\gamma_J/2>0$ on $J$, so $\max_x|t_q(x,E)|\ge e^{q(\ell(E)-\eta_1/2)}$ and $\max_x|T_{q,E}(x)|\ge\frac12e^{q(\ell(E)-\eta_1/2)}$ by Lemma~\ref{lem:TR}. Lemma~\ref{lem:EF}(5) gives $m(T_{q,E})\ge q(\ell(E)-\eta_1/2)-\log2-K\log4\ge q(\ell(E)-\eta_1)$ for large $q$.
\end{proof}

\subsection{Geometry at separated phases}\label{sec:NG}

\begin{lemma}\label{lem:NGe}
Let $K$ be as in Lemma~\ref{lem:TR}, and put
\[
a_0=\frac{1+\gamma_J}{8K},\qquad\rho_K=\frac1{148(2K-1)},\qquad\lambda_K=\frac12\Bigl(\frac{c_0\rho_K}8\Bigr)^{2K}.
\]
Let $0<\eps<a_0/2$ and $0<\eta\le\gamma_J/2$. For all large $q$, every $x_0\in\T\setminus B_q(\eps)$ and every $E\in J_q$, with $\tau=e^{-\eps q}$ and $S_0=\{x\in\T:|x-x_0|\le\tau/16\}$, one of the following holds.
\begin{enumerate}
\item[(i)] $|t_q(x,E)|\ge\lambda_Ke^{q(\ell(E)-\eta)}\tau^{2K}$ for every $x\in S_0$.
\item[(ii)] $t_q(\cdot,E)$ has a simple real zero $z_E$ with $|z_E-x_0|<\tau/2$. On $I_E=\{x\in\R:|x-z_E|\le\rho_K\tau/4\}$ the derivative $\partial_xt_q(\cdot,E)$ has a constant sign and modulus at least $S_E:=\frac\pi4e^{q(\ell(E)-\eta)}(c_0\tau/2)^{2K-1}$, and $|t_q(x,E)|\ge\lambda_Ke^{q(\ell(E)-\eta)}\tau^{2K}$ for $x\in S_0\setminus I_E$.
\end{enumerate}
\end{lemma}

\begin{proof}
Let $q$ be large, $x_0\notin B_q(\eps)$, $E\in J_q$, and write $T=T_{q,E}$, of exact degree $K'\le K$. By Lemma~\ref{lem:LB}, $m(T)\ge q(\ell(E)-\eta)$, and $\ell(E)-\eta\ge\gamma_J/2$. If $K'=0$, then $T$ is a constant of modulus $e^{m(T)}$, $t_q(\cdot,E)$ has no zeros in $|\operatorname{Im}x|\le1/2$, and (i) holds. Assume $K'\ge1$.

\emph{Matching of zeros.} Let $h=e^{-a_0q}$ and let $U$ be the union of the open disks of radius $h$ about the zeros of $T$ in $\C/\Z$. On the boundary of $U$, and outside $U$, every zero of $T$ is at distance at least $h$, so by Lemma~\ref{lem:EF}(2),(3), for $|\operatorname{Im}x|\le1/2$,
\[
|T(x)|\ge4^{K'}e^{m(T)-\pi K'}(c_0h)^{2K'}\ge e^{-\pi K}c_0^{2K}e^{q(\gamma_J/2-(1+\gamma_J)/4)}>e^{-q}\ge|t_q(x,E)-T(x)|
\]
for large $q$, since $\gamma_J/2-(1+\gamma_J)/4>-1$. A component $W$ of $U$ has diameter at most $4Kh$. If it meets $\{|\operatorname{Im}x|\le1/4\}$, its closure lies in $\{|\operatorname{Im}x|<1/2\}$, and Rouch\'e's theorem shows that $t_q(\cdot,E)$ and $T$ have the same number of zeros in $W$. Moreover $t_q(\cdot,E)$ has no zeros in $\{|\operatorname{Im}x|\le1/4\}\setminus U$. Since $\eps<a_0/2$, $8Kh\le\tau/8$ for large $q$.

\emph{Case A: $t_q(\cdot,E)$ has no zero $z$ with $|z-x_0|<\tau/2$.} A zero $\zeta$ of $T$ with $|\zeta-x_0|<3\tau/8$ would lie in a component meeting $\{|\operatorname{Im}x|\le1/4\}$, which would contain a zero of $t_q(\cdot,E)$ within $3\tau/8+4Kh<\tau/2$ of $x_0$. So every zero of $T$ is at distance at least $3\tau/8$ from $x_0$, hence at least $\tau/4$ from $S_0$. Lemma~\ref{lem:EF}(3) gives $|T|\ge e^{q(\ell(E)-\eta)}(c_0\tau/4)^{2K}$ on $S_0$, and $|t_q|\ge|T|-e^{-q}$ gives (i), since $\ell(E)-\eta-2K\eps>-1$.

\emph{Case B: $t_q(\cdot,E)$ has a zero $z$ with $|z-x_0|<\tau/2$.} Since $x_0\notin B_q(\eps)$, no other entry $z'$ of $Z_q(E)$ satisfies $|z-z'|<\tau$. In particular $z$ is simple, and $z$ is real: otherwise $\bar z$ would be a second entry with $|\bar z-z|=2|\operatorname{Im}z|<\tau$. Let $W$ be the component of $U$ containing $z$. It contains no other zero of $t_q(\cdot,E)$, because its diameter is less than $\tau$, so it contains exactly one zero $\zeta$ of $T$, which is simple, with $|\zeta-z|\le4Kh$. Since $U$ is symmetric under conjugation and $W$ contains the real point $z$, $\overline W=W$, so $\bar\zeta=\zeta$ is real. If another zero $\zeta'$ of $T$ had $|\zeta'-\zeta|<\tau-8Kh$, its component would differ from $W$ and contain a zero $z''\ne z$ of $t_q(\cdot,E)$ with $|z''-z|<\tau$, which is excluded. So every other zero of $T$ is at distance at least $s:=\tau-8Kh\ge7\tau/8$ from $\zeta$.

Write $T(x)=\sin\pi(x-\zeta)\,G(x)$. For real $x$ the function $G$ is real, and $G'/G=\pi\sum_{k\ge2}\cot\pi(x-\zeta_k)$. If $\dist(u,\Z)\ge s/2$, then $|\cot\pi u|\le1.01$ when $|\operatorname{Im}u|\ge1$, and $|\cot\pi u|\le\cosh\pi/(2\dist(u,\Z))\le11.6/s$ otherwise, by the bound in the proof of Lemma~\ref{lem:EF}(2). Hence $|G'/G|\le37(2K-1)/s$ for real $x$ with $|x-\zeta|\le s/2$. On $I_T=\{x\in\R:|x-\zeta|\le\rho_Ks\}$ we get $|\sin\pi(x-\zeta)|\,|G'/G|\le\pi/4$, $\pi\cos\pi(x-\zeta)\ge0.99\pi$ and $|G(x)|\ge e^{-1/4}|G(\zeta)|$. So $T'=G\,(\pi\cos\pi(x-\zeta)+\sin\pi(x-\zeta)\,G'/G)$ has the constant sign of $G$ and $|T'|\ge0.74\pi e^{-1/4}|G(\zeta)|\ge|T'(\zeta)|/2$ on $I_T$. By Lemma~\ref{lem:EF}(4) and $\mathsf g(\zeta-\zeta_k)\ge c_0s$,
\[
|T'(\zeta)|\ge\pi4^{K'}e^{m(T)}(c_0s)^{2K'-1}\ge\pi e^{q(\ell(E)-\eta)}(c_0\tau/2)^{2K-1}=4S_E.
\]
Since $|\partial_x(t_q-T)|\le2e^{-q}$ and $S_E\gg e^{-q}$, $\partial_xt_q(\cdot,E)$ has the sign of $T'$ and modulus at least $|T'(\zeta)|/4\ge S_E$ on $I_T$. Put $z_E=z$. Then $I_E\subset I_T$, because $\rho_K\tau/4+4Kh\le\rho_Ks$.

Finally let $x\in S_0\setminus I_E$. Then $|x-\zeta|\ge\rho_K\tau/4-4Kh\ge\rho_K\tau/8$. Every other zero of $T$ is at distance at least $s-|x-\zeta|\ge7\tau/8-(\tau/16+\tau/2+4Kh)\ge\tau/4$ from $x$. So all zeros of $T$ are at distance at least $\rho_K\tau/8$ from $x$, and Lemma~\ref{lem:EF}(3) gives $|T(x)|\ge e^{q(\ell(E)-\eta)}(c_0\rho_K\tau/8)^{2K}=2\lambda_Ke^{q(\ell(E)-\eta)}\tau^{2K}$. Subtracting $e^{-q}$ gives the last claim of (ii). In Case~A the same bound holds, since $c_0\tau/4\ge c_0\rho_K\tau/8$.
\end{proof}

\begin{corollary}\label{cor:WXStoW}
Let $\xi,\zeta>0$ and $0<\eps\le\min\{\zeta/(16K),\,a_0/4\}$. For all large $q$ and every $x_0\in\T\setminus B_q(\eps)$, (W) holds at $(q,x_0,e^{-2K\eps q})$ with these $\xi,\zeta$.
\end{corollary}

\begin{proof}
Let $\sigma=\tau^{2K}=e^{-2K\eps q}$, so $e^{-\zeta q/4}\le\sigma\le1$. Apply Lemma~\ref{lem:NGe} with $\eta=\min\{\zeta,\gamma_J\}/2$, and let $E\in\mathcal N_q(\xi)\subset J_q$. A point $x$ with $|x-x_0|\le\sigma/16$ lies in $S_0$. In case (i), and in case (ii) for $x\notin I_E$, $|t_q(x,E)|\ge\lambda_Ke^{q(\ell(E)-\eta)}\tau^{2K}\ge\sigma^2e^{q(\ell(E)-\zeta)}$ for large $q$, since $\eta\le\zeta$ and $\tau^{2K}\to0$; so such $x$ do not lie in the set $Z_E$ of Definition~\ref{def:tI-W}. In case (i) therefore $Z_E=\emptyset$. In case (ii) every $x\in Z_E$ lies in $I_E$, where $t_q(\cdot,E)$ is monotone with slope at least $S_E$ and vanishes at $z_E$, so
\[
|x-z_E|\le\frac{|t_q(x,E)|}{S_E}<\frac{\sigma^2e^{q(\ell(E)-\zeta)}}{S_E}\le C_K\tau^{2K+1}\le\frac\sigma{16}
\]
for large $q$. Let $J_E=\{|x-z_E|\le\sigma/8\}$, an arc of length $\sigma/4$ centred at $z_E$. It lies in $I_E$ for large $q$, and on it $t_q(\cdot,E)$ is strictly monotone with $|\partial_xt_q|\ge S_E\ge\sigma e^{q(\ell(E)-\zeta)}$. This is (W).
\end{proof}

\subsection{The conditional law}\label{sec:CP}

\begin{theorem}\label{thm:M2CP}
Let $v$ be real analytic and $H_{v,\alpha}$ supercritical. Fix $c>0$, $0<d<B$, $\eps_{\mathrm{out}}>0$ and $m>0$, let $\xi,\zeta$ be as in Theorem~\ref{thm:tI-window} with $\eps_{\mathrm{out}}$ in place of $\eps$, and let $0<\eps\le\min\{\zeta/(16K),a_0/4\}$. There are $C$ and $n_1$ such that for every $n\ge n_1$ with $b_n\ge c$, every phase $\theta$ with $q_n\theta\bmod1\notin B_{q_n}(\eps)$, and every $T$ with $e^{dq_n}\le T\le e^{Bq_n}$,
\[
P(R_n(T),T)\le CT^{-m},\qquad R_n(T)=e^{\eps_{\mathrm{out}}q_n}\Bigl[1+\min\bigl(Te^{-\ell_*q_n},\,e^{(b_n-\ell_*)_+q_n}\bigr)\Bigr].
\]
\end{theorem}

\begin{proof}
By Corollary~\ref{cor:WXStoW}, (W) holds at $(q_n,q_n\theta,e^{-2K\eps q_n})$ for all large $n$, uniformly in $\theta$, and $e^{-2K\eps q_n}\ge e^{-\zeta q_n/4}$. Theorem~\ref{thm:tI-window}, which does not use the type~I hypothesis, applies.
\end{proof}

\begin{theorem}\label{thm:CP}
Let $v$ be real analytic, $H_{v,\alpha}$ supercritical, and assume WXS$(v,\alpha)$. Put $\beta=\beta(\alpha)$.
\begin{enumerate}
\item If $\beta<\infty$, then for almost every $\theta$, $P(T^a,T)=O(T^{-m})$ for every $a>F_\beta(\ell_*)$ and $m>0$, and $\au(H_{v,\alpha,\theta})=F_\beta(\ell_*)$.
\item If $\beta=0$ or $\beta=\infty$, then $\au(H_{v,\alpha,\theta})=F_\beta(\ell_*)$ for every $\theta$.
\item $\au(H_{v,\alpha,\theta})\ge F_\beta(\ell_*)$ for every $\theta$.
\end{enumerate}
\end{theorem}

\begin{proof}
For $k\in\N$ let $G_k$ be the full-measure set of phases given by WXS$(v,\alpha)$ with $\eps=c=1/k$, and let $G=\bigcap_kG_k$. Fix $\theta\in G$ and $\beta<\infty$. Theorem~\ref{thm:tI-M1} gives (M1). For (M2$_{\ell_*}$), fix $c,d,B,\eps_{\mathrm{out}},m$, let $\eps$ be given by Theorem~\ref{thm:M2CP}, and choose $k$ with $1/k\le\min\{\eps,c\}$. Since $\theta\in G_k$, $q_n\theta\notin B_{q_n}(1/k)$ for all large $n$ with $b_n\ge1/k$, in particular for all large $n$ with $b_n\ge c$. Since $B_{q_n}(\cdot)$ decreases, $B_{q_n}(\eps)\subset B_{q_n}(1/k)$. So Theorem~\ref{thm:M2CP} applies at all large $n$ with $b_n\ge c$, which is (M2$_{\ell_*}$). Theorem~\ref{thm:covering} gives the upper bound in part~1. The lower bound in part~1 and part~3 for $0<\beta<\infty$ are Theorem~\ref{thm:L}, which applies because $\ell$ is continuous. For $\beta=0$ use Corollary~\ref{cor:tI-beta0}, which does not use the type~I hypothesis, and for $\beta=\infty$ use~\cite[Theorem~3(2)]{JZ2022} as in the proof of Theorem~\ref{thm:tI-law}.
\end{proof}

\begin{theorem}[The complete law under weak separation]\label{thm:complete}
Let $v$ be real analytic with $L>0$ on its spectrum. Let $s\ge1$ be the greatest common divisor of the frequencies $k\ne0$ with $\hat v(k)\ne0$, and let $w$ be the real-analytic function with $v(x)=w(sx)$. If WXS$(w,s\alpha)$ holds, then for almost every $\theta$,
\[
\au(H_{v,\alpha,\theta})=F_{\beta(s\alpha)}(\ell_*),
\]
where $\ell_*$ is the minimum of the Lyapunov exponent on the spectrum. If $\beta(\alpha)\in\{0,\infty\}$, this holds for every $\theta$, without any hypothesis.
\end{theorem}

\begin{proof}
A constant potential has zero Lyapunov exponent on its spectrum, so $v$ is not constant and $s$ is defined. By Theorem~\ref{thm:tI-NP} and its proof, $H_{v,\alpha,\theta}=H_{w,s\alpha,s\theta}$, and the two families have the same spectrum and the same Lyapunov exponent; so $H_{w,s\alpha}$ is supercritical with the same $\ell_*$. Theorem~\ref{thm:CP} applies to $(w,s\alpha)$ and gives a full-measure set $G$ of phases. The map $\theta\mapsto s\theta$ preserves Lebesgue measure, so $\{\theta:s\theta\in G\}$ has full measure. By Lemma~\ref{lem:tI-betas}, $\beta(s\alpha)\in\{0,\infty\}$ exactly when $\beta(\alpha)\in\{0,\infty\}$, and then Theorem~\ref{thm:CP}(2) needs no hypothesis.
\end{proof}

If $w$ is of type~I at the frequency $s\alpha$, the conclusion of Theorem~\ref{thm:complete} already follows from Theorem~\ref{thm:tI-NP}, without WXS. At every supercritical spectral energy the acceleration is a positive integer (proof of Lemma~\ref{lem:tI-T}). So WXS$(w,s\alpha)$ is the missing hypothesis exactly for the primitive potentials $w$ whose acceleration is at least two at some point of the spectrum.

\begin{remark}[Why the primitive part]\label{rem:primitive}
WXS$(v,\alpha)$ itself can fail when $v$ is not primitive. Let $v=w(s\cdot)$ with $s\ge2$, assume that $H_{w,s\alpha}$ is supercritical and of type~I, and assume that infinitely many convergent denominators $q_n$ of $\alpha$ are divisible by $s$ and satisfy $b_n\ge c$ for a fixed $c>0$. At such a scale, $q'=q_n/s$ is a convergent denominator of $s\alpha$ with the same $\delta_n$ (Proposition~\ref{prop:tI-cheb}), and $t^{(v)}_{q_n}=2T_s(t^{(w)}_{q'}/2)$. Let $\theta$ satisfy $s\theta\in G^w_{s\alpha}$, a full-measure condition, and let $x_0=q_n\theta=q'(s\theta)$, which is good for $w$ at scale $q'$. Let $E$ be a band centre of $w$ at the phase $x_0$; it lies within $Ce^{-b_nq_n/2}$ of $\Sigma$ by Lemma~\ref{lem:tI-loc}(1). By Lemma~\ref{lem:tI-Wgood} for $w$, $t^{(w)}_{q'}(\cdot,E)$ is monotone near $x_0$ with slope at least $q'^{-3}e^{q'(\ell(E)-\zeta)}$, so it takes the $s$ values $2\cos((2j-1)\pi/(2s))$ at $s$ distinct real points within $4q'^3e^{-q'(\ell(E)-\zeta)}$ of $x_0$. These are zeros of $t^{(v)}_{q_n}(\cdot,E)$. Hence $q_n\theta\in B_{q_n}(\eps)$ for every $\eps<\ell_*/(2s)$ at all large such $n$, and WXS$(v,\alpha)$ fails. This failure needs the resonant scales: if $\beta(\alpha)=0$, WXS holds trivially for every potential. It also does not locate a change of the exponent, since WXS$(v,\alpha)$ can fail while both $F_{\beta(\alpha)}(\ell_*)$ and $F_{\beta(s\alpha)}(\ell_*)$ vanish.
\end{remark}

\subsection{An explicit two-harmonic family}\label{sec:twoharmonic}

Let
\[
v(x)=2\lambda\cos2\pi x+2\mu\cos4\pi x,\qquad\lambda\ne0,\ \mu\ne0,
\]
and assume that $H_{v,\alpha}$ is supercritical. This potential is primitive and even. By Lemma~\ref{lem:trigpot} and the evenness of $t_q$ (Section~\ref{sec:tI-disc}),
\[
t_q(x,E)=P_q(E)+2C_q(E)\cos2\pi x+2B_q\cos4\pi x,\qquad B_q=(-\mu)^q,
\]
with real polynomials $P_q$ of degree $q$ and $C_q$ of degree at most $q-1$. With $c=\cos2\pi x$, $t_q=4B_qc^2+2C_qc+(P_q-2B_q)$, whose roots in $c$ are $c_\pm=\bar c_q\pm\frac12\sqrt{\tilde\Delta_q}$, where
\[
\bar c_q(E)=-\frac{C_q(E)}{4B_q},\qquad\tilde\Delta_q(E)=\Bigl(\frac{C_q(E)}{2B_q}\Bigr)^2-\frac{P_q(E)-2B_q}{B_q}=(c_+-c_-)^2 .
\]
The four phase zeros of $t_q(\cdot,E)$ are $\pm x_+$ and $\pm x_-$, with $\cos2\pi x_\pm=c_\pm$. The polynomial $\tilde\Delta_q$, of degree at most $2q$, measures how far $t_q$ is from a perfect square in $\cos2\pi x$.

\begin{definition}\label{def:EDelta}
We say that \emph{(E-$\Delta$) holds} if for every $c>0$ and every $\eps>0$ there is $\eps'>0$ such that for all large $n$ with $b_n\ge c$,
\[
\bigl|\bigl\{E\in J:\ |\tilde\Delta_{q_n}(E)|\le e^{-\eps q_n},\ |\bar c_{q_n}(E)|\le2\bigr\}\bigr|\le e^{-\eps'q_n}.
\]
\end{definition}

\begin{proposition}\label{prop:EDelta}
For the family above, (E-$\Delta$) implies WXS$(v,\alpha)$. Consequently, under (E-$\Delta$), $\au(H_{v,\alpha,\theta})=F_{\beta(\alpha)}(\ell_*)$ for almost every $\theta$.
\end{proposition}

\begin{proof}
Fix $c,\eps>0$, let $q=q_n$ with $b_n\ge c$ be large, $\tau=e^{-\eps q}$, and let $x_0\in B_q(\eps)$, with witnesses $E\in J_q$ and entries $z,z'$ of $Z_q(E)$. The entries are two of the four zeros $\pm x_\pm$.

\emph{Mirror pairs.} If $\{z,z'\}=\{x_+,-x_+\}$ or $\{x_-,-x_-\}$, then $\dist(2z,\Z)=|z-z'|<\tau$, so $z$ is within $\tau/2$ of $\{0,\frac12\}$ and $x_0$ is within $2\tau$ of $\{0,\frac12\}$. These phases form a set of measure at most $8\tau$.

\emph{Cross pairs.} Otherwise $\cos2\pi z=c_+$ and $\cos2\pi z'=c_-$, or the reverse. Since $|\operatorname{Im}z|<\tau$ and $|\operatorname{Im}z'|<2\tau$, the derivative of $\cos2\pi\cdot$ is at most $2\pi\cosh4\pi\tau$ on the segment between them, so $|c_+-c_-|\le7\tau$ for large $q$. Hence $|\tilde\Delta_q(E)|\le49\tau^2\le e^{-\eps q}$, $|\cos2\pi x_0-\bar c_q(E)|\le|\cos2\pi x_0-\cos2\pi z|+|c_+-c_-|/2\le11\tau$, and $|\bar c_q(E)|\le2$. So $E$ lies in
\[
E_{\mathrm{bad}}=\{E\in J:|\tilde\Delta_q(E)|\le e^{-\eps q},\ |\bar c_q(E)|\le2\},
\]
and $\cos2\pi x_0$ lies in the $11\tau$-neighbourhood $S'$ of $S=\bar c_q(E_{\mathrm{bad}})$.

\emph{The set $S$.} Since $\tilde\Delta_q$ and $\bar c_q$ are real polynomials of degrees at most $2q$ and $q-1$, $E_{\mathrm{bad}}$ is a union of at most $3q+2$ intervals. On $E_{\mathrm{bad}}$ the roots satisfy $|c_\pm|\le2+e^{-\eps q/2}/2\le2.5$, so the four phase zeros have $|\operatorname{Im}x|\le\operatorname{arsinh}(2.5)/(2\pi)\le0.3$. By Lemma~\ref{lem:EF}(2), with $|c_{q,\pm2}|=|\mu|^q$, $m(t_q(\cdot,E))\le q\log|\mu|+1.2\pi\le q\log|\mu|+3.8$. Let $\eta,\zeta>0$. Lemma~\ref{lem:LB} gives $\ell(E)\le\log|\mu|+\eta+3.8/q$ on $E_{\mathrm{bad}}$, for large $q$. Since $C_q=(t_q(0,\cdot)-t_q(\frac12,\cdot))/4$, \eqref{eq:UB} gives $|\bar c_q'(E)|=|C_q'(E)|/(4|\mu|^q)\le Cqe^{(\eta+\zeta)q}$ on $E_{\mathrm{bad}}$. Integrating over each component, $S$ is a union of at most $3q+2$ intervals with $|S|\le Cqe^{(\eta+\zeta)q}|E_{\mathrm{bad}}|$, and $|S'|\le|S|+22\tau(3q+2)$.

\emph{Measure of the bad phases.} For $0<r\le1$, the set $\{x\in\T:|\sin2\pi x|<r\}$ has measure at most $4r$. On its complement, $x\mapsto\cos2\pi x$ is at most two-to-one with derivative of modulus at least $2\pi r$, so $\{x:\cos2\pi x\in S',\ |\sin2\pi x|\ge r\}$ has measure at most $|S'|/(\pi r)$. Altogether
\[
|B_q(\eps)|\le8\tau+4r+\frac1{\pi r}\Bigl[Cqe^{(\eta+\zeta)q}|E_{\mathrm{bad}}|+22\tau(3q+2)\Bigr].
\]
Let $\eps'$ be given by (E-$\Delta$) for $c$ and $\eps$; decreasing it, we may assume $\eps'<2\eps$. Take $r=e^{-\eps'q/2}$ and $\eta+\zeta<\eps'/4$. Every term is at most $Cqe^{-c'q}$ with $c'=\min\{\eps'/4,\eps-\eps'/2\}>0$. Since $\sum_nq_ne^{-c'q_n}<\infty$, the Borel--Cantelli lemma gives WXS$(v,\alpha)$. The last statement is Theorem~\ref{thm:CP}.
\end{proof}

For the non-primitive member $\lambda=0$ and even $q$, the Chebyshev identity of Proposition~\ref{prop:tI-cheb} gives exactly $\tilde\Delta_q=2/\mu^q$, which is exponentially small when $|\mu|>1$; this is the structural clustering of Remark~\ref{rem:primitive}. Whether (E-$\Delta$) holds for $\lambda\ne0$ is open.

\begin{remark}[Numerical observations]\label{rem:numerics}
The following finite computations are not used in any proof. For the family above with $\mu=1.8$, $\lambda\in\{0.6,\,0.02\}$ and the frequency $161/242$, and with $\lambda=0.6$ and the frequency $321/482$, we computed, at $32$, $32$ and $16$ equally spaced phases respectively and for each of the two lowest bands, the band centres, their slopes, the four phase zeros and the loss $\Phi$ of Remark~\ref{rem:bandproduct}, in arithmetic with $100$ decimal digits. All phase zeros were real. The vertical gaps between the two lowest band centres were as small as $4.0\cdot10^{-12}$, $1.1\cdot10^{-11}$ and $3.6\cdot10^{-21}$, and the slopes $|E_b'|$ at most $3.1\cdot10^{-10}$, $7.9\cdot10^{-10}$ and $9.3\cdot10^{-20}$. Yet the largest loss $\Phi$ was $4.48$, $3.69$ and $3.54$, and the smallest distance between two phase zeros was $0.0085$, $0.0098$ and $0.0098$. The identity of Remark~\ref{rem:bandproduct} held to a relative error below $10^{-76}$. These are observations at finitely many phases of two rational frequencies, which belong to the sequence $(2k+1)/(3k+2)$ converging to the rational number $2/3$, not to the convergents of a fixed irrational frequency. At these samples the tiny vertical gaps did not produce horizontal clusters. They say nothing about the behaviour along the convergents of a fixed irrational frequency. Whether exponentially small vertical gaps recur, at a positive exponential rate, along the convergents of a fixed irrational frequency, and what they then do to the horizontal separation, is part of the open question of WXS.
\end{remark}

\section{Discussion and open problems}\label{sec:discussion}

\paragraph{The dynamical picture.} Theorems~\ref{thm:A} and~\ref{thm:B} identify the upper Abel-averaged transport exponents of every supercritical type~I analytic operator, for every moment, at every phase of the full-measure set $G^v_\alpha$: all of them equal $\max\{0,1-\ell_*/\beta(\alpha)\}$ (Corollary~\ref{cor:allmoments}). For the almost Mathieu operator, Damanik and Tcheremchantsev proved that the lower transport exponents vanish for every $\lambda>1$ and every frequency~\cite{DT2007}, so the fastest lower and upper exponents differ in the window $1<\lambda<e^{\beta}$. The proofs suggest a picture of strong intermittency. Near the resonant times $T\approx|q_n\alpha-p_n|^{-1}$ of strongly Liouville scales, the part of the packet at energy $E$ makes a burst, moving with speed about $e^{-L(E)q_n}$ and reaching distance $T^{1-L(E)/\beta_n}$; a fixed fraction of the packet takes part (Theorem~\ref{thm:B}(1)). The exponents record the largest of these bursts. We prove nothing about the proportion of times spent in each regime, and nothing about the lower exponents of potentials other than the almost Mathieu operator.

\paragraph{The open hypothesis: weak separation beyond type~I.} Theorem~\ref{thm:D} reduces the law $\max\{0,1-\ell_*/\beta(s\alpha)\}$ for every real-analytic potential with positive Lyapunov exponent on its spectrum to one hypothesis, WXS$(w,s\alpha)$ for its primitive part $w$, and the hypothesis is needed only when $w$ has acceleration at least two at some point of the spectrum. The open problem is therefore exactly this: \emph{does WXS hold for primitive real-analytic potentials with acceleration at least two?} WXS is a horizontal condition. It asks that, at each fixed energy near the spectrum, no two zeros of the rational discriminant in the phase variable, real or complex, cluster at an exponential rate near the drifting phase $q_n\theta$. It is not a condition on vertical gaps between band centres at a fixed phase: by Remark~\ref{rem:bandproduct}, vertical gaps enter the horizontal loss only through a product over all bands. Whether exponentially small vertical gaps recur, at a positive exponential rate, along the convergents of a fixed irrational frequency, and whether they then force horizontal clusters, is part of the open question. For the two-harmonic family $2\lambda\cos2\pi x+2\mu\cos4\pi x$, Proposition~\ref{prop:EDelta} reduces WXS to the phase-free estimate (E-$\Delta$) on a single polynomial of degree at most $2q$, which we also leave open. A proof of WXS, or a fixed frequency at which the law fails together with the corrected arithmetic invariant, would complete the picture for real-analytic potentials. The law could also hold where WXS fails; our methods do not decide this.

\paragraph{Exceptional phases.} The upper bounds are proved for $\theta\in G^v_\alpha$, the phases whose local phase $q_n\theta$ stays polynomially far from the turning points of the band centres. At other phases a turning point near the origin changes the geometry of the discriminant, and the window estimate must be modified. Whether the exponent can be larger at such phases is open.

\paragraph{Lower local dimensions.} Theorem~\ref{thm:B}(3) determines the upper local dimension of $\mu_\theta$. The lower local dimension and the Hausdorff dimension of spectral measures are not addressed here.

\paragraph{Critical and mixed spectrum.} The law concerns potentials whose Lyapunov exponent is positive on the whole spectrum. When the Lyapunov exponent vanishes somewhere on the spectrum and $\beta(\alpha)>0$, transport is quasiballistic by the theorem of Jitomirskaya and Zhang~\cite[Theorem~3]{JZ2022}. At critical energies with $\beta(\alpha)=0$ our methods give nothing: the resonant scales disappear, and the speed bound $e^{-L(E)q}$ becomes vacuous.

\paragraph{Physical meaning.} Laboratory realizations of the Aubry--Andr\'e model use frequencies such as the golden mean, for which $\beta=0$; Corollary~\ref{cor:tI-beta0} confirms the absence of power-law transport for every analytic potential whose Lyapunov exponent is positive on the spectrum. The theorems show that this absence is arithmetic, not a consequence of positive Lyapunov exponent alone, and that it does not depend on the details of the potential. The proofs indicate that a frequency with one exceptionally large partial quotient after a denominator $q$ produces a transient burst over the time window $e^{L(E)q}\lesssim T\lesssim|q\alpha-p|^{-1}$, energy by energy. The exponents themselves are infinite-time quantities of an infinite system.

\subsection*{How this work was done}

The authors developed the results and proofs with assistance from AI models from more than one developer. They take responsibility for every claim. The authors chose the question, directed the work and decided which lines to pursue; the models did most of the exploration, derivation and drafting.

The proof was developed over several long working sessions, each building on the previous ones and recording its proofs, failed routes and self-checks. Each step was then checked by a separate session that had not seen the material, with instructions to refute it and to read the primary sources it relies on. The repairs found in these checks are incorporated here. Priority and literature searches were AI-assisted. No computation is used in any proof.


\bibliographystyle{plainurl}
\bibliography{references}
\end{document}
